\chapter{变换、矩与表示：通往泛函分析}\label{ch:07}

到目前为止，概率分布、卷积半群和矩序列都已经出现。它们有一个共同的困难：无限对象本身看不见，
我们只能让它作用在一族测试对象上。若测试族选得太小，不同对象可能给出同样答案；若测试族选得太大，
存在性证明又可能依赖尚未建立的抽象定理。于是问题不再只是“怎样计算”，而是怎样选择一族既可计算又足以
识别对象的观测。

本章考察三类格外有效的观测。复参数核把测度变成全纯函数，复指数把平移和卷积变成乘法，正泛函则
记录一个未知测度对测试函数的全部积分。它们各自也有典型的失效情形：点态极限未必保留可微性，Fourier
部分和未必收敛，所有矩都相同也未必决定测度，而一个看似像积分的线性规则也未必来自测度。修补这些
失败所需的工具分别是局部一致控制、正核与完备化、紧性与准解析性，以及正性和正则性。

第一节从复可微的定义出发建立 Cauchy 理论、正规族与参数积分，并把它们用于 Stieltjes 变换和有界
算子的预解式。第二节发展圆周与欧氏空间上的 Fourier 分析，再把热半群和 Poisson 半群统一为
Fourier 乘子。第三节讨论 Hankel/Toeplitz 正性、Helly 选择、矩确定性和 Carleman 条件。第四节证明
Daniell--Stone 与 Riesz--Markov 表示，构造原子和经验逼近，并用三个综合问题贯通全书的存在、识别
与演化机制。

所需的复分析工具都在本章建立。泛函分析方面只使用第三章已证明的 Hahn--Banach
定理；Banach--Alaoglu 定理、一般弱星紧性、自伴算子与谱测度留待后续的泛函分析课程。

\section{复分析工具：全纯极限与参数积分}\label{sec:07-01}

\subsection{全纯性、幂级数与局部一致极限}\label{subsec:7-1-1}
\index{复可微}
\index{全纯函数}
\index{局部一致收敛}
\index{正规族}

实变量函数在一点可微，只说明它在该点附近可由一条直线作一阶逼近；复变量函数的差商却必须对
平面中的所有趋近方向给出同一极限。这个额外约束会逐步产生任意阶导数、幂级数展开和唯一性。
不过，逐点收敛本身不足以保留这种刚性。序列 $p_n(z)=z^n$ 已经提示：每个固定点上的信息既不
包含紧集上的统一误差控制，也不足以保证极限与求导可以交换。

Euler 时代的形式幂级数提供了大量可计算公式；Cauchy 的积分方法解释了这些公式为何适用于所有
全纯函数；Weierstrass 随后把一致收敛放到理论中心。这里先从幂级数这个无需积分公式的模型出发，
再引入紧集上一致收敛。Cauchy 理论将在下一小节说明，这个拓扑为何恰好能保留全纯性和全部导数。

\begin{definition}[复可微与全纯]\label{def:7-1-1-1}
设 $\Omega\subset\C$ 为开集，$f:\Omega\to\C$。若极限
\[
  f'(z_0)=\lim_{h\to0}\frac{f(z_0+h)-f(z_0)}{h}
\]
存在，其中 $h\in\C\setminus\{0\}$ 可沿任意方向趋于 $0$，则称 $f$ 在 $z_0$ \emph{复可微}。
若 $f$ 在 $\Omega$ 的每一点复可微，则称 $f$ 在 $\Omega$ 上\emph{全纯}，并写作
$f\in\mathcal O(\Omega)$。若 $f$ 在每一点附近都可表示为收敛幂级数，则称 $f$ \emph{解析}。
\end{definition}

把 $h$ 限制在实轴和虚轴，便得到 Cauchy--Riemann 方程；反过来，只有实偏导存在还不够，通常还需
连续性等条件才能恢复复可微。以下不沿偏微分方程路线展开，而是直接从差商和积分建立所需结论。

\begin{definition}[局部一致收敛]\label{def:7-1-1-2}
函数列 $f_n:\Omega\to\C$ \emph{局部一致收敛}到 $f$，是指对每个紧集
$K\Subset\Omega$，
\[
  \sup_{z\in K}|f_n(z)-f(z)|\longrightarrow0.
\]
这正是 $\mathcal C(\Omega)$ 上紧开拓扑（compact--open topology）的序列收敛。
\end{definition}

局部一致中的“局部”不能理解为：对每个点临时选一个可随 $n$ 改变的邻域。量词顺序是先固定任意
紧集 $K$，再让同一个上确界误差趋于零。

\begin{definition}[正规族]\label{def:7-1-1-3}
本章采用有限值版本。若全纯函数族
$\mathcal F\subset\mathcal O(\Omega)$ 的每个序列都含有一个在 $\Omega$ 上局部一致收敛、且极限
仍属于 $\mathcal O(\Omega)$ 的子列，则称 $\mathcal F$ 为\emph{正规族}，或称它在
$\mathcal O(\Omega)$ 中相对紧。
\end{definition}

有些复分析教材使用 Riemann 球面度量，并允许子列趋于常值 $\infty$。那是另一种正规族定义；这里
始终只允许有限全纯极限，因而不能把两种版本混用。

\begin{example}[几何级数与预解式原型]\label{ex:7-1-1-1}
若 $|z|<1$，则有限几何和满足
\[
  \sum_{n=0}^{N}z^n=\frac{1-z^{N+1}}{1-z},
  \qquad
  \frac1{1-z}-\sum_{n=0}^{N}z^n=\frac{z^{N+1}}{1-z}.
\]
因此
\[
  \frac1{1-z}=\sum_{n=0}^{\infty}z^n,
  \qquad
  \left(\frac1{1-z}\right)'=\frac1{(1-z)^2}
  =\sum_{n=1}^{\infty}nz^{n-1}.
\]
后一等式也可直接核对，因为
$\sum_{n=1}^{N}nz^{n-1}=(1-(N+1)z^N+Nz^{N+1})/(1-z)^2$，余下两项在
$|z|<1$ 时趋于零。第三小节把标量 $z$ 换成有界算子 $T/z$；同一个几何级数将给出
$(zI-T)^{-1}$。
\end{example}

\begin{example}[局部一致但非全局一致]\label{ex:7-1-1-2}
固定 $0<r<1$。在闭圆盘 $|z|\le r$ 上，几何级数的尾项满足
\[
 \sup_{|z|\le r}\left|\sum_{n=N+1}^{\infty}z^n\right|
 \le\frac{r^{N+1}}{1-r}\longrightarrow0.
\]
所以它在单位圆盘上局部一致收敛。它却不在整个开圆盘上一致收敛：一致收敛的必要条件是
$z^n\to0$ 一致成立，而
\[
  \sup_{|z|<1}|z^n|=1
\]
对每个 $n$ 都成立。复分析中的自然控制因此是“在每个紧子集上”，而不是对整个非紧开域取一次
上确界。
\end{example}

\begin{example}[边界点态观测不足]\label{ex:7-1-1-3}
令 $p_n(z)=z^n$。对任意 $r<1$，
\[
  \sup_{|z|\le r}|p_n(z)|=r^n\longrightarrow0,
  \qquad
  \sup_{|z|\le r}|p_n'(z)|=nr^{n-1}\longrightarrow0.
\]
但在闭单位圆盘的边界上，$p_n(1)=1$，而 $p_n(-1)=(-1)^n$ 甚至不收敛；限制到 $[0,1]$ 时，
点态极限为
\[
  p(x)=\begin{cases}0,&0\le x<1,\\1,&x=1,
  \end{cases}
\]
它在 $1$ 处不连续。这个计算不是开域上“点态极限非全纯”的反例；它准确指出的是，逐点记录本身
不提供紧集上的共同误差，因而不能单独支持连续性或导数交换。
\end{example}

\begin{theorem}[幂级数逐项微分]\label{theorem:7-1-1-1}
设幂级数
\[
  \sum_{n=0}^{\infty}a_n(z-z_0)^n
\]
的收敛半径为 $R\in[0,\infty]$。它在 $D(z_0,R)$ 内定义一个全纯函数 $f$，并且
\[
  f'(z)=\sum_{n=1}^{\infty}na_n(z-z_0)^{n-1}.
  \tag{7.1.1}\label{eq:7-1-1-power-derivative}
\]
导数级数的收敛半径仍为 $R$；因而幂级数可在每个闭小圆盘上任意次逐项微分。
\end{theorem}

\begin{proof}
先设 $0<\rho<R$，并取 $s$ 满足 $\rho<s<R$。级数在 $z_0+s$ 处收敛，所以数列
$|a_n|s^n$ 有界；记其上界为 $M$。于是对 $|z-z_0|\le\rho$，
\[
 |a_n(z-z_0)^n|\le M(\rho/s)^n,
\]
从而原级数由一致收敛的 $M$ 判别法收敛。又有
\[
 n|a_n||z-z_0|^{n-1}
 \le \frac{M}{\rho}\,n(\rho/s)^n
\]
（$z=z_0$ 时单独读取各项），而 $\sum n(\rho/s)^n<\infty$，故形式导数级数也在闭
$\rho$ 圆盘上一致收敛到某个连续函数 $g$。

记部分和为 $S_N$。若 $z,z+h$ 都位于闭 $\rho$ 圆盘内部，则一元实变量微积分给出
\[
 S_N(z+h)-S_N(z)
 =h\int_0^1S_N'(z+th)\dd t.
\]
这里把 $t\mapsto S_N(z+th)$ 看成 $[0,1]$ 上的复值 $C^1$ 函数；逐项使用实变量链式法则有
\[
 \frac{\dd}{\dd t}S_N(z+th)=hS_N'(z+th),
\]
再分别对实部、虚部应用微积分基本定理，便得到上一式。因此这一步没有预先使用尚待证明的
“全纯函数可沿路径积分”结论。
让 $N\to\infty$；$S_N$ 与 $S_N'$ 在包含线段的闭小圆盘上一致收敛，故
\[
 f(z+h)-f(z)=h\int_0^1g(z+th)\dd t.
\]
除以 $h$ 并令 $h\to0$，由 $g$ 的连续性得到 $f'(z)=g(z)$，即
\eqref{eq:7-1-1-power-derivative}。

先设 $R<\infty$。若 $|w-z_0|>R$ 而 $a_n(w-z_0)^n\to0$，则这个项列有界。任取
$R<s<|w-z_0|$，便有
\[
 |a_n|s^n
 =|a_n(w-z_0)^n|\left(\frac{s}{|w-z_0|}\right)^n.
\]
右端由一个收敛几何级数逐项支配，因而原幂级数在 $|z-z_0|=s$ 上绝对收敛，与其收敛半径为
$R$ 矛盾。故每个 $|w-z_0|>R$ 都使 $a_n(w-z_0)^n$ 不趋于零。若导数级数在 $w$ 处收敛，则
$na_n(w-z_0)^{n-1}\to0$，进而
$a_n(w-z_0)^n=(w-z_0)na_n(w-z_0)^{n-1}/n\to0$，矛盾。因此导数级数在每个
$|w-z_0|>R$ 处发散；结合前一段，它的半径正是 $R$。对导数级数重复同一论证即可任意次逐项
微分。$R=0$ 时关于开圆盘的断言为空；$R=\infty$ 时上述论证对任意有限 $\rho$ 成立。
\end{proof}

幂级数情形提示了下面三条一般结论。先予陈述；它们将在下一小节分别由 Morera 定理、
Cauchy 估计和恒等定理证明。

\begin{theorem}[Weierstrass 极限定理]\label{theorem:7-1-1-2}
若 $f_n\in\mathcal O(\Omega)$ 且 $f_n$ 在 $\Omega$ 上局部一致收敛到 $f$，则
$f\in\mathcal O(\Omega)$。此外，对每个 $k\ge1$，
\[
  f_n^{(k)}\longrightarrow f^{(k)}
\]
也在 $\Omega$ 上局部一致成立。
\end{theorem}

\begin{theorem}[Montel 定理]\label{theorem:7-1-1-3}
设 $\mathcal F\subset\mathcal O(\Omega)$。若对每个紧集 $K\Subset\Omega$ 都有
\[
  \sup_{f\in\mathcal F}\sup_{z\in K}|f(z)|<\infty,
  \tag{7.1.2}\label{eq:7-1-1-locally-bounded-family}
\]
则 $\mathcal F$ 是正规族。
\end{theorem}

这里要求的是整个函数族在每个紧集上有同一个界；“每个 $f$ 各自在局部有界”只是连续函数的普通
性质，完全不足以推出函数族紧性。

\begin{proposition}[Vitali 收敛定理（选读）]\label{proposition:7-1-1-4}
设 $\Omega$ 连通，$(f_n)\subset\mathcal O(\Omega)$ 满足局部一致有界条件
\eqref{eq:7-1-1-locally-bounded-family}。若存在集合 $E\subset\Omega$，它在 $\Omega$ 内有聚点，
并且对每个 $z\in E$ 数列 $f_n(z)$ 收敛，则 $(f_n)$ 在 $\Omega$ 上局部一致收敛到某个全纯函数。
\end{proposition}

三条结论之间的几何关系可由图~\ref{fig:7-1-1-exhaustion-and-power-series} 读取。紧耗尽把一个非紧开域
拆成可使用 \ArzelaAscoli{} 定理的紧层；在第 $m$ 层抽取的子列又作为第 $m+1$ 层的起点。对角项
最终在每个固定层上都落入相应子列的尾部。右侧的收敛圆则说明，幂级数的控制也总是先固定
$r<R$，再在闭小圆上作一致估计。

\begin{figure}[htbp]
\centering
\begin{tikzpicture}[scale=0.92]
  \begin{scope}[xshift=0cm]
    \draw[book region] (0,0) ellipse (2.05 and 1.55);
    \draw[BookAccent!75,semithick] (0,0) ellipse (1.48 and 1.05);
    \draw[BookWarm,semithick] (0,0) ellipse (0.88 and 0.58);
    \node[book label] at (0,1.28) {$\Omega$};
    \node[book label] at (1.10,0.74) {$K_3$};
    \node[book label] at (0.73,0.34) {$K_2$};
    \node[book label] at (0,0) {$K_1$};
    \node[book label] at (0,-1.92) {$K_1\Subset K_2\Subset K_3\Subset\cdots$};
  \end{scope}
  \begin{scope}[xshift=4.25cm,yshift=0.82cm,node distance=7mm]
    \node[book node] (s0) {$(f_n)$};
    \node[book node,below=of s0] (s1) {$(f_n^{[1]})$：在 $K_1$ 收敛};
    \node[book node,below=of s1] (s2) {$(f_n^{[2]})$：在 $K_2$ 收敛};
    \node[book node,below=of s2] (s3) {$(f_n^{[3]})$：在 $K_3$ 收敛};
    \draw[book arrow] (s0)--(s1);
    \draw[book arrow] (s1)--(s2);
    \draw[book arrow] (s2)--(s3);
    \node[book warm node,right=9mm of s2] (diag) {对角列\\$f_n^{[n]}$};
    \draw[book accent arrow] (s1.east) to[bend left=18] (diag.north);
    \draw[book accent arrow] (s2.east)--(diag.west);
    \draw[book accent arrow] (s3.east) to[bend right=18] (diag.south);
  \end{scope}
  \begin{scope}[xshift=10.1cm]
    \draw[BookAccent,semithick,fill=BookAccent!5] (0,0) circle (1.48);
    \draw[BookWarm,semithick,dashed] (0,0) circle (0.88);
    \fill[BookInk] (0,0) circle (1.4pt);
    \draw[book arrow] (0,0)--node[above,book label]{$r$}(0.88,0);
    \draw[book arrow] (0,0)--node[below,book label]{$R$}(1.48,0);
    \node[book label] at (0,-1.83) {先固定 $r<R$};
    \node[book label] at (-0.28,0.22) {$z_0$};
  \end{scope}
\end{tikzpicture}
\caption{紧耗尽上的逐层抽取与幂级数闭小圆控制。两者都把非紧或临界边界问题化为固定紧集上的
一致估计。}
\label{fig:7-1-1-exhaustion-and-power-series}
\end{figure}

这些结论将在不同方向反复出现：预解式级数和 Stieltjes 变换需要局部一致极限保持全纯；有理
逼近需要从局部有界函数族中提取子列；而 Montel 的证明把第一章的 \ArzelaAscoli{} 定理与复分析
的导数估计接到了一起。三者共同表达同一个原则：有限点值观测必须先升级为紧集一致控制，复解析刚性
才能开始工作。

\begin{exercise}
\begin{enumerate}
  \item 对 $|z-z_0|\le r<s<R$，从 $|a_n|s^n$ 有界出发，分别写出原级数及前两阶导数级数的
        $M$ 判别控制数列。
  \item 在单位圆盘上令 $f_n(z)=nz$。验证 $f_n(0)$ 收敛，但函数族在紧圆盘
        $|z|\le1/2$ 上不一致有界；说明只在一个没有域内聚点的集合上观察收敛，为何不足以应用
        Vitali 定理。
  \item 设 $f_n(z)=z^n$。计算它在 $|z|\le r<1$ 上第 $k$ 阶导数的上确界，并证明该上确界趋于
        $0$。
  \item 使用 Cauchy 估计~\eqref{eq:7-1-2-cauchy-estimate}，推出满足
        \eqref{eq:7-1-1-locally-bounded-family} 的函数族在每个紧子集上等度连续。
\end{enumerate}
\end{exercise}

曲线积分将把边界上的一致信息传回内部。下一步先从 Goursat 的三角形细分出发建立 Cauchy 公式，
再回来完成上述三条极限与紧性结论的证明。


\subsection{曲线积分、Cauchy 公式、Morera 与唯一性}\label{subsec:7-1-2}
\index{复曲线积分}
\index{Cauchy--Goursat 定理}
\index{Cauchy 积分公式}
\index{Morera 定理}
\index{恒等定理}
\index{最大模原理}

复导数由一点附近的差商定义，Cauchy 公式却能从一条圆周上的函数值恢复圆内一点的全部导数。这种
由边界决定内部的现象，是全纯函数远比一般光滑函数刚性的根源。反方向上，Morera 定理又把闭路积分
为零变成全纯性的判据；以后遇到含参数积分时，这往往比直接控制复差商更方便。

\begin{definition}[复曲线积分]\label{def:7-1-2-1}
设 $\gamma:[a,b]\to\Omega$ 为分段 $C^1$ 曲线，$f:\Omega\to\C$ 在 $\gamma$ 的像上连续。
定义
\[
  \int_\gamma f(z)\dd z
  =\int_a^b f(\gamma(t))\gamma'(t)\dd t,
  \tag{7.1.3}\label{eq:7-1-2-curve-integral}
\]
其中右侧按实参数逐段积分。若 $\gamma^{-}(t)=\gamma(a+b-t)$ 表示反向参数化，则
$\int_{\gamma^-}f\dd z=-\int_\gamma f\dd z$；首尾相接曲线的拼接积分等于各段积分之和。
\end{definition}

若 $\gamma$ 为闭曲线，仅凭 $f$ 在曲线附近有复导数，并不能在任意开域上断言积分为零；域中的孔洞
可能被积分检测到。原函数恰好把这个全局障碍编码出来。

\begin{definition}[原函数]\label{def:7-1-2-2}
若 $F\in\mathcal O(\Omega)$ 且 $F'=f$，则称 $F$ 是 $f$ 在 $\Omega$ 上的\emph{原函数}。
\end{definition}

\begin{proposition}[原函数与路径无关]\label{proposition:7-1-2-primitive-criterion}
设 $f:\Omega\to\C$ 连续。则 $f$ 在 $\Omega$ 的每个连通分支上有原函数，当且仅当它沿该分支内
每条分段 $C^1$ 闭曲线的积分都为零。特别地，若 $\Omega$ 关于 $z_*$ 星形，且每个内含三角形的
边界积分为零，则
\[
  F(z)=\int_{[z_*,z]}f(\zeta)\dd\zeta
  \tag{7.1.4}\label{eq:7-1-2-radial-primitive}
\]
是 $f$ 的原函数。
\end{proposition}

\begin{proof}
若 $F'=f$，则在每个 $C^1$ 曲线段上，实变量链式法则给
\[
 \frac{\dd}{\dd t}F(\gamma(t))=f(\gamma(t))\gamma'(t).
\]
逐段积分可得 $\int_\gamma f\dd z=F(\gamma(b))-F(\gamma(a))$，所以闭曲线积分为零。

反过来，复平面中开连通集是折线路径连通的。在一个连通分支中固定 $z_*$；对任意 $z$ 选一条从
$z_*$ 到 $z$ 的分段 $C^1$ 路径 $\gamma_z$，令 $F(z)=\int_{\gamma_z}f\dd z$。两条这样的路径
首尾拼接成闭曲线，故题设保证 $F$ 与路径选择无关。取 $r>0$ 使 $D(z,r)\subset\Omega$。当
$|h|<r$ 时，可把 $\gamma_z$ 与线段 $[z,z+h]$ 拼接为 $\gamma_{z+h}$，于是
\[
 \frac{F(z+h)-F(z)}h
 =\int_0^1 f(z+th)\dd t\longrightarrow f(z),
\]
最后一步来自 $f$ 在 $z$ 连续。因此 $F'=f$。

若 $\Omega$ 关于 $z_*$ 星形，式~\eqref{eq:7-1-2-radial-primitive} 中的线段都在 $\Omega$ 内。
对充分小的 $h$，三条线段 $[z_*,z]$、$[z,z+h]$、$[z+h,z_*]$ 围成的三角形位于 $\Omega$：
边 $[z,z+h]$ 位于一个小圆盘内，而从 $z_*$ 到这条边上每一点的线段由星形性包含于 $\Omega$。
三角形边界积分为零，因而仍有
$F(z+h)-F(z)=\int_{[z,z+h]}f\dd z$；上面的差商计算完成证明。
\end{proof}

\begin{theorem}[Cauchy--Goursat]\label{theorem:7-1-2-1}
若 $f$ 在包含闭三角形 $T$ 的开集上全纯，则
\[
  \int_{\partial T}f(z)\dd z=0,
  \tag{7.1.5}\label{eq:7-1-2-goursat-triangle}
\]
其中 $\partial T$ 取正向。因而，全纯函数在星形域中有原函数，并沿域内每条分段 $C^1$ 闭曲线
积分为零。
\end{theorem}

\begin{proof}
记 $I(T)=\int_{\partial T}f(z)\dd z$。连接三边中点，把 $T=T_0$ 分成四个相似小三角形；给各
边界正向后，内部边的积分两两抵消，所以四个小三角形的边界积分之和等于 $I(T_0)$。其中至少一个
小三角形 $T_1$ 满足
\[
  |I(T_1)|\ge\frac14|I(T_0)|.
\]
重复选择，得到嵌套闭三角形 $T_n$，并且
\[
 |I(T_n)|\ge4^{-n}|I(T_0)|,
 \qquad
 \operatorname{per}(T_n)=2^{-n}\operatorname{per}(T_0).
 \tag{7.1.6}\label{eq:7-1-2-goursat-scale}
\]
它们的直径趋于零，故嵌套紧集定理给出唯一点
$z_*\in\bigcap_nT_n$。

由 $f$ 在 $z_*$ 复可微，可写成
\[
 f(z)=f(z_*)+f'(z_*)(z-z_*)+r(z),
 \qquad
 \frac{|r(z)|}{|z-z_*|}\longrightarrow0\quad(z\to z_*).
 \tag{7.1.7}\label{eq:7-1-2-goursat-remainder}
\]
常数项沿闭折线的积分为零；线性项也满足
\[
 \int_{\partial T_n}(z-z_*)\dd z
 =\frac12\int_{\partial T_n}\dd (z-z_*)^2=0.
\]
把函数
\[
 \eta(z)=\frac{|r(z)|}{|z-z_*|}\quad(z\ne z_*),
 \qquad \eta(z_*)=0
\]
连续延拓到 $z_*$，并令 $\varepsilon_n=\max_{z\in T_n}\eta(z)$。由
式~\eqref{eq:7-1-2-goursat-remainder}，对任意 $\varepsilon>0$ 可取 $\delta_\varepsilon>0$，使
$|z-z_*|<\delta_\varepsilon$ 时 $\eta(z)<\varepsilon$。再由 $z_*\in T_n$ 和
$\diam(T_n)\to0$，充分大的 $n$ 满足 $T_n\subset D(z_*,\delta_\varepsilon)$，从而
$\varepsilon_n\leq\varepsilon$；故 $\varepsilon_n\to0$。于是
\[
 |I(T_n)|
 \le \sup_{\partial T_n}|r(z)|\operatorname{per}(T_n)
 \le \varepsilon_n\diam(T_n)\operatorname{per}(T_n)
 \le \varepsilon_n\operatorname{per}(T_n)^2.
 \tag{7.1.8}\label{eq:7-1-2-goursat-upper}
\]
另一方面，式~\eqref{eq:7-1-2-goursat-scale} 给出
\[
 \frac{|I(T_n)|}{\operatorname{per}(T_n)^2}
 \ge\frac{|I(T_0)|}{\operatorname{per}(T_0)^2}.
\]
与式~\eqref{eq:7-1-2-goursat-upper} 合并并令 $n\to\infty$，得到 $I(T_0)=0$。
这里必须比较周长的平方：余项是
$o(\diam T_n)\operatorname{per}(T_n)=o(\operatorname{per}(T_n)^2)$；若误写成除以一次周长，
缩放量纲便不相容。

现在设 $\Omega$ 关于 $z_*$ 星形。刚证的三角形积分为零，命题
\ref{proposition:7-1-2-primitive-criterion} 给出原函数。对任意分段 $C^1$ 闭曲线应用原函数的
链式法则，积分即为零。
\end{proof}

\begin{theorem}[Cauchy 积分公式]\label{theorem:7-1-2-2}
设 $f\in\mathcal O(\Omega)$ 且 $\overline{D(a,r)}\subset\Omega$。若 $|w-a|<r$，则
\[
 f(w)=\frac1{2\pi i}\int_{|z-a|=r}\frac{f(z)}{z-w}\dd z.
 \tag{7.1.9}\label{eq:7-1-2-cauchy-value}
\]
函数 $f$ 在圆内无限次复可微，并且对每个 $n\ge0$，
\[
 f^{(n)}(w)=\frac{n!}{2\pi i}
 \int_{|z-a|=r}\frac{f(z)}{(z-w)^{n+1}}\dd z.
 \tag{7.1.10}\label{eq:7-1-2-cauchy-derivative}
\]
特别地，在 $w=a$ 便得到以圆心为中心的公式。每个全纯函数都是解析函数。
\end{theorem}

\begin{proof}
先记录 Goursat 证明的一个可去点版本：若 $h$ 在闭三角形 $T$ 上连续，并在 $T$ 附近除至多一点
$w$ 外全纯，则 $\int_{\partial T}h\dd z=0$。令 $M=\max_T|h|$，并对 $T$ 作 $m$ 次中点四分。
每个第 $m$ 层小三角形的周长都是 $2^{-m}\operatorname{per}(T)$；由三角网格的局部组合结构，包含
$w$ 的小三角形至多有六个。不接触 $w$ 的小三角形在一个避开 $w$ 的开邻域内全纯，故由
定理~\ref{theorem:7-1-2-1} 得零积分。把全部小三角形的正向边界积分相加，内部边两两抵消，于是
\[
 \left|\int_{\partial T}h(z)\dd z\right|
 \leq 6M\,2^{-m}\operatorname{per}(T)\longrightarrow0.
\]
这就证明可去点版本，而没有预先假设 $h$ 在 $w$ 复可微。

由 $\overline{D(a,r)}$ 紧含于 $\Omega$，可取 $r_1>r$ 使
$\overline{D(a,r_1)}\subset\Omega$。固定 $w\in D(a,r)$，定义
\[
 h_w(z)=
 \begin{cases}
 \dfrac{f(z)-f(w)}{z-w},&z\ne w,\\[6pt]
 f'(w),&z=w.
 \end{cases}
\]
复可微的定义说明 $h_w$ 在 $w$ 连续；它在 $w$ 外全纯。上一段说明它在 $D(a,r_1)$ 内每个三角形上的
边界积分为零。这个圆盘是星形域，所以命题~\ref{proposition:7-1-2-primitive-criterion} 的构造给出
$h_w$ 的原函数，进而
\[
 \int_{|z-a|=r}h_w(z)\dd z=0.
 \tag{7.1.11}\label{eq:7-1-2-divided-difference-zero}
\]

直接计算 Cauchy 核的积分。令 $q=(w-a)/r$，则 $|q|<1$，沿
$z=a+re^{it}$ 有
\begin{align*}
 \int_{|z-a|=r}\frac{\dd z}{z-w}
 &=i\int_0^{2\pi}\frac{1}{1-qe^{-it}}\dd t\\
 &=i\sum_{k=0}^{\infty}q^k\int_0^{2\pi}e^{-ikt}\dd t
 =2\pi i.
\end{align*}
这里几何级数因 $|q|<1$ 在 $t\in[0,2\pi]$ 上一致收敛，故可逐项积分。将
\eqref{eq:7-1-2-divided-difference-zero} 展开，就得到
\eqref{eq:7-1-2-cauchy-value}。

现在把积分圆周固定，只让 $w$ 在其内部变化。对任意 $\rho<r$，当 $|w-a|\le\rho$ 时，
$|z-w|\ge r-\rho$。若 $m\ge1$ 且 $|h|<(r-\rho)/2$，则在圆周上一致地有
\[
 \frac{(z-w-h)^{-m}-(z-w)^{-m}}h
 =m\int_0^1(z-w-th)^{-m-1}\dd t,
\]
其绝对值不超过 $m((r-\rho)/2)^{-m-1}$，并在 $h\to0$ 时一致趋于
$m(z-w)^{-m-1}$。因此曲线参数积分的差商可在一致极限下通过积分号，而且
\[
 \frac{\partial}{\partial w}(z-w)^{-m}
 =m(z-w)^{-m-1}.
\]
从 \eqref{eq:7-1-2-cauchy-value} 开始归纳，得到
\eqref{eq:7-1-2-cauchy-derivative}。

最后，若 $|w-a|\le\rho<r$，则圆周上一致地有
\[
 \frac1{z-w}
 =\frac1{z-a}\sum_{n=0}^{\infty}\left(\frac{w-a}{z-a}\right)^n.
\]
代入式~\eqref{eq:7-1-2-cauchy-value} 并逐项积分，得到
\[
 f(w)=\sum_{n=0}^{\infty}c_n(w-a)^n,
 \qquad
 c_n=\frac1{2\pi i}\int_{|z-a|=r}
       \frac{f(z)}{(z-a)^{n+1}}\dd z
     =\frac{f^{(n)}(a)}{n!}.
\]
所以 $f$ 在每一点附近等于其 Taylor 级数，即 $f$ 解析。
\end{proof}

\begin{theorem}[Morera 定理]\label{theorem:7-1-2-3}
设 $f:\Omega\to\C$ 连续。若对每个满足 $\overline T\subset\Omega$ 的三角形 $T$ 都有
\[
  \int_{\partial T}f(z)\dd z=0,
\]
则 $f\in\mathcal O(\Omega)$。
\end{theorem}

\begin{proof}
固定 $a\in\Omega$，取 $r>0$ 使 $D(a,r)\subset\Omega$。圆盘是星形域；由题设及
命题~\ref{proposition:7-1-2-primitive-criterion}，
\[
 F(z)=\int_{[a,z]}f(\zeta)\dd\zeta
\]
在 $D(a,r)$ 上满足 $F'=f$。于是 $F$ 全纯。定理~\ref{theorem:7-1-2-2} 已说明全纯函数无限次
复可微，故 $f=F'$ 也全纯。点 $a$ 任意，所以 $f\in\mathcal O(\Omega)$。
\end{proof}

\begin{corollary}[刚性推论]\label{corollary:7-1-2-4}
设 $f\in\mathcal O(\Omega)$。
\begin{enumerate}
  \item 若 $\overline{D(a,r)}\subset\Omega$，并记
        $M_r=\max_{|z-a|=r}|f(z)|$，则对每个 $n\ge0$，
        \[
          |f^{(n)}(a)|\le\frac{n!}{r^n}M_r.
          \tag{7.1.12}\label{eq:7-1-2-cauchy-estimate}
        \]
  \item 每个有界整函数都是常数（Liouville 定理）；每个非常数复系数多项式都在 $\C$ 中有零点
        （基本代数定理）。
  \item 若 $\Omega$ 连通且 $|f|$ 在域内一点取得局部最大值，则 $f$ 为常数（最大模原理）。
  \item 若 $\Omega$ 连通，$f,g\in\mathcal O(\Omega)$，且集合
        $\{z\in\Omega:f(z)=g(z)\}$ 在 $\Omega$ 内有聚点，则 $f\equiv g$（恒等定理）。
\end{enumerate}
\end{corollary}

\begin{proof}
由 Cauchy 导数公式和曲线积分估计，
\[
 |f^{(n)}(a)|
 \le\frac{n!}{2\pi}\,(2\pi r)\frac{M_r}{r^{n+1}}
 =\frac{n!}{r^n}M_r,
\]
这就是式~\eqref{eq:7-1-2-cauchy-estimate}。

若 $f$ 是满足 $|f|\le M$ 的整函数，则对任意 $a\in\C$、$R>0$，Cauchy 估计给
$|f'(a)|\le M/R$。令 $R\to\infty$ 得 $f'(a)=0$；沿任意线段积分可知 $f$ 为常数。

设非常数多项式 $p(z)=a_mz^m+\cdots+a_0$ 没有零点。则 $1/p$ 是整函数。取 $R$ 充分大，使当
$|z|\ge R$ 时
\[
 |a_{m-1}z^{m-1}+\cdots+a_0|
 \le\frac{|a_m|}{2}|z|^m;
\]
于是 $|p(z)|\ge |a_m||z|^m/2$，所以 $1/p$ 在圆盘外有界。它在紧圆盘 $|z|\le R$ 上也因连续
且分母不为零而有界。Liouville 定理迫使 $1/p$、从而 $p$ 为常数，矛盾。这证明基本代数定理。

先证明恒等定理。令 $h=f-g$，设它的零点列 $z_j\ne a$ 趋于域内点 $a$。若 $h$ 在 $a$ 的 Taylor
级数有首个非零系数 $c_m$，则
\[
 h(z)=(z-a)^m\bigl(c_m+(z-a)c_{m+1}+\cdots\bigr).
\]
括号中的全纯函数在 $a$ 取非零值，因连续性在 $a$ 的某邻域内不为零；这与 $z_j\to a$ 且
$h(z_j)=0$ 矛盾。因此 $h$ 的全部 Taylor 系数为零，$h$ 在 $a$ 的某邻域恒为零。

令
\[
 A=\{z\in\Omega:h\text{ 在 }z\text{ 的某邻域恒为零}\}.
\]
$A$ 非空且开。若 $z_j\in A$ 且 $z_j\to z\in\Omega$，则对每个 $k\ge0$，
$h^{(k)}(z_j)=0$。Cauchy 公式保证各阶导数连续，所以 $h^{(k)}(z)=0$；$h$ 在 $z$ 的 Taylor 级数
恒为零，故 $z\in A$。于是 $A$ 在 $\Omega$ 中也闭。连通性给 $A=\Omega$，即 $f\equiv g$。

最后证明最大模原理。设 $|f(z)|\le|f(a)|$ 在某个圆盘 $D(a,r)$ 内成立。若 $f(a)=0$，则 $f$ 在
该圆盘恒为零，恒等定理给出全域结论。以下设 $f(a)\ne0$，取单位复数 $\eta$ 使
$\eta f(a)=|f(a)|$。对每个 $0<s<r$，Cauchy 公式给出平均值等式
\[
 |f(a)|=\eta f(a)
 =\frac1{2\pi}\int_0^{2\pi}\eta f(a+se^{it})\dd t.
\]
连续函数
\[
 u_s(t)=|f(a)|-\operatorname{Re}\bigl(\eta f(a+se^{it})\bigr)
\]
非负且积分为零，故处处为零。再结合
$|f(a+se^{it})|\le|f(a)|$，得到
$\eta f(a+se^{it})=|f(a)|$。任意 $z\in D(a,r)\setminus\{a\}$ 都在某条这样的圆周上，故
$f=f(a)$ 于整个圆盘；恒等定理把等式推广到连通域 $\Omega$。
\end{proof}

\begin{definition}[零点阶数]\label{def:7-1-2-3}
设 $z_0\in\Omega$，$f\in\mathcal O(\Omega)$，并且 $f$ 在 $z_0$ 所在的连通分支上不恒为零。若
$f(z_0)=0$，则恒等定理保证下列 Taylor 展开不可能所有系数都为零。写
\[
 f(z)=\sum_{n=0}^{\infty}a_n(z-z_0)^n
\]
在 $z_0$ 附近成立，并令 $m\ge1$ 为满足 $a_m\ne0$ 的最小指标；称 $m$ 为 $f$ 在 $z_0$ 的
\emph{零点阶数}。若 $f(z_0)\ne0$，则约定其消失阶数为 $0$，但不称 $z_0$ 为零点。
\end{definition}

在零点情形，把 Taylor 级数的前 $m$ 项提出，便有
\[
 f(z)=(z-z_0)^m g(z),
 \qquad g\in\mathcal O(\Omega_0),\quad g(z_0)=a_m\ne0
\]
在 $z_0$ 的某邻域 $\Omega_0$ 内成立。由 $g$ 的连续性，缩小 $\Omega_0$ 后可使 $g$ 处处非零，
所以 $z_0$ 是孤立零点。于是，一个在连通域上不恒为零的全纯函数，其零点只能是孤立的；若零点在
域内有聚点，恒等定理便迫使函数恒为零。

Morera 定理识别局部一致极限，Cauchy 估计提供函数族的紧性，恒等定理则识别
所得子列极限。下面依次完成三项证明。

\begin{proof}[定理~\ref{theorem:7-1-1-2} 的证明]
局部一致极限 $f$ 连续。事实上，每一点都有闭小圆盘紧含于 $\Omega$，而连续函数在该紧集上的一致
极限连续。若 $\overline T\subset\Omega$ 为三角形，则
\[
 \left|\int_{\partial T}(f_n-f)(z)\dd z\right|
 \le \operatorname{per}(T)\sup_{z\in\partial T}|f_n(z)-f(z)|\longrightarrow0.
\]
由 Cauchy--Goursat 定理，$\int_{\partial T}f_n\dd z=0$；故
$\int_{\partial T}f\dd z=0$。Morera 定理给出 $f\in\mathcal O(\Omega)$。

固定紧集 $K\Subset\Omega$ 和 $k\ge1$。可取 $\delta>0$，使闭 $2\delta$ 邻域
\[
 L=\{z\in\C:\dist(z,K)\le2\delta\}
\]
紧含于 $\Omega$；当 $\Omega=\C$ 时任取充分小的 $\delta$ 即可。对每个 $a\in K$，在圆周
$|z-a|=\delta$ 上应用 Cauchy 导数公式，得到
\[
 |f_n^{(k)}(a)-f^{(k)}(a)|
 \le\frac{k!}{\delta^k}\sup_{z\in L}|f_n(z)-f(z)|.
\]
右侧与 $a$ 无关并趋于零，所以 $f_n^{(k)}\to f^{(k)}$ 在 $K$ 上一致。$K$ 与 $k$ 任意，定理得证。
\end{proof}

\begin{proof}[定理~\ref{theorem:7-1-1-3} 的证明]
由命题~\ref{proposition:1-4-2-2}，平面开集 $\Omega$ 有紧耗尽
\[
 K_1\subset\operatorname{int}K_2\subset K_2
 \subset\operatorname{int}K_3\subset\cdots,
 \qquad \bigcup_{m\ge1}K_m=\Omega.
 \tag{7.1.13}\label{eq:7-1-2-compact-exhaustion}
\]
而且每个紧子集都包含在某个 $K_m$ 中。

固定 $m$。取 $\delta>0$，使 $K_m$ 的闭 $2\delta$ 邻域 $L_m$ 紧含于 $\Omega$。由局部一致
有界性，存在 $M_m$ 使 $|f|\le M_m$ 对所有 $f\in\mathcal F$、$z\in L_m$ 成立。当
$\dist(\zeta,K_m)\le\delta$ 时，圆 $|z-\zeta|=\delta$ 位于 $L_m$；Cauchy 估计给
\[
 |f'(\zeta)|\le\frac{M_m}{\delta}.
\]
若 $z,w\in K_m$ 且 $|z-w|<\delta$，连接它们的线段位于 $K_m$ 的闭 $\delta$ 邻域，因此
\[
 |f(z)-f(w)|
 \le\frac{M_m}{\delta}|z-w|
 \qquad(f\in\mathcal F).
 \tag{7.1.14}\label{eq:7-1-2-montel-equicontinuity}
\]
故 $\mathcal F|_{K_m}$ 等度连续且逐点有界。\ArzelaAscoli{} 定理
\ref{theorem:1-3-2-2} 说明它在 $C(K_m)$ 中相对紧。

任取序列 $(f_n)\subset\mathcal F$。先抽取在 $K_1$ 上一致收敛的子列，再从中抽取在 $K_2$ 上
一致收敛的子列，如此继续。取第 $m$ 层子列的第 $m$ 项组成对角列 $(g_m)$。对每个固定 $j$，
$(g_m)_{m\ge j}$ 是第 $j$ 层子列的子列，故在 $K_j$ 上一致收敛；各层极限在交叠处相同，因而
定义出 $g:\Omega\to\C$。式~\eqref{eq:7-1-2-compact-exhaustion} 保证收敛在每个紧子集上一致。
刚证的 Weierstrass 极限定理给 $g\in\mathcal O(\Omega)$。所以 $\mathcal F$ 是正规族。
\end{proof}

\begin{proof}[命题~\ref{proposition:7-1-1-4} 的证明]
由 Montel 定理，$(f_n)$ 有局部一致收敛子列；记其极限为 $g\in\mathcal O(\Omega)$。若另一子列
局部一致收敛到 $h$，则对每个 $z\in E$，原数列 $f_n(z)$ 的收敛性给出
\[
 g(z)=\lim_nf_n(z)=h(z).
\]
$E$ 在域内有聚点，恒等定理遂给 $g\equiv h$。

还需排除原序列不收敛到 $g$ 的可能。若局部一致收敛失败，则存在紧集 $K\Subset\Omega$、
$\varepsilon>0$ 和子列 $(f_{n_j})$，使
\[
 \sup_K|f_{n_j}-g|\ge\varepsilon\qquad(j\ge1).
\]
Montel 定理又从这条子列抽出局部一致收敛的子列，其极限必须等于 $g$，与上式矛盾。因此整个序列
局部一致收敛到 $g$。
\end{proof}

\begin{theorem}[Osgood 定理（选读）]\label{theorem:7-1-2-osgood}
若 $f_n\in\mathcal O(\Omega)$ 只逐点收敛到函数 $f$，则不能断言 $f$ 在整个 $\Omega$ 上全纯；不过
$f$ 必在某个稠密开子集上全纯。
\end{theorem}

因此，逐点收敛仍保留一个稠密的局部解析区域，却不提供该区域之外的共同局部控制。

\begin{proof}
任取非空开集 $U\subset\Omega$，并在其中取一个正半径的闭圆盘 $K\Subset U$。逐点收敛说明
$(f_n(z))_n$ 对每个 $z\in K$ 有界。令
\[
 E_m=\{z\in K:|f_n(z)|\le m\text{ 对所有 }n\},\qquad m\ge1.
\]
每个 $E_m$ 都闭，且 $K=\bigcup_mE_m$。紧度量空间 $K$ 完备；Baire 定理
\ref{theorem:1-4-1-1} 给出某个 $E_m$ 在 $K$
中有非空相对内点；该相对开集与 $\operatorname{int}K$ 相交，因而可取开圆盘
$D\Subset\operatorname{int}K$，使 $|f_n|\le m$ 在 $D$ 上对所有 $n$ 成立。
于是 $(f_n)$ 在 $D$ 上局部一致有界，而它在那里处处收敛。命题~\ref{proposition:7-1-1-4} 给出
$f_n\to f$ 在 $D$ 上局部一致，故 $f|_D$ 全纯。

令 $V$ 为 $f$ 在某邻域全纯的所有点之集。按定义 $V$ 开；刚才对任意非空 $U$ 都找到了
$D\subset U\cap V$，所以 $V$ 稠密。
\end{proof}

\begin{example}[绕原点积分]\label{ex:7-1-2-1}
令 $\gamma(t)=re^{it}$，$0\le t\le2\pi$，其中 $r>0$。由定义直接得到
\[
 \int_{|z|=r}\frac{\dd z}{z}
 =\int_0^{2\pi}\frac{ire^{it}}{re^{it}}\dd t
 =2\pi i.
\]
因此 $1/z$ 不可能在 $\C\setminus\{0\}$ 上有全局原函数，否则闭曲线积分必须为零。另一方面，
$1/z$ 在任何不绕过原点的小圆盘内都有局部原函数。这个计算区分了局部全纯与有孔域上的全局
路径无关。
\end{example}

\begin{example}[Cauchy 核]\label{ex:7-1-2-2}
若 $\overline{D(a,r)}\subset\Omega$，则
\[
 f(a)=\frac1{2\pi i}\int_{|z-a|=r}\frac{f(z)}{z-a}\dd z,
 \qquad
 f^{(n)}(a)=\frac{n!}{2\pi i}\int_{|z-a|=r}
                  \frac{f(z)}{(z-a)^{n+1}}\dd z.
\]
例如取 $f(z)=e^z$、$a=0$，便得到
\[
 \int_{|z|=r}\frac{e^z}{z^{n+1}}\dd z
 =\frac{2\pi i}{n!}f^{(n)}(0)=\frac{2\pi i}{n!}.
\]
圆周数据由同一个核恢复中心值及全部导数。若把圆周换成一般闭曲线，则公式还需绕数和适当的域条件；
仅说“点在曲线里面”并不是严格假设。
\end{example}

\begin{example}[Morera 参数积分]\label{ex:7-1-2-3}
设 $\mu$ 为有限复测度，$x\mapsto q(z,x)$ 对每个 $z\in\Omega$ 可测，而对 $|\mu|$-几乎每个
$x$，函数 $z\mapsto q(z,x)$ 全纯。再假设对每个紧集 $K\Subset\Omega$，存在
$g_K\in L^1(|\mu|)$ 使
\[
 |q(z,x)|\le g_K(x)\qquad(z\in K)
\]
在一个共同零集外成立。令 $Q(z)=\int q(z,x)\mu(\dd x)$。若 $z_n\to z$，可把 $z$ 与充分靠后的
$z_n$ 放进同一紧小圆盘；支配收敛给 $Q(z_n)\to Q(z)$，故 $Q$ 连续。

对任意 $\overline T\subset\Omega$，取包含 $\partial T$ 的紧集 $K$。曲线长度测度与 $|\mu|$ 的
乘积下，
\[
 \int_{\partial T}\int |q(z,x)|\,|\mu|(\dd x)|\dd z|
 \le \operatorname{per}(T)\int g_K\dd|\mu|<\infty.
\]
在全纯性成立的共同满测集之外把 $q$ 改定义为零。对 $\partial T$ 的每个 $C^1$ 参数段
$\gamma$，函数 $(t,x)\mapsto q(\gamma(t),x)$ 对 $t$ 连续、对 $x$ 可测；用有限网格的左端点
阶梯函数逼近 $t$，可知它关于乘积 $\sigma$-代数可测。因此 Fubini
定理~\ref{theorem:3-2-1-2} 确实适用，并给出
\[
 \int_{\partial T}Q(z)\dd z
 =\int\left(\int_{\partial T}q(z,x)\dd z\right)\mu(\dd x)=0.
\]
Morera 定理推出 $Q$ 全纯。第三小节将把这一计算整理成参数积分定理。
\end{example}

Goursat 证明与 Cauchy 公式使用两种互补的几何压缩。图
\ref{fig:7-1-2-goursat-and-cauchy} 左侧不断缩小三角形，比较的是边界积分与周长平方；右侧则把固定
圆周上的信息经 Cauchy 核送到任意内点。一个从局部差商出发消灭闭路积分，另一个从闭路积分恢复
局部导数。

\begin{figure}[htbp]
\centering
\begin{tikzpicture}[scale=0.96]
  \begin{scope}[xshift=0cm]
    \coordinate (A) at (0,0);
    \coordinate (B) at (3.4,0);
    \coordinate (C) at (0.7,2.6);
    \coordinate (AB) at (1.7,0);
    \coordinate (AC) at (0.35,1.3);
    \coordinate (BC) at (2.05,1.3);
    \fill[BookAccent!5] (A)--(B)--(C)--cycle;
    \draw[book line] (A)--(B)--(C)--cycle;
    \draw[BookAccent,semithick] (AB)--(AC)--(BC)--cycle;
    \draw[BookAccent,semithick] (AB)--(BC);
    \draw[BookAccent,semithick] (AC)--(BC);
    \coordinate (AAB) at (0.85,0);
    \coordinate (AAC) at (0.175,0.65);
    \coordinate (AACAB) at (1.025,0.65);
    \fill[BookWarm!18] (A)--(AAB)--(AAC)--cycle;
    \draw[BookWarm,semithick] (AAB)--(AAC)--(AACAB)--cycle;
    \fill[BookInk] (0.22,0.28) circle (1.4pt);
    \node[book label] at (0.48,0.28) {$z_*$};
    \node[book label] at (1.7,-0.38) {$T_0\supset T_1\supset T_2\supset\cdots$};
    \node[book label] at (1.7,-0.78) {$\operatorname{per}(T_n)=2^{-n}\operatorname{per}(T_0)$};
  \end{scope}
  \begin{scope}[xshift=7.0cm,yshift=0.15cm]
    \draw[BookAccent,semithick,fill=BookAccent!5] (0,1.05) circle (1.55);
    \fill[BookInk] (0,1.05) circle (1.5pt);
    \node[book label,below left=1mm and 1mm] at (0,1.05) {$a$};
    \fill[BookWarm] (0.55,1.45) circle (1.7pt);
    \node[book label,above right=0mm and 1mm] at (0.55,1.45) {$w$};
    \draw[book accent arrow] (-1.35,1.82) to[bend right=15]
      node[above,book label]{$\dfrac{f(z)}{z-w}\dd z$} (0.45,1.48);
    \draw[book arrow] (1.35,0.28) to[bend left=15]
      node[below,book label]{$f^{(n)}(w)$} (0.58,1.38);
    \node[book label] at (0,-0.88) {圆周观测 $\longrightarrow$ 内部全部导数};
  \end{scope}
\end{tikzpicture}
\caption{Goursat 的三角形细分与 Cauchy 核的边界恢复。左图中的积分尺度与周长平方同阶，而不是
与周长一次方同阶。}
\label{fig:7-1-2-goursat-and-cauchy}
\end{figure}

\begin{remark}[三项不能省略的边界]
闭曲线积分为零需要原函数或适当域条件；$1/z$ 说明有孔域上局部原函数不等于全局原函数。Morera
定理把连续性列为假设，若只给几乎处处可积性，则需另行证明一个弱版本。恒等定理中的聚点必须位于
连通域内部；零点只在边界聚集并不迫使函数恒为零，例如单位圆盘中的 $f(z)=\sin(1/(1-z))$ 在
$z=1$ 附近有无穷多个域内零点，但这些零点在域内没有聚点。
\end{remark}

\begin{exercise}
\begin{enumerate}
  \item 对整数 $m$ 计算 $\int_{|z|=r}z^m\dd z$，分别处理 $m=-1$ 与 $m\ne-1$。
  \item 从 Cauchy 导数公式重新推导式~\eqref{eq:7-1-2-cauchy-estimate}，并证明若 $d\ge0$ 且整函数满足
        $|f(z)|\le C(1+|z|)^d$，则它是次数不超过 $\lfloor d\rfloor$ 的多项式。
  \item 设 $f_n\in\mathcal O(\Omega)$ 局部一致收敛到 $f$。不用差商，逐步写出“连续性、三角形
        积分、Morera、Cauchy 公式”四个环节，证明 $f_n'\to f'$ 局部一致。
  \item 补全 Montel 证明的对角步骤：证明式~\eqref{eq:7-1-2-compact-exhaustion} 中每个紧集
        $K\Subset\Omega$ 都包含在某个 $K_m$ 中，并验证各层一致极限在交叠处相同。
\end{enumerate}
\end{exercise}

至此，局部一致极限、函数族紧性和唯一性都已获得完整证明。下面把这些工具用于含测度参数的核
$(z-x)^{-1}$，并比较同一几何级数在测度变换与有界算子预解式中的两种形态。


\subsection{参数全纯积分、Stieltjes 变换与预解式}\label{subsec:7-1-3}
\index{参数全纯积分}
\index{Cauchy--Stieltjes 变换}
\index{Stieltjes 反演}
\index{预解集}
\index{预解式恒等式}

核 $(z-x)^{-1}$ 同时出现在测度的 Cauchy--Stieltjes 变换与线性算子的预解式中。离开测度支集或
算子谱以后，分母有统一下界，于是积分、级数和导数可以交换；在无穷远处展开几何级数，又会把变换
的解析数据变成矩或算子幂。这个平行关系不需要谱定理，只需要有限复测度、Bochner 积分和有界算子
的基本代数。

Stieltjes 最初在连分数和矩问题中使用这类变换；Fredholm 的积分方程则把参数逆算子推向现代谱
理论。Dirichlet 问题、Fredholm 方程和 Fourier 无穷矩阵在历史上共同推动了 Hilbert 空间语言的形成。
这些联系在有限复测度、Bochner 积分和有界算子的层面即可建立；自伴算子与谱测度属于后续的算子理论。

\begin{example}[Dirac 与有限原子]\label{ex:7-1-3-1}
对 $a\in\R$，Dirac 测度满足
\[
  G_{\delta_a}(z)=\int_\R\frac1{z-x}\delta_a(\dd x)=\frac1{z-a}.
\]
更一般地，若 $a_1,\ldots,a_m$ 互异、$c_j\ne0$，且
$\mu=\sum_{j=1}^mc_j\delta_{a_j}$，则
\[
  G_\mu(z)=\sum_{j=1}^m\frac{c_j}{z-a_j}.
\]
这是一个有理函数；$a_j$ 是极点，而
\[
 \lim_{z\to a_j}(z-a_j)G_\mu(z)=c_j.
\]
所以极点位置记录原子位置，留数记录复质量。离散谱产生有理预解式正是同一代数图景。
\end{example}

\begin{example}[无穷远矩展开]\label{ex:7-1-3-2}
设 $R\ge0$，$\mu$ 为支撑包含于 $[-R,R]$ 的有限复测度，并记
$m_n=\int x^n\mu(\dd x)$。当 $|z|>R$ 时，
\[
 \frac1{z-x}=\frac1z\sum_{n=0}^{\infty}\left(\frac{x}{z}\right)^n
\]
关于 $x\in[-R,R]$ 一致收敛。因此
\[
 G_\mu(z)=\sum_{n=0}^{\infty}\frac{m_n}{z^{n+1}}.
 \tag{7.1.15}\label{eq:7-1-3-moment-expansion}
\]
截到第 $N$ 项的误差可直接估计为
\[
 \left|G_\mu(z)-\sum_{n=0}^{N}\frac{m_n}{z^{n+1}}\right|
 \le \frac{\|\mu\|_{TV}R^{N+1}}
 {|z|^{N+2}(1-R/|z|)}.
 \tag{7.1.16}\label{eq:7-1-3-moment-remainder}
\]
若支集无界而高阶矩不存在，这个展开便没有意义；不能仅凭测度有限就把它写在无穷远处。
\end{example}

\begin{example}[Neumann 预解式]\label{ex:7-1-3-3}
设 $T$ 是复 Banach 空间 $B$ 上的有界算子。若 $|z|>\|T\|$，则算子范数级数
\[
 S(z)=\frac1z\sum_{n=0}^{\infty}\left(\frac{T}{z}\right)^n
\]
绝对收敛。有限和满足
\[
 (zI-T)\frac1z\sum_{n=0}^{N}\left(\frac{T}{z}\right)^n
 =I-\left(\frac{T}{z}\right)^{N+1},
\]
右乘也有同一等式。令 $N\to\infty$，得到
\[
 (zI-T)^{-1}=\frac1z\sum_{n=0}^{\infty}\frac{T^n}{z^n},
 \qquad
 \|(zI-T)^{-1}\|\le\frac1{|z|-\|T\|}.
 \tag{7.1.17}\label{eq:7-1-3-neumann-resolvent}
\]
这与式~\eqref{eq:7-1-3-moment-expansion} 平行：标量矩 $m_n$ 被算子幂 $T^n$ 取代。
\end{example}

\begin{definition}[Cauchy--Stieltjes 变换]\label{def:7-1-3-1}
设 $\mu$ 为 $\R$ 上的有限复 Borel 测度。定义
\[
 G_\mu(z)=\int_\R\frac1{z-x}\mu(\dd x),
 \qquad z\in\C\setminus\supp|\mu|,
 \tag{7.1.18}\label{eq:7-1-3-stieltjes-definition}
\]
称为 $\mu$ 的 \emph{Cauchy--Stieltjes 变换}，也简称 Stieltjes 变换。本章固定核
$(z-x)^{-1}$；若采用 $(x-z)^{-1}$，所有涉及虚部的符号都要反向。
\end{definition}

\begin{definition}[预解集与预解式]\label{def:7-1-3-2}
设 $T\in\mathcal L(B)$ 是复 Banach 空间上的有界算子。若 $zI-T$ 有有界双边逆，则称
$z$ 属于 $T$ 的\emph{预解集} $\rho(T)$，并记
\[
  R(z,T)=(zI-T)^{-1},
\]
称其为 $T$ 在 $z$ 处的\emph{预解式}。补集 $\C\setminus\rho(T)$ 称为谱，但这里不使用任何
谱分解定理。
\end{definition}

\begin{definition}[局部一致支配]\label{def:7-1-3-3}
设 $\mu$ 是可测空间 $(X,\mathcal A)$ 上的有限复测度，并且对每个 $z\in\Omega$，截面
$x\mapsto q(z,x)$ 可测。若对每个紧集 $K\Subset\Omega$，存在
$g_K\in L^1(|\mu|)$，使得在一个可以随 $K$ 改变、但对所有 $z\in K$ 共同的 $|\mu|$-零集外，
\[
  |q(z,x)|\le g_K(x)\qquad(z\in K),
\]
则称参数族 $q$ \emph{局部一致受支配}。
\end{definition}

\begin{theorem}[参数积分全纯定理]\label{theorem:7-1-3-1}
设 $\mu$ 为可测空间 $(X,\mathcal A)$ 上的有限复测度，$q:\Omega\times X\to\C$ 满足：
\begin{enumerate}
  \item 对每个 $z\in\Omega$，$x\mapsto q(z,x)$ 可测；
  \item 存在 $N\in\mathcal A$、$|\mu|(N)=0$，使对每个 $x\in X\setminus N$，
        $z\mapsto q(z,x)$ 在 $\Omega$ 上全纯；
  \item $q$ 局部一致受支配。
\end{enumerate}
则
\[
  Q(z)=\int_Xq(z,x)\mu(\dd x)
\]
在 $\Omega$ 上全纯。若 $\partial_zq$ 也局部一致受某个 $L^1(|\mu|)$ 函数支配，则
\[
  Q'(z)=\int_X\partial_zq(z,x)\mu(\dd x).
  \tag{7.1.19}\label{eq:7-1-3-differentiate-under-integral}
\]
\end{theorem}

\begin{proof}
先在共同零集 $N$ 上把 $q$ 改定义为零。若 $\gamma:[a,b]\to\Omega$ 连续，则
$(t,x)\mapsto q(\gamma(t),x)$ 关于
$\Borel([a,b])\otimes\mathcal A$ 联合可测。事实上，把 $[a,b]$ 分成 $2^m$ 个等长区间，并令
$s_m(t)$ 为 $t$ 所在区间的左端点；函数
$q(\gamma(s_m(t)),x)$ 是有限个可测 $x$-截面乘区间指标之和。由 $s_m(t)\to t$ 及每个
$x$-截面的连续性，它们逐点趋于 $q(\gamma(t),x)$，故后者联合可测。这样，以下曲线积分与测度
积分的交换不仅有绝对可积估计，也满足 Fubini 所需的乘积可测性。

先证连续性。若 $z_n\to z$，可取紧圆盘 $K\Subset\Omega$ 包含 $z$ 及所有充分靠后的 $z_n$。
对几乎每个 $x$，全纯性给 $q(z_n,x)\to q(z,x)$；而
$|q(z_n,x)-q(z,x)|\le2g_K(x)$。对全变差测度应用支配收敛定理
\ref{theorem:3-1-3-5}，得到 $Q(z_n)\to Q(z)$。复平面为度量空间，所以 $Q$ 连续。

令 $\overline T\subset\Omega$ 为三角形，并取包含 $\partial T$ 的紧集 $K$。由支配条件，
\[
 \int_{\partial T}\int_X|q(z,x)|\,|\mu|(\dd x)|\dd z|
 \le\operatorname{per}(T)\int_Xg_K\dd|\mu|<\infty.
\]
Fubini 定理因而允许交换曲线参数积分与测度积分。对几乎每个 $x$ 使用 Cauchy--Goursat 定理，得
\[
 \int_{\partial T}Q(z)\dd z
 =\int_X\left(\int_{\partial T}q(z,x)\dd z\right)\mu(\dd x)=0.
\]
Morera 定理给出 $Q\in\mathcal O(\Omega)$。

最后固定 $z$。先取正实数列 $h_m\downarrow0$，使 $z+h_m\in\Omega$。对共同满测集中的每个
$x$，可测差商
\[
 x\longmapsto\frac{q(z+h_m,x)-q(z,x)}{h_m}
\]
逐点趋于 $\partial_zq(z,x)$；因此导数截面 $x\mapsto\partial_zq(z,x)$ 可测。再取闭小圆盘
$K\Subset\Omega$，使它包含所有 $z+th$，其中 $0\le t\le1$ 且 $h$ 充分小。对几乎每个 $x$，
沿复线段使用实变量微积分，
\[
 \frac{q(z+h,x)-q(z,x)}h
 =\int_0^1\partial_zq(z+th,x)\dd t.
\]
若 $|\partial_zq(\zeta,x)|\le h_K(x)$ 于 $\zeta\in K$，则差商由 $h_K\in L^1(|\mu|)$ 支配。
支配收敛把 $h\to0$ 送入积分号内，得到
\eqref{eq:7-1-3-differentiate-under-integral}。这样，可测性是在应用支配收敛之前由差商极限得到的，
没有把“导数可测”作为未写出的附加假设。
\end{proof}

\begin{proposition}[Stieltjes 变换的基本估计]\label{proposition:7-1-3-2}
设 $\mu$ 为 $\R$ 上的有限复 Borel 测度。则 $G_\mu$ 在
$\C\setminus\supp|\mu|$ 上全纯，并且
\[
 |G_\mu(z)|
 \le\frac{\|\mu\|_{TV}}{\dist(z,\supp|\mu|)}.
 \tag{7.1.20}\label{eq:7-1-3-stieltjes-bound}
\]
对每个 $n\ge0$，
\[
 G_\mu^{(n)}(z)
 =(-1)^nn!\int_\R\frac1{(z-x)^{n+1}}\mu(\dd x).
 \tag{7.1.21}\label{eq:7-1-3-stieltjes-derivatives}
\]
当 $\mu=0$ 时约定 $\supp|\mu|=\varnothing$，估计
\eqref{eq:7-1-3-stieltjes-bound} 的右端按 $0$ 理解。
\end{proposition}

\begin{proof}
若 $\mu=0$，结论直接成立。以下设 $\supp|\mu|\ne\varnothing$。固定紧集
$K\Subset\C\setminus\supp|\mu|$，并令
\[
 d_K=\dist(K,\supp|\mu|)>0.
\]
由于 $|\mu|(\R\setminus\supp|\mu|)=0$，对 $z\in K$ 及几乎每个 $x$ 有
\[
 \left|\frac1{z-x}\right|\le d_K^{-1},
 \qquad
 \left|\frac{\partial^n}{\partial z^n}\frac1{z-x}\right|
 \le n!d_K^{-n-1}.
\]
这些常数函数属于 $L^1(|\mu|)$。定理~\ref{theorem:7-1-3-1} 可反复应用，给出全纯性和
式~\eqref{eq:7-1-3-stieltjes-derivatives}。在固定 $z$ 处直接对核取绝对值并积分，则得到
式~\eqref{eq:7-1-3-stieltjes-bound}。
\end{proof}

\begin{theorem}[Stieltjes 唯一性基础]\label{theorem:7-1-3-3}
记 $\C_+=\{z:\operatorname{Im}z>0\}$、
$\C_-=\{z:\operatorname{Im}z<0\}$。
\begin{enumerate}
  \item 若 $\mu,\nu$ 是 $\R$ 上的有限实有符号 Borel 测度，且
        $G_\mu=G_\nu$ 于 $\C_+$，则 $\mu=\nu$。
  \item 若 $\mu,\nu$ 是有限复 Borel 测度，且
        $G_\mu=G_\nu$ 同时在 $\C_+$ 与 $\C_-$ 成立，则 $\mu=\nu$。
\end{enumerate}
对有限实有符号测度和任意 $a<b$，还有边界反演式
\[
 -\frac1\pi\lim_{\varepsilon\downarrow0}
 \int_a^b\operatorname{Im}G_\mu(t+i\varepsilon)\dd t
 =\mu((a,b))+\frac12\mu(\{a\})+\frac12\mu(\{b\}).
 \tag{7.1.22}\label{eq:7-1-3-stieltjes-inversion}
\]
特别地，若 $a,b$ 都不是原子，则左侧恢复 $\mu((a,b))$。
\end{theorem}

\begin{proof}
令
\[
 P_\varepsilon(u)=\frac1\pi\frac{\varepsilon}{u^2+\varepsilon^2},
 \qquad \varepsilon>0.
\]
它非负且积分为 $1$。若 $\sigma$ 是有限实有符号测度，则实线性允许把虚部送入积分，因而
\[
 -\frac1\pi\operatorname{Im}G_\sigma(t+i\varepsilon)
 =\int_\R P_\varepsilon(t-x)\sigma(\dd x)
 =:(P_\varepsilon*\sigma)(t).
 \tag{7.1.23}\label{eq:7-1-3-poisson-boundary}
\]
若 $G_\sigma=0$ 于 $\C_+$，则 $P_\varepsilon*\sigma=0$ 对所有 $\varepsilon>0$ 成立。

取 $\varphi\in C_c(\R)$。Fubini 定理给出
\[
 0=\int_\R\varphi(t)(P_\varepsilon*\sigma)(t)\dd t
  =\int_\R(P_\varepsilon*\varphi)(x)\sigma(\dd x),
 \tag{7.1.24}\label{eq:7-1-3-poisson-test}
\]
其中利用了 $P_\varepsilon(-u)=P_\varepsilon(u)$。Poisson 核是近似恒等。为完整核验这一点，令
$\omega_\varphi(\delta)=\sup_{|u-v|\le\delta}|\varphi(u)-\varphi(v)|$。则
\begin{align*}
 \sup_x|(P_\varepsilon*\varphi)(x)-\varphi(x)|
 &\le \omega_\varphi(\delta)
 +2\|\varphi\|_\infty
   \int_{|u|>\delta}P_\varepsilon(u)\dd u.
\end{align*}
先固定小的 $\delta$，再令 $\varepsilon\downarrow0$；尾积分趋于零，而
$\omega_\varphi(\delta)\to0$。所以 $P_\varepsilon*\varphi\to\varphi$ 一致。让
$\varepsilon\downarrow0$ 于式~\eqref{eq:7-1-3-poisson-test}，得到
$\int\varphi\dd\sigma=0$ 对每个 $\varphi\in C_c(\R)$ 成立。对 $\sigma$ 的 Jordan 正负部使用
定理~\ref{theorem:3-1-1-4} 的 $C_c$ 测试恢复，便有 $\sigma=0$。取
$\sigma=\mu-\nu$，即证明第一项。

现在令 $\sigma$ 为有限复测度。直接作核的上下边界差，得到
\begin{align*}
 G_\sigma(t-i\varepsilon)-G_\sigma(t+i\varepsilon)
 &=\int_\R\left(\frac1{t-i\varepsilon-x}
                 -\frac1{t+i\varepsilon-x}\right)\sigma(\dd x)\\
 &=2\pi i\,(P_\varepsilon*\sigma)(t).
 \tag{7.1.25}\label{eq:7-1-3-two-sided-jump}
\end{align*}
若 $G_\sigma$ 在两个半平面都为零，则 $P_\varepsilon*\sigma=0$。仍用
式~\eqref{eq:7-1-3-poisson-test} 和一致近似，得到
$\int\varphi\dd\sigma=0$。分别取实部和虚部，归结为两个有限实有符号测度，故 $\sigma=0$。

最后证明反演式。由式~\eqref{eq:7-1-3-poisson-boundary} 与 Fubini，
\[
 -\frac1\pi\int_a^b\operatorname{Im}G_\mu(t+i\varepsilon)\dd t
 =\int_\R H_\varepsilon(x)\mu(\dd x),
\]
其中
\[
 H_\varepsilon(x)
 =\int_a^bP_\varepsilon(t-x)\dd t
 =\frac1\pi\left[
   \arctan\frac{b-x}{\varepsilon}
  -\arctan\frac{a-x}{\varepsilon}\right].
\]
有 $0\le H_\varepsilon\le1$，且当 $\varepsilon\downarrow0$ 时，它逐点趋于
$\1_{(a,b)}+\tfrac12\1_{\{a\}}+\tfrac12\1_{\{b\}}$。对全变差测度使用支配收敛，便得到
式~\eqref{eq:7-1-3-stieltjes-inversion}。
\end{proof}

\begin{remark}[复测度为何需要另一半边界数据]
实有符号测度满足
$G_\mu(\bar z)=\overline{G_\mu(z)}$，所以上半平面自动决定下半平面。复测度没有这个对称性，而且
仅给上半平面数据确实可能丢失非零测度。令
\[
  \mu(\dd x)=\frac1{(x-i)^2}\dd x.
\]
因为 $|(x-i)^{-2}|=(x^2+1)^{-1}$，这是非零有限复测度。固定 $z\in\C_+$，并取
$R>2\max\{|z|,1\}$。把
\[
 \zeta\longmapsto\frac1{(z-\zeta)(\zeta-i)^2}
\]
沿实线段 $[-R,R]$ 与从 $R$ 回到 $-R$ 的下半圆组成的闭曲线积分。令
$\eta=\tfrac12\min\{\operatorname{Im}z,1\}$；凸开集
\[
 U_R=\{\zeta:|\zeta|<R+1,\ \operatorname{Im}\zeta<\eta\}
\]
包含该闭曲线而排除两个极点 $z,i$，所以被积函数在星形域 $U_R$ 全纯，Cauchy--Goursat 定理给
闭路积分为零。在下半圆弧上，$|z-\zeta|\ge R/2$、$|\zeta-i|\ge R/2$，故被积函数绝对值不超过
$8R^{-3}$；弧长为 $\pi R$，所以弧积分绝对值不超过 $8\pi R^{-2}$，随 $R\to\infty$ 趋于零。因此
\[
 G_\mu(z)=\int_\R\frac{1}{(z-x)(x-i)^2}\dd x=0
 \qquad(z\in\C_+),
\]
尽管 $\mu\ne0$。这说明第二项中的双半平面假设不是实测度证明里可随意删除的装饰。
\end{remark}

\begin{proposition}[预解式恒等式]\label{proposition:7-1-3-4}
设 $T\in\mathcal L(B)$。对任意 $z,w\in\rho(T)$，
\[
 R(z,T)-R(w,T)=(w-z)R(z,T)R(w,T).
 \tag{7.1.26}\label{eq:7-1-3-resolvent-identity}
\]
预解集 $\rho(T)$ 是开集，$z\mapsto R(z,T)$ 在其上按算子范数全纯，并满足
\[
  \frac{\dd}{\dd z}R(z,T)=-R(z,T)^2.
 \tag{7.1.27}\label{eq:7-1-3-resolvent-derivative}
\]
此外，$\{|z|>\|T\|\}\subset\rho(T)$，且该区域上的预解式由
式~\eqref{eq:7-1-3-neumann-resolvent} 给出。
\end{proposition}

\begin{proof}
一般可逆算子 $A,B$ 满足
\[
 A^{-1}-B^{-1}=A^{-1}(B-A)B^{-1},
\]
因为右侧展开为 $A^{-1}BB^{-1}-A^{-1}AB^{-1}$。取
$A=zI-T$、$B=wI-T$，便得到
式~\eqref{eq:7-1-3-resolvent-identity}。这也顺带说明不同参数的预解式彼此交换。

固定 $z_0\in\rho(T)$，记 $R_0=R(z_0,T)$。因
\[
 zI-T=(z_0I-T)\bigl[I+(z-z_0)R_0\bigr],
\]
当 $|z-z_0|\|R_0\|<1$ 时，Neumann 级数给出
\[
 R(z,T)
 =\bigl[I+(z-z_0)R_0\bigr]^{-1}R_0
 =\sum_{n=0}^{\infty}(-1)^n(z-z_0)^nR_0^{n+1},
 \tag{7.1.28}\label{eq:7-1-3-local-resolvent-series}
\]
并且级数按算子范数收敛。这同时证明 $\rho(T)$ 开及 $R(\,\cdot\,,T)$ 的算子范数全纯性。逐项微分
式~\eqref{eq:7-1-3-local-resolvent-series}，或令 $w\to z$ 于
式~\eqref{eq:7-1-3-resolvent-identity}，得到
式~\eqref{eq:7-1-3-resolvent-derivative}。最后，$|z|>\|T\|$ 时的双边逆及范数收敛已在
例~\ref{ex:7-1-3-3} 中逐项核验。
\end{proof}

图~\ref{fig:7-1-3-kernel-and-resolvent} 把两条计算链并排放置。对测度而言，距离支集的正下界产生
积分支配，无穷远展开的系数是矩；对算子而言，$|z|>\|T\|$ 产生范数收敛，展开系数是 $T^n$。
Stieltjes 边界值还可恢复测度，而预解式恒等式则控制参数改变时逆算子的变化。

\begin{figure}[htbp]
\centering
\begin{tikzpicture}[node distance=12mm and 15mm,scale=0.94,transform shape]
  \begin{scope}[xshift=0cm]
    \draw[book line,-{Stealth[length=2mm]}] (-2.0,0)--(2.2,0) node[right,book label]{$\R$};
    \draw[BookWarm,very thick] (-1.25,0)--(1.15,0);
    \foreach \x in {-1.05,-0.45,0.15,0.85}
      \fill[BookWarm] (\x,0) circle (1.6pt);
    \fill[BookAccent] (0.55,1.55) circle (2pt);
    \node[book label,above right=0mm and 1mm] at (0.55,1.55) {$z$};
    \draw[book accent arrow] (0.5,1.45)--node[left,book label]{$z-x$}(0.05,0.16);
    \node[book label] at (0,-0.48) {$\supp|\mu|$};
    \node[book warm node] at (0,-1.25) {$G_\mu(z)=\int(z-x)^{-1}\dd\mu(x)$};
  \end{scope}
  \node[book strong node] (mom) at (5.05,0.85)
    {无穷远几何展开\\[1mm]$G_\mu(z)=\sum m_nz^{-n-1}$};
  \node[book node] (bound) at (5.05,-0.7)
    {边界逼近\\$P_\varepsilon*\mu\to\mu$};
  \draw[book accent arrow] (2.25,0.9)--(mom.west);
  \draw[book arrow] (2.25,-0.15)--(bound.west);
  \node[book strong node] (T) at (9.2,0.85)
    {有界算子 $T$\\[1mm]$R(z,T)=\sum T^nz^{-n-1}$};
  \node[book node] (id) at (9.2,-0.7)
    {参数变化\\$R(z)-R(w)=(w-z)R(z)R(w)$};
  \draw[book accent arrow] (mom.east)--node[above,book label]{矩 $m_n\leftrightarrow T^n$}(T.west);
  \draw[book arrow] (bound.east)--(id.west);
\end{tikzpicture}
\caption{同一核机制的测度版与算子版。离开支集或在算子范数意义下远离原点，都会产生可逐项运算的
几何级数；边界恢复与预解式恒等式则分别编码唯一性和参数稳定性。}
\label{fig:7-1-3-kernel-and-resolvent}
\end{figure}

这些公式也解释了前面若干对象之间的联系。半群的 Laplace--Bochner 积分在
命题~\ref{proposition:4-4-3-4} 中已经给出生成元的预解式；这里的参数全纯定理可进一步控制复参数。
概率分布的 Stieltjes 变换通过定理~\ref{theorem:7-1-3-3} 唯一识别其分布律，并可用 Poisson 核边界
逼近连续测试积分。Fredholm 方程中的参数逆算子则服从
式~\eqref{eq:7-1-3-resolvent-identity}，无需先有谱分解就能比较两个参数处的解。

\begin{remark}[三项适用范围]
核的符号约定决定式~\eqref{eq:7-1-3-poisson-boundary} 的正负号；改用 $(x-z)^{-1}$ 时必须整体
改号。矩展开需要有界支集，或至少需要足够多矩并另证余项控制；有限测度本身不保证全部矩存在。
预解式的定义要求 $zI-T$ 有有界双边逆；本小节只处理处处定义的有界算子，未把结论扩大到无界
生成元的一般谱理论。
\end{remark}

\begin{exercise}
\begin{enumerate}
  \item 设 $\mu=2\delta_{-1}-(1+i)\delta_2$。计算 $G_\mu$、全部极点及其留数，并写出
        $|z|>2$ 时式~\eqref{eq:7-1-3-moment-expansion} 的前三项。
  \item 从有限几何和出发证明式~\eqref{eq:7-1-3-moment-remainder}，并指出当 $R=0$ 时公式如何
        退化。
  \item 对有限实有符号测度，逐项计算
        $H_\varepsilon(x)=\int_a^bP_\varepsilon(t-x)\dd t$ 的极限，重新证明
        式~\eqref{eq:7-1-3-stieltjes-inversion}；解释端点原子为何只贡献一半。
  \item 由式~\eqref{eq:7-1-3-resolvent-identity} 证明
        $R(z,T)R(w,T)=R(w,T)R(z,T)$，再由局部 Neumann 级数
        \eqref{eq:7-1-3-local-resolvent-series} 计算前两阶导数。
\end{enumerate}
\end{exercise}

下一节转向圆周与欧氏空间上的 Fourier 分析。复指数将把平移和卷积变成乘法，而三角矩问题会把
这里的“变换识别测度”改写成一族离散频率观测能否确定整体分布的问题。

\section{Fourier 分析：从三角矩到热核}\label{sec:07-02}

\subsection{圆周 Fourier 级数、\Fejer{} 逼近与 Plancherel}\label{subsec:7-2-1}
\index{圆周 Fourier 级数}\index{Fejer@\Fejer{} 核}\index{Plancherel 定理}

把圆周写成 $\T=\R/2\pi\Z$，并用
\[
  \dd m(x)=\frac{\dd x}{2\pi}
\]
表示归一化 Haar 测度。指数函数 $e_n(x)=e^{inx}$ 满足
$\int_\T e_n\overline{e_k}\,\dd m=\1_{\{n=k\}}$。因此，若只观测函数对所有
$e_n$ 的积分，所得数列正是函数关于这族正交方向的坐标。困难不在于坐标是否存在，而在于这些坐标
能否把原函数恢复出来。

最先想到的恢复算子是对称部分和
\[
  S_Nf(x)=\sum_{|n|\leq N}\widehat f(n)e^{inx}.
\]
下一个例子说明，不能把有限维正交投影的直觉直接升级为逐点恢复。

\begin{example}[Gibbs 现象与 Dirichlet 部分和的失效]\label{ex:7-2-1-1}
令 $f$ 是圆周上的方波，取代表
\[
 f(x)=
 \begin{cases}
   -1,&-\pi<x<0,\\
   0,&x=0\text{ 或 }x=\pi\equiv-\pi,\\
   1,&0<x<\pi.
 \end{cases}
\]
直接积分得到：当 $n\neq0$ 时
\[
 \widehat f(n)
 =\frac1{2\pi}\left(\int_0^\pi-\int_{-\pi}^0\right)e^{-inx}\dd x
 =\frac{1-(-1)^n}{\pi i n}.
\]
于是，以奇数频率 $2N+1$ 截断时，
\begin{equation}\label{eq:7-2-gibbs-partial-sum}
 S_{2N+1}f(x)=\frac4\pi\sum_{k=0}^N\frac{\sin((2k+1)x)}{2k+1}.
\end{equation}
对右端求导，并使用有限等比和，得到
\[
 \frac{\dd}{\dd x}S_{2N+1}f(x)
 =\frac4\pi\sum_{k=0}^N\cos((2k+1)x)
 =\frac2\pi\frac{\sin(2(N+1)x)}{\sin x}.
\]
取 $x_N=\pi/(2(N+1))$，由 $S_{2N+1}f(0)=0$ 及换元
$u=2(N+1)x$，有
\[
 S_{2N+1}f(x_N)
 =\frac2\pi\int_0^\pi
   \frac{\sin u}{2(N+1)\sin(u/(2(N+1)))}\dd u
 \longrightarrow \frac2\pi\int_0^\pi\frac{\sin u}{u}\dd u.
\]
这里的支配可取常数倍的 $\sin u/u$：在 $0\leq v\leq\pi/2$ 上
$\sin v\geq 2v/\pi$。极限严格大于 $1$。例如把
$\sin u/u$ 展开并使用交错级数余项，
\[
 \int_0^\pi\frac{\sin u}{u}\dd u
 >\pi-\frac{\pi^3}{18}+\frac{\pi^5}{600}
       -\frac{\pi^7}{35280}>\frac\pi2.
\]
数值上该极限约为 $1.17898$。因此观察点虽趋向跳点，过冲却不趋于零；这就是 Gibbs 现象。

这种失效不只来自跳跃。Dirichlet 核
\[
 D_N(x)=\sum_{|n|\leq N}e^{inx}
       =\frac{\sin((N+\tfrac12)x)}{\sin(x/2)}
\]
的 $L^1(m)$ 范数按 $\log N$ 增长。下面给出带显式区间的下界。对
$1\leq k\leq\lfloor N/2\rfloor$，置
\[
 I_{k,N}=\left[
   \frac{(k+1/3)\pi}{N+1/2},
   \frac{(k+2/3)\pi}{N+1/2}
 \right].
\]
这些区间两两不交。若 $x\in I_{k,N}$，则
\[
 \left|\sin\bigl((N+\tfrac12)x\bigr)\right|\geq\frac{\sqrt3}{2},
 \qquad 0<\sin(x/2)\leq\frac{x}{2},
\]
因而 $|D_N(x)|\geq\sqrt3/x$。使用归一化测度 $\dd m=\dd x/(2\pi)$，逐段得到
\begin{align*}
 \int_{I_{k,N}}|D_N(x)|\dd m(x)
 &\geq\frac{\sqrt3}{2\pi}
       \frac{|I_{k,N}|}{\sup I_{k,N}}\\
 &=\frac{\sqrt3}{6\pi(k+2/3)}
 \geq\frac{\sqrt3}{10\pi}\frac1k.
\end{align*}
因此
\[
 \|D_N\|_{L^1(m)}
 \geq\frac{\sqrt3}{10\pi}
       \sum_{k=1}^{\lfloor N/2\rfloor}\frac1k
 \geq c'\log N
\]
（$N$ 充分大），其中 $c'>0$ 与 $N$ 无关。

令 $L_N:C(\T)\to\C$ 为 $L_Ng=S_Ng(0)=\int g(-y)D_N(y)\dd m(y)$。
三角不等式先给 $\|L_N\|\leq\|D_N\|_1$。反向固定 $\eta>0$。$D_N$ 只有有限多个零点，可在这些
零点周围取有限个小弧，其并记为 $O$，并使
$\int_O|D_N|\dd m<\eta$。在补集上令
$g(-y)=\operatorname{sgn}D_N(y)$，再在每个小弧上作模不超过 $1$ 的线性连接，得到
$g\in C(\T)$、$\|g\|_\infty\leq1$。于是
\[
 |L_Ng|
 \geq\int_{O^c}|D_N|\dd m-\int_O|D_N|\dd m
 \geq\|D_N\|_1-2\eta.
\]
令 $\eta\downarrow0$，得到 $\|L_N\|=\|D_N\|_1$。若每个 $g\in C(\T)$ 的数列
$(L_Ng)$ 都有界，则闭集
\[
 E_r=\{g:\sup_N|L_Ng|\leq r\}\qquad(r\in\N)
\]
覆盖完备空间 $C(\T)$。Baire 定理给出某个球
$B(g_0,\rho)\subset E_r$。对 $\|h\|_\infty\leq1$，把
$g_0+(\rho/2)h$ 与 $g_0$ 相减，便得
$\sup_N|L_Nh|\leq4r/\rho$，这与 $\|L_N\|\to\infty$ 矛盾。
所以存在连续函数 $g$，其 Fourier 部分和甚至在固定点 $0$ 发散。
\end{example}

Dirichlet 核的问题是绝对质量越来越大而且正负振荡。\Cesaro{} 平均把前 $N+1$ 个部分和再平均；
这看似只是数列求和法的改变，却恰好把振荡核换成了正核。

\begin{example}[\Fejer{} 正核]\label{ex:7-2-1-2}
有限和的次数可以交换，因此
\begin{align*}
 F_N(x)
 &=\frac1{N+1}\sum_{k=0}^ND_k(x)\\
 &=\sum_{|n|\leq N}\left(1-\frac{|n|}{N+1}\right)e^{inx}\\
 &=\frac1{N+1}\left|\sum_{k=0}^Ne^{ikx}\right|^2
 =\frac1{N+1}\left(\frac{\sin((N+1)x/2)}{\sin(x/2)}\right)^2.
\end{align*}
在 $x\in2\pi\Z$ 处取连续延拓值 $N+1$。平方表示说明 $F_N\geq0$；常数 Fourier
系数为 $1$，故 $\int_\T F_N\dd m=1$。固定 $0<\delta<\pi$ 后，若
$\delta\leq|x|\leq\pi$，则
\[
 F_N(x)\leq\frac1{(N+1)\sin^2(\delta/2)},
 \qquad
 \int_{\delta\leq|x|\leq\pi}F_N(x)\dd m(x)\longrightarrow0.
\]
所以全部质量集中到圆周原点。并且若 $f\geq0$ a.e.，则
$f*F_N\geq0$；这项保正性正是 Dirichlet 部分和所没有的稳定性。
\end{example}

Fourier 系数也可以对测度定义。此时系数的高频行为能排除某些对象，却不能单独确定测度的全部
几何性质。

\begin{example}[原子测度]\label{ex:7-2-1-3}
对 $x_0\in\T$，
\[
 \widehat{\delta_{x_0}}(n)
 =\int_\T e^{-inx}\dd\delta_{x_0}(x)=e^{-inx_0},
 \qquad |\widehat{\delta_{x_0}}(n)|=1.
\]
因此原子测度的 Fourier 系数不趋零。推论~\ref{corollary:7-2-1-2} 与圆周上的
Riemann--Lebesgue 结论还将说明：这种测度不可能有 $L^1(m)$ 密度。

反过来，“系数不趋零”并不推出存在原子。令
\[
 X=2\pi\sum_{j=1}^\infty\frac{2\varepsilon_j}{3^j},
 \qquad \Prob(\varepsilon_j=0)=\Prob(\varepsilon_j=1)=\frac12,
\]
并令 $\mu$ 为 $X$ 在圆周上的分布。这是无原子的 Cantor 型奇异测度，而独立性给出
\[
 |\widehat\mu(n)|
 =\prod_{j=1}^\infty\left|\cos\frac{2\pi n}{3^j}\right|.
\]
取 $n=3^k$，前 $k$ 个因子绝对值为 $1$，故
\[
 |\widehat\mu(3^k)|
 =\prod_{r=1}^\infty\left|\cos\frac{2\pi}{3^r}\right|>0.
\]
先核对“无原子”。一个 Cantor 点至多对应两个只含数字 $0,2$ 的三进制展开；任一指定数字串的前
$N$ 位同时出现的概率为 $2^{-N}$，故该完整数字串的概率至多
$\lim_{N\to\infty}2^{-N}=0$。因此对每个 $x$，
$\mu(\{x\})\leq2\lim_N2^{-N}=0$。

最后的无穷乘积为正也可逐项核验。令 $\theta_r=2\pi/3^r$。当 $r\geq2$ 时
$0<\cos\theta_r<1$，并且
\[
 u_r:=1-\cos\theta_r\leq\frac{\theta_r^2}{2},
 \qquad \sum_{r=2}^\infty u_r<\infty.
\]
又有 $u_r<1/2$，所以 $-\log(1-u_r)\leq2u_r$；从而
$\sum_{r\geq2}-\log(\cos\theta_r)<\infty$，即
$\prod_{r\geq2}\cos\theta_r>0$。乘上第一个因子的绝对值
$|\cos(2\pi/3)|=1/2$，便得所写无穷乘积严格为正。这给出非原子而 Fourier 系数不衰减的具体例子。
\end{example}

\begin{definition}[Fourier 系数]\label{def:7-2-1-1}
若 $f\in L^1(\T,m)$，定义
\[
 \widehat f(n)=\int_\T f(x)e^{-inx}\dd m(x),\qquad n\in\Z.
\]
若 $\mu$ 是 $\T$ 上有限复 Borel 测度，则定义
$\widehat\mu(n)=\int_\T e^{-inx}\dd\mu(x)$。对可积函数，圆周卷积约定为
\[
 (f*g)(x)=\int_\T f(x-y)g(y)\dd m(y).
\]
在 $L^2(\T)$ 中取内积 $\langle f,g\rangle=\int f\overline g\dd m$，于是
$\widehat f(n)=\langle f,e_n\rangle$。
\end{definition}

\begin{definition}[Dirichlet 与 \Fejer{} 核]\label{def:7-2-1-2}
对 $N\geq0$，定义
\[
 D_N=\sum_{|n|\leq N}e_n,
 \qquad
 F_N=\frac1{N+1}\sum_{k=0}^ND_k.
\]
相应的 Fourier 部分和与 \Fejer{} 平均为
\[
 S_Nf=f*D_N,
 \qquad
 \sigma_Nf=\frac1{N+1}\sum_{k=0}^NS_kf=f*F_N.
\]
\end{definition}

例~\ref{ex:7-2-1-2} 已证明 $F_N\geq0$、$\int F_N\dd m=1$，且其远离零点的质量趋于零；
即 $(F_N)$ 是圆周上的正近似恒等族。

\begin{definition}[三角矩问题]\label{def:7-2-1-3}
给定复数列 $(c_n)_{n\in\Z}$，三角矩问题是寻找函数 $f$ 或有限测度 $\mu$，使
\[
 c_n=\int_\T e^{-inx}f(x)\dd m(x)
 \quad\text{或}\quad
 c_n=\int_\T e^{-inx}\dd\mu(x)
 \qquad(n\in\Z).
\]
\end{definition}

\begin{proposition}[正定性的必要性]\label{proposition:7-2-1-positive-definite-necessity}
若正测度 $\mu$ 解三角矩问题，则对任意 $r\geq1$、整数 $n_1,\ldots,n_r$ 与复数
$z_1,\ldots,z_r$，都有
\[
 \sum_{j,k=1}^r c_{n_j-n_k}z_j\overline{z_k}\geq0.
\]
\end{proposition}

\begin{proof}
把矩表示代入有限和，得到
\[
 \sum_{j,k=1}^r c_{n_j-n_k}z_j\overline{z_k}
 =\int_\T\left|\sum_{j=1}^rz_je^{-in_jx}\right|^2\dd\mu(x)\geq0,
\]
因为被积函数与测度均为非负。
\end{proof}

因此正定性是正测度三角矩的必要条件。第~\ref{subsec:7-3-1} 小节将证明相应的充分性；本小节先
解决 Fourier 系数所决定测度的唯一性。

\begin{theorem}[\Fejer{} 定理]\label{theorem:7-2-1-1}
若 $f\in C(\T)$，则 $\sigma_Nf\to f$ 一致收敛。若
$f\in L^p(\T)$ 且 $1\leq p<\infty$，则
$\sigma_Nf\to f$ 于 $L^p(\T)$。
\end{theorem}

\begin{proof}
先设 $f\in C(\T)$。由 $\int F_N\dd m=1$，
\begin{equation}\label{eq:7-2-fejer-difference}
 \sigma_Nf(x)-f(x)
 =\int_\T\bigl(f(x-y)-f(x)\bigr)F_N(y)\dd m(y).
\end{equation}
给定 $\varepsilon>0$，圆周紧性使 $f$ 一致连续；可取 $0<\delta<\pi$，使
$|f(x-y)-f(x)|<\varepsilon/2$ 对所有 $x$ 及 $|y|<\delta$ 成立。把积分分成近端与远端，并使用
$F_N\geq0$，得到
\begin{align*}
 \|\sigma_Nf-f\|_\infty
 &\leq \frac\varepsilon2\int_{|y|<\delta}F_N(y)\dd m(y)
 +2\|f\|_\infty\int_{|y|\geq\delta}F_N(y)\dd m(y)\\
 &\leq\frac\varepsilon2
 +\frac{2\|f\|_\infty}{(N+1)\sin^2(\delta/2)}.
\end{align*}
当 $N$ 足够大时第二项小于 $\varepsilon/2$，证明一致收敛。

再设 $f\in L^p$，$1\leq p<\infty$。积分型 Minkowski 不等式与平移等距性给出
\begin{equation}\label{eq:7-2-fejer-lp}
 \|\sigma_Nf-f\|_p
 \leq\int_\T\|f(\,\cdot-y)-f\|_pF_N(y)\dd m(y).
\end{equation}
这里平移在 $L^p(\T)$ 中强连续：先用连续函数 $g$ 逼近 $f$，再由
\[
 \|\tau_yf-f\|_p
 \leq2\|f-g\|_p+\|\tau_yg-g\|_p
 \leq2\|f-g\|_p+\|\tau_yg-g\|_\infty
\]
及 $g$ 的一致连续性令 $y\to0$。因此给定 $\varepsilon>0$，可取 $\delta>0$ 使
$\|f(\cdot-y)-f\|_p<\varepsilon/2$ 对 $|y|<\delta$ 成立。远端上该范数不超过
$2\|f\|_p$。于是与上面相同的分裂给出
\[
 \|\sigma_Nf-f\|_p
 \leq\frac\varepsilon2
 +2\|f\|_p\int_{|y|\geq\delta}F_N(y)\dd m(y)<\varepsilon
\]
，只要 $N$ 足够大。
\end{proof}

\begin{corollary}[Stone--Weierstrass 的 Fourier 版]\label{corollary:7-2-1-2}
三角多项式在 $C(\T)$ 中关于一致范数稠密，也在 $L^p(\T)$ 中关于 $L^p$ 范数稠密，
$1\leq p<\infty$。更具体地，$\sigma_Nf$ 就是一列只使用 $|n|\leq N$ 频率的逼近。
\end{corollary}

\begin{proof}
由例~\ref{ex:7-2-1-2} 的有限 Fourier 展开，
\[
 \sigma_Nf(x)
 =\sum_{|n|\leq N}\left(1-\frac{|n|}{N+1}\right)\widehat f(n)e^{inx},
\]
所以 $\sigma_Nf$ 是三角多项式。定理~\ref{theorem:7-2-1-1} 分别给出一致范数与
$L^p$ 范数收敛，因而给出两种稠密性。
\end{proof}

这里还得到圆周上的 Riemann--Lebesgue 结论。若 $f\in L^1(\T)$，给定
$\varepsilon>0$，由本推论取三角多项式 $p$ 使 $\|f-p\|_1<\varepsilon$。当 $|n|$ 超过
$p$ 的最高频率时，$\widehat p(n)=0$，所以
\[
 |\widehat f(n)|=|\widehat{f-p}(n)|\leq\|f-p\|_1<\varepsilon.
\]
因此 $\widehat f(n)\to0$ 当 $|n|\to\infty$。例~\ref{ex:7-2-1-3} 中的原子测度不可能有
$L^1$ 密度，正是这项结论的应用。

现在可以把 Bessel 的单向估计升级为 Parseval 等式。关键不是再作一次形式求和，而是利用刚刚得到的
稠密性证明指数族没有遗漏任何 $L^2$ 方向。

\begin{theorem}[$\T$ 上的 Plancherel/Parseval]\label{theorem:7-2-1-3}
Fourier 映射
\[
 \mathcal F_\T:L^2(\T)\longrightarrow\ell^2(\Z),
 \qquad f\longmapsto(\widehat f(n))_{n\in\Z},
\]
是满射线性等距同构。特别地，
\begin{equation}\label{eq:7-2-circle-parseval}
 \|f\|_2^2=\sum_{n\in\Z}|\widehat f(n)|^2,
 \qquad
 \langle f,g\rangle
 =\sum_{n\in\Z}\widehat f(n)\overline{\widehat g(n)}.
\end{equation}
并且对称部分和 $S_Nf=\sum_{|n|\leq N}\widehat f(n)e^{inx}$ 在 $L^2$ 中收敛到 $f$。
\end{theorem}

\begin{proof}
对有限集 $E\subset\Z$，令 $p_E=\sum_{n\in E}\widehat f(n)e_n$。正交性给出
\begin{align*}
 \|f-p_E\|_2^2
 &=\|f\|_2^2-2\operatorname{Re}\langle f,p_E\rangle+\|p_E\|_2^2\\
 &=\|f\|_2^2-\sum_{n\in E}|\widehat f(n)|^2.
\end{align*}
左边非负；令 $E$ 递增耗尽 $\Z$，得到 Bessel 不等式
$\sum_n|\widehat f(n)|^2\leq\|f\|_2^2$。

取 $E_N=\{-N,\ldots,N\}$。若 $q$ 是频率都落在 $E_N$ 中的三角多项式，则
$f-S_Nf$ 与 $q-S_Nf$ 正交，因而
\[
 \|f-q\|_2^2
 =\|f-S_Nf\|_2^2+\|S_Nf-q\|_2^2
 \geq\|f-S_Nf\|_2^2.
\]
给定 $\varepsilon>0$，推论~\ref{corollary:7-2-1-2} 给出三角多项式 $q$ 使
$\|f-q\|_2<\varepsilon$；当 $N$ 包含 $q$ 的全部频率时，上式说明
$\|f-S_Nf\|_2<\varepsilon$。故 $S_Nf\to f$ 于 $L^2$。在
\[
 \|S_Nf\|_2^2=\sum_{|n|\leq N}|\widehat f(n)|^2
\]
中令 $N\to\infty$，便得 Parseval 范数等式。对 $f+g$、$f+ig$ 使用范数等式，或直接把
$S_Nf,S_Ng$ 的有限和内积取极限，得到式~\eqref{eq:7-2-circle-parseval} 的内积版本。

尚需证明满射。任取 $a=(a_n)\in\ell^2(\Z)$，令
$p_N=\sum_{|n|\leq N}a_ne_n$。若 $M<N$，正交性给出
\[
 \|p_N-p_M\|_2^2=\sum_{M<|n|\leq N}|a_n|^2\longrightarrow0.
\]
所以 $(p_N)$ 在 $L^2(\T)$ 中为 Cauchy。$L^2$ 的完备性给出
$f\in L^2$ 使 $p_N\to f$。固定 $n$ 并取 $N\geq|n|$；Cauchy--Schwarz 不等式给出
\[
 |\widehat f(n)-a_n|
 =|\langle f-p_N,e_n\rangle|
 \leq\|f-p_N\|_2\longrightarrow0.
\]
因此 $\widehat f(n)=a_n$ 对每个 $n$ 成立。这一步同时使用了 $L^2$ 完备性，并证明 Fourier
映射的像不是 $\ell^2$ 的某个真闭子空间。
\end{proof}

\begin{proposition}[Fourier 唯一性]\label{proposition:7-2-1-4}
若 $\mu$ 是 $\T$ 上有限复 Borel 测度，且
$\widehat\mu(n)=0$ 对每个 $n\in\Z$ 成立，则 $\mu=0$。因此两个有限复测度具有相同的全部
三角矩，当且仅当它们相等。
\end{proposition}

\begin{proof}
若 $p(x)=\sum_{|n|\leq N}c_ne^{inx}$，则
\[
 \int_\T p\dd\mu=\sum_{|n|\leq N}c_n\widehat\mu(-n)=0.
\]
给定 $h\in C(\T)$，由推论~\ref{corollary:7-2-1-2} 可取三角多项式 $p_k$ 使
$\|p_k-h\|_\infty\to0$。全变差估计给出
\[
 \left|\int_\T h\dd\mu\right|
 \leq\left|\int_\T(h-p_k)\dd\mu\right|
 \leq\|h-p_k\|_\infty |\mu|(\T)\longrightarrow0.
\]
故 $\mu$ 消去所有连续函数。为从这里恢复测度本身，写极分解
$\dd\mu=u\dd|\mu|$，其中 $|u|=1$ a.e.。连续函数在
$L^1(|\mu|)$ 中稠密，故可取 $h_k\in C(\T)$ 使
$h_k\to\overline u$ 于 $L^1(|\mu|)$。于是
\[
 0=\lim_{k\to\infty}\int h_k\dd\mu
  =\int\overline u\,u\dd|\mu|=|\mu|(\T).
\]
所以 $|\mu|=0$，从而 $\mu=0$。对两个测度应用此结论于其差，即得最后一句。
\end{proof}

图~\ref{fig:7-2-dirichlet-fejer} 把两个核的差异画在同一尺度上。Dirichlet 核的正负旁瓣使
$L^1$ 质量累积；\Fejer{} 核只有非负主峰，峰外质量随 $N$ 消失。恢复机制不是“删除高频”，而是给
第 $n$ 个已观测矩乘上线性权重 $(1-|n|/(N+1))_+$。

\begin{figure}[htbp]
\centering
\begin{tikzpicture}[x=0.9cm,y=0.24cm]
  \begin{scope}
    \clip (-3.35,-7) rectangle (3.35,18.5);
    \draw[->,BookInk!70] (-3.3,0)--(3.3,0) node[right] {$x$};
    \draw[->,BookInk!70] (0,-6.5)--(0,18) node[above] {$D_8(x)$};
    \draw[BookWarning,thick,domain=-3.14:-0.04,samples=180]
      plot (\x,{sin(8.5*deg(\x))/sin(0.5*deg(\x))});
    \draw[BookWarning,thick,domain=0.04:3.14,samples=180]
      plot (\x,{sin(8.5*deg(\x))/sin(0.5*deg(\x))});
    \fill[BookWarning] (0,17) circle[radius=1.4pt];
    \node[book warning node,anchor=north west] at (-3.2,17.2)
      {符号振荡\\$\|D_N\|_1\asymp\log N$};
  \end{scope}
  \begin{scope}[xshift=8.2cm]
    \clip (-3.35,-1.3) rectangle (3.35,10.3);
    \draw[->,BookInk!70] (-3.3,0)--(3.3,0) node[right] {$x$};
    \draw[->,BookInk!70] (0,-1)--(0,10) node[above] {$F_8(x)$};
    \draw[BookAccent,thick,domain=-3.14:-0.04,samples=180]
      plot (\x,{(sin(4.5*deg(\x))/sin(0.5*deg(\x)))^2/9});
    \draw[BookAccent,thick,domain=0.04:3.14,samples=180]
      plot (\x,{(sin(4.5*deg(\x))/sin(0.5*deg(\x)))^2/9});
    \fill[BookAccent] (0,9) circle[radius=1.4pt];
    \node[book strong node,anchor=north west] at (-3.2,9.7)
      {正峰\\$\int F_N\dd m=1$};
  \end{scope}
  \node[book node] (mom) at (1.1,-10.8) {三角矩\\$\widehat f(n)$};
  \node[book warm node,right=1.4cm of mom] (weight)
    {\Fejer{} 权重\\$\left(1-|n|/(N+1)\right)_+$};
  \node[book strong node,right=1.4cm of weight] (rec) {恢复\\$\sigma_Nf\to f$};
  \draw[book arrow] (mom)--(weight);
  \draw[book accent arrow] (weight)--(rec);
\end{tikzpicture}
\caption{Dirichlet 振荡与 \Fejer{} 正平均。下方流程说明正核如何把三角矩变成稳定恢复。}
\label{fig:7-2-dirichlet-fejer}
\end{figure}

\begin{proposition}[周期热流与循环随机游走]\label{proposition:7-2-1-periodic-evolution}
\begin{enumerate}
  \item 对 $f\in L^2(\T)$ 和 $t\geq0$，由
  \[
    \widehat{H_tf}(n)=e^{-n^2t}\widehat f(n)
  \]
  定义的算子族 $(H_t)_{t\geq0}$ 是 $L^2(\T)$ 上的强连续收缩半群。对每个 $t>0$，
  $u(t,\cdot)=H_tf$ 有光滑代表元，并满足
  \[
    \partial_tu=\partial_x^2u,
    \qquad H_tf\longrightarrow f\quad\text{于 }L^2(\T)\text{ 中}\quad(t\downarrow0).
  \]
  \item 设 $M\geq1$、$G=\Z/M\Z$，$p$ 是 $G$ 上的概率分布，并令
  \[
    (Pf)(x)=\sum_{y\in G}p(y)f(x+y),
    \qquad \chi_k(x)=e^{2\pi ikx/M}.
  \]
  若
  \[
    \widehat f(k)=\frac1M\sum_{x\in G}f(x)\overline{\chi_k(x)},
    \qquad
    \lambda_k=\sum_{y\in G}p(y)\chi_k(y),
  \]
  则对每个整数 $m\geq0$，
  \begin{equation}\label{eq:7-2-cyclic-walk-diagonalization}
    P^mf=\sum_{k=0}^{M-1}\lambda_k^m\widehat f(k)\chi_k.
  \end{equation}
\end{enumerate}
\end{proposition}

\begin{proof}
Plancherel 等式给出
\[
 \norm{H_tf}_2^2
 =\sum_{n\in\Z}e^{-2n^2t}|\widehat f(n)|^2
 \leq\norm f_2^2.
\]
乘子恒等式
$e^{-n^2(s+t)}=e^{-n^2s}e^{-n^2t}$ 给出 $H_{s+t}=H_sH_t$。又由支配收敛，
\[
 \norm{H_tf-f}_2^2
 =\sum_{n\in\Z}|e^{-n^2t}-1|^2|\widehat f(n)|^2\longrightarrow0.
\]
半群律与收缩性把零时刻的连续性推广到每个 $t\geq0$。
若 $t>0$ 且 $r\geq0$，Cauchy--Schwarz 不等式给
\[
 \sum_{n\in\Z}|n|^re^{-n^2t}|\widehat f(n)|
 \leq\left(\sum_{n\in\Z}|n|^{2r}e^{-2n^2t}\right)^{1/2}\norm f_2<\infty.
\]
因此 Fourier 级数可任意次逐项对 $x$ 微分；在任意 $t\geq t_0>0$ 上还可逐项对 $t$
微分。乘子的导数为 $-n^2e^{-n^2t}$，故 $\partial_tH_tf=\partial_x^2H_tf$。

在有限群 $G$ 上，对 $0\leq k,\ell<M$，有限几何和给出
\[
 \frac1M\sum_{x\in G}\chi_k(x)\overline{\chi_\ell(x)}
 =\begin{cases}
 1,&k=\ell,\\
 \displaystyle\frac{1-e^{2\pi i(k-\ell)}}
 {M\bigl(1-e^{2\pi i(k-\ell)/M}\bigr)}=0,&k\ne\ell.
 \end{cases}
\]
因此这些特征标构成 $M$ 维空间 $\C^G$ 在归一化计数测度下的正交基，从而
$f=\sum_k\widehat f(k)\chi_k$。直接计算得
\[
 P\chi_k(x)=\sum_{y\in G}p(y)\chi_k(x+y)
 =\lambda_k\chi_k(x).
\]
对有限 Fourier 展开逐项应用 $P^m$，即得
\eqref{eq:7-2-cyclic-walk-diagonalization}。例如，当 $M\geq3$ 且
$p(1)=p(-1)=1/2$ 时，
\[
 \lambda_k=\cos\frac{2\pi k}{M}.
\]
若 $M$ 为奇数，则 $k\neq0$ 时 $|\lambda_k|<1$，故 $P^mf$ 趋于 $f$ 的平均值；若 $M$ 为偶数，则
$\lambda_{M/2}=-1$，对应模式显示出二周期振荡。
\end{proof}

\begin{exercise}
\begin{enumerate}
  \item 从 $S_Nf=f*D_N$ 与 $\sigma_Nf=(N+1)^{-1}\sum_{j=0}^NS_jf$ 出发，逐步推出
        \[
        D_N(x)=\frac{\sin((N+\tfrac12)x)}{\sin(x/2)},\qquad
        F_N(x)=\frac1{N+1}\left(\frac{\sin((N+1)x/2)}{\sin(x/2)}\right)^2,
        \]
        并在 $x=0$ 处用连续延拓解释公式。
  \item 对每个 $\delta\in(0,\pi)$ 给出
        $\sup_{\delta\le |x|\le\pi}F_N(x)$ 的显式上界，由此证明
        $\int_{|x|\ge\delta}F_N\dd m\to0$；同时直接核对 $F_N\ge0$ 与
        $\int_\T F_N\dd m=1$。
  \item 给定 $f\in C(\T)$ 与 $\varepsilon>0$，把
        $\sigma_Nf(x)-f(x)$ 的积分分成 $|t|<\delta$ 与 $|t|\ge\delta$ 两段，完整证明
        $\norm{\sigma_Nf-f}_\infty\to0$。再指出把一致连续性替换成 $L^p$ 平移连续性后，哪一步给出
        $1\le p<\infty$ 的 $L^p$ 收敛。
  \item 设有限复 Borel 测度 $\mu$ 的全部 Fourier 系数为零。先证明它消去每个三角多项式，再用
        \Fejer{} 稠密性与全变差估计证明它消去 $C(\T)$；最后写出极分解逼近步骤，推出
        $|\mu|(\T)=0$。
\end{enumerate}
\end{exercise}

圆周上的频率是整数，Fourier 映射把 $L^2(\T)$ 完整地离散化为 $\ell^2(\Z)$。在
$\R^d$ 上，频率变成连续变量；相应的求和、矩阵坐标和离散 Parseval 等式将分别变成积分、乘子和
Plancherel 等式。


\subsection{\texorpdfstring{$\R^d$}{R d} Fourier 变换、卷积与反演}\label{subsec:7-2-2}
\index{Fourier 变换}\index{Riemann--Lebesgue 定理}\index{Fourier 反演}

平移不变算子的共同特征是：指数函数把平移变成一个相位因子。若
$\tau_af(x)=f(x-a)$，则形式上
\[
 \int_{\R^d}f(x-a)e^{-2\pi ix\cdot\xi}\dd x
 =e^{-2\pi ia\cdot\xi}\int_{\R^d}f(x)e^{-2\pi ix\cdot\xi}\dd x.
\]
因此，卷积在频率侧应成为乘法，微分应成为多项式乘子。以下三个计算先展示这套规则的力量及其
边界；本小节始终使用指数 $e^{-2\pi ix\cdot\xi}$。

\begin{example}[Gaussian]\label{ex:7-2-2-1}
先在一维令 $g(x)=e^{-\pi x^2}$，并记
$G(\xi)=\int_\R g(x)e^{-2\pi ix\xi}\dd x$。对 $|h|\leq1$，差商恒等式给
\[
 \left|g(x)e^{-2\pi ix\xi}
       \frac{e^{-2\pi ixh}-1}{h}\right|
 \leq2\pi|x|e^{-\pi x^2},
\]
而右端可积，故支配收敛允许对 $\xi$ 求导。又因
$g'(x)=-2\pi xg(x)$，分部积分的边界项为零，故
\begin{align*}
 G'(\xi)
 &=-2\pi i\int_\R xg(x)e^{-2\pi ix\xi}\dd x\\
 &=i\int_\R g'(x)e^{-2\pi ix\xi}\dd x
 =-2\pi\xi G(\xi).
\end{align*}
Gaussian 积分给出 $G(0)=\int e^{-\pi x^2}\dd x=1$。把微分方程乘以
$e^{\pi\xi^2}$，可见
\[
 \frac{\dd}{\dd\xi}\bigl(e^{\pi\xi^2}G(\xi)\bigr)=0;
\]
代入初值得到
\[
 G(\xi)=e^{-\pi\xi^2}.
\]
对 $g_a(x)=e^{-\pi a|x|^2}$，$a>0$，逐坐标使用 Fubini，再作伸缩
$y=\sqrt a\,x$，得到
\begin{equation}\label{eq:7-2-gaussian-transform}
 \widehat{g_a}(\xi)=a^{-d/2}e^{-\pi|\xi|^2/a}.
\end{equation}
Gaussian 因而在空间侧与频率侧具有相同形状，宽度则互为倒数。
\end{example}

\begin{example}[区间指标]\label{ex:7-2-2-2}
在一维令 $I=[-1/2,1/2]$。当 $\xi\neq0$ 时
\[
 \widehat{\1_I}(\xi)
 =\int_{-1/2}^{1/2}e^{-2\pi ix\xi}\dd x
 =\frac{e^{-\pi i\xi}-e^{\pi i\xi}}{-2\pi i\xi}
 =\frac{\sin(\pi\xi)}{\pi\xi};
\]
在 $\xi=0$ 处连续延拓值为 $1$。该 sinc 函数不属于 $L^1(\R)$。事实上，对每个整数
$k\geq1$，在 $[k+1/6,k+5/6]$ 上有 $|\sin(\pi\xi)|\geq1/2$，从而
\[
 \int_1^\infty\left|\frac{\sin(\pi\xi)}{\pi\xi}\right|\dd\xi
 \geq\frac1{2\pi}\sum_{k=1}^\infty
   \frac{2/3}{k+5/6}=\infty.
\]
所以 $f\in L^1$ 并不保证 $\widehat f\in L^1$。区间边界的跳跃只产生 $|\xi|^{-1}$ 级衰减，
未经正则化的反演积分不绝对收敛。
\end{example}

\begin{example}[卷积平滑]\label{ex:7-2-2-3}
仍令 $I=[-1/2,1/2]$。对每个 $x$，两个区间 $I$ 与 $x-I$ 的交长给出
\[
 (\1_I*\1_I)(x)
 =\int_\R\1_I(x-y)\1_I(y)\dd y
 =(1-|x|)_+.
\]
这次卷积把两个跳跃函数变成连续的帐篷函数。另一方面，稍后证明的卷积公式给出
\[
 \widehat{(1-|\cdot|)_+}(\xi)
 =\left(\frac{\sin(\pi\xi)}{\pi\xi}\right)^2.
\]
空间侧多一次卷积，频率侧就多一个 sinc 衰减因子；平方后的函数已经属于 $L^1(\R)$。
\end{example}

\begin{definition}[Fourier 变换约定]\label{def:7-2-2-1}
对 $f\in L^1(\R^d)$，定义
\[
 \widehat f(\xi)=\mathcal Ff(\xi)
 =\int_{\R^d}f(x)e^{-2\pi ix\cdot\xi}\dd x.
\]
对 $g\in L^1(\R^d)$，定义反 Fourier 积分
\[
 \check g(x)=\mathcal F^{-1}g(x)
 =\int_{\R^d}g(\xi)e^{2\pi ix\cdot\xi}\dd\xi.
\]
对有限复 Borel 测度 $\mu$，相同约定给出
$\widehat\mu(\xi)=\int e^{-2\pi ix\cdot\xi}\dd\mu(x)$。
\end{definition}

文献中常把 $2\pi$ 放在逆变换的常数里。三种常用写法的换算如下；在本书余下部分只使用第一行。
\begin{center}
\renewcommand{\arraystretch}{1.25}
\small
\begin{tabularx}{\textwidth}{@{}p{2.0cm}XX@{}}
\toprule
约定 & 正变换 & 逆变换\\
\midrule
本书 & $\widehat f(\xi)=\int f(x)e^{-2\pi ix\cdot\xi}\dd x$
     & $f(x)=\int\widehat f(\xi)e^{2\pi ix\cdot\xi}\dd\xi$\\
角频率 & $F_0(\omega)=\int f(x)e^{-ix\cdot\omega}\dd x$
     & $f(x)=(2\pi)^{-d}\int F_0(\omega)e^{ix\cdot\omega}\dd\omega$\\
酉角频率 & $F_u(\omega)=(2\pi)^{-d/2}\int f(x)e^{-ix\cdot\omega}\dd x$
     & $f(x)=(2\pi)^{-d/2}\int F_u(\omega)e^{ix\cdot\omega}\dd\omega$\\
\bottomrule
\end{tabularx}
\end{center}
其中 $F_0(\omega)=\widehat f(\omega/(2\pi))$，且
$F_u=(2\pi)^{-d/2}F_0$。常数的所有差异都来自换元 $\omega=2\pi\xi$。

\begin{definition}[调制、平移与伸缩]\label{def:7-2-2-2}
对 $a,b\in\R^d$ 及可逆线性映射 $A:\R^d\to\R^d$，定义
\[
 \tau_af(x)=f(x-a),\qquad
 M_bf(x)=e^{2\pi ib\cdot x}f(x),\qquad
 D_Af(x)=f(Ax).
\]
\end{definition}

\begin{definition}[可反演 $L^1$ 类]\label{def:7-2-2-3}
称 $f$ 属于可反演 $L^1$ 类，如果
\[
 f\in L^1(\R^d),\qquad \widehat f\in L^1(\R^d).
\]
\end{definition}

\begin{theorem}[Riemann--Lebesgue]\label{theorem:7-2-2-1}
若 $f\in L^1(\R^d)$，则 $\widehat f\in C_0(\R^d)$，且
\[
 \|\widehat f\|_\infty\leq\|f\|_1.
\]
\end{theorem}

\begin{proof}
三角不等式立即给出
$|\widehat f(\xi)|\leq\int|f|=\|f\|_1$。若 $\xi_k\to\xi$，则积分中的指数逐点收敛，且
\[
 |f(x)e^{-2\pi ix\cdot\xi_k}|\leq|f(x)|.
\]
支配收敛定理给出 $\widehat f(\xi_k)\to\widehat f(\xi)$，故 $\widehat f$ 连续。

先设 $g\in C_c(\R^d)$。给 $\xi\neq0$，选坐标 $j$ 使
$|\xi_j|\geq|\xi|/\sqrt d$，并令 $h=e_j/(2\xi_j)$。于是
$e^{-2\pi ih\cdot\xi}=-1$。在线性换元后，
\begin{align*}
 2\widehat g(\xi)
 &=\int_{\R^d}\bigl(g(x)-g(x-h)\bigr)e^{-2\pi ix\cdot\xi}\dd x,\\
 2|\widehat g(\xi)|&\leq\|g-\tau_hg\|_1.
\end{align*}
当 $|\xi|\to\infty$ 时，$|h|\leq\sqrt d/(2|\xi|)\to0$；$L^1$ 平移连续性使右端趋于零。

对一般 $f\in L^1$，给定 $\varepsilon>0$，取 $g\in C_c$ 使
$\|f-g\|_1<\varepsilon$。由第一步的上界，
\[
 |\widehat f(\xi)|
 \leq\|f-g\|_1+|\widehat g(\xi)|<2\varepsilon
\]
对所有充分大的 $|\xi|$ 成立。因此 $\widehat f$ 在无穷远消失。
\end{proof}

\begin{proposition}[代数公式]\label{proposition:7-2-2-2}
在下列各项所写的可积条件下，Fourier 变换满足：
\begin{enumerate}
  \item 若 $f,g\in L^1$，则
  $\widehat{f*g}=\widehat f\,\widehat g$；
  \item 若 $f\in W^{1,1}(\R^d)$，则
  $\widehat{\partial_jf}(\xi)=2\pi i\xi_j\widehat f(\xi)$；
  \item 若 $f,x_jf\in L^1$，则 $\widehat f$ 关于 $\xi_j$ 连续可微，且
  \[
    \widehat{x_jf}(\xi)=-\frac1{2\pi i}\partial_{\xi_j}\widehat f(\xi);
  \]
  \item 若 $f\in L^1(\R^d)$，则 $\tau_af,M_bf,D_Af\in L^1(\R^d)$，并且
  \[
    \widehat{\tau_af}(\xi)=e^{-2\pi ia\cdot\xi}\widehat f(\xi),
    \qquad
    \widehat{M_bf}(\xi)=\widehat f(\xi-b),
  \]
  \[
    \widehat{D_Af}(\xi)=|\det A|^{-1}\widehat f(A^{-T}\xi).
  \]
\end{enumerate}
\end{proposition}

\begin{proof}
若 $f,g\in L^1$，则
\[
 \int\!\!\int |f(x-y)g(y)|\dd y\dd x=\|f\|_1\|g\|_1<\infty.
\]
所以 Fubini 定理与换元 $z=x-y$ 给出
\begin{align*}
 \widehat{f*g}(\xi)
 &=\int\!\!\int f(x-y)g(y)e^{-2\pi ix\cdot\xi}\dd y\dd x\\
 &=\int\!\!\int f(z)g(y)e^{-2\pi i(z+y)\cdot\xi}\dd z\dd y
 =\widehat f(\xi)\widehat g(\xi).
\end{align*}

若 $f\in W^{1,1}$，取 $\chi\in C_c^\infty$ 使
$0\leq\chi\leq1$、$\chi=1$ 于单位球，并令 $\chi_R(x)=\chi(x/R)$。弱导数的积分分部公式给出
\begin{align*}
 \int \partial_jf(x)e^{-2\pi ix\cdot\xi}\chi_R(x)\dd x
 &=2\pi i\xi_j\int f(x)e^{-2\pi ix\cdot\xi}\chi_R(x)\dd x\\
 &\quad-\int f(x)e^{-2\pi ix\cdot\xi}\partial_j\chi_R(x)\dd x.
\end{align*}
前两项由支配收敛分别趋于 $\widehat{\partial_jf}(\xi)$ 与
$2\pi i\xi_j\widehat f(\xi)$；最后一项绝对值不超过
\[
 R^{-1}\|\partial_j\chi\|_\infty\|f\|_1,
\]
故趋于零。这证明微分公式。

若 $x_jf\in L^1$，差商恒等式以及
$|(e^{-it}-1)/t|\leq1$ 给出一个由 $2\pi|x_jf(x)|$ 控制的支配函数。因此
\[
 \partial_{\xi_j}\widehat f(\xi)
 =-2\pi i\int x_jf(x)e^{-2\pi ix\cdot\xi}\dd x
 =-2\pi i\widehat{x_jf}(\xi).
\]
同一支配收敛论证说明导数连续。

最后，平移、调制与线性伸缩的 $L^1$ 归属分别由平移不变性、指数因子的模为 $1$ 以及线性换元
得到。前两式只需分别作 $y=x-a$ 和合并指数。在线性伸缩中令 $y=Ax$，便有
\begin{align*}
 \widehat{D_Af}(\xi)
 &=|\det A|^{-1}\int f(y)e^{-2\pi i(A^{-1}y)\cdot\xi}\dd y\\
 &=|\det A|^{-1}\widehat f(A^{-T}\xi).
\end{align*}
这里 $|\det A|^{-1}$ 来自换元公式中的雅可比行列式因子，而 $A^{-T}$ 由
$(A^{-1}y)\cdot\xi=y\cdot A^{-T}\xi$ 决定。
\end{proof}

\begin{theorem}[Fourier 反演]\label{theorem:7-2-2-3}
若 $f,\widehat f\in L^1(\R^d)$，则在所选代表 $f$ 的每个 Lebesgue 点 $x$，
\begin{equation}\label{eq:7-2-fourier-inversion}
 f(x)=\int_{\R^d}\widehat f(\xi)e^{2\pi ix\cdot\xi}\dd\xi.
\end{equation}
特别地，等式在每个连续点成立，而且 $f$ 由 $\widehat f$ 唯一决定。
\end{theorem}

\begin{proof}
对 $\varepsilon>0$，令
\[
 \phi_\varepsilon(x)=\varepsilon^{-d/2}e^{-\pi|x|^2/\varepsilon}.
\]
由例~\ref{ex:7-2-2-1}，
$\widehat{\phi_\varepsilon}(\xi)=e^{-\pi\varepsilon|\xi|^2}$，且
$\int\phi_\varepsilon=1$。由于
\[
 \int\!\!\int |f(y)|e^{-\pi\varepsilon|\xi|^2}\dd\xi\dd y
 =\varepsilon^{-d/2}\|f\|_1<\infty,
\]
Fubini 定理允许交换积分，得到 Gaussian 正则化恒等式
\begin{align}
 I_\varepsilon(x)
 &:=\int\widehat f(\xi)e^{2\pi ix\cdot\xi}
                 e^{-\pi\varepsilon|\xi|^2}\dd\xi\notag\\
 &=\int f(y)\left(\int e^{2\pi i(x-y)\cdot\xi}
                 e^{-\pi\varepsilon|\xi|^2}\dd\xi\right)\dd y\notag\\
 &=(f*\phi_\varepsilon)(x).
 \label{eq:7-2-gaussian-regularization}
\end{align}
写 $\delta=\sqrt\varepsilon$，便有
$\phi_\varepsilon(x)=\delta^{-d}e^{-\pi|x/\delta|^2}$。所以它正是质量一的径向递减伸缩核；
第四章的 Lebesgue 点近似恒等定理
\ref{theorem:4-4-2-2} 因而给出
$(f*\phi_\varepsilon)(x)\to f(x)$，只要 $x$ 是 $f$ 的 Lebesgue 点。
另一方面，$\widehat f\in L^1$，且 Gaussian 乘子绝对值不超过 $1$ 并逐点趋于 $1$；支配收敛给出
\[
 I_\varepsilon(x)\longrightarrow
 \int\widehat f(\xi)e^{2\pi ix\cdot\xi}\dd\xi.
\]
在式~\eqref{eq:7-2-gaussian-regularization} 两边取极限，即得反演公式。

右端是连续函数，因为对 $x_k\to x$ 可用 $|\widehat f|$ 作支配。Lebesgue 微分定理说明几乎每点是
Lebesgue 点，故这个连续函数与 $f$ a.e. 相等。两个连续函数若 a.e. 相等，则其差的非零点集既开又为
零测，只能为空；所以连续代表唯一。若两个可反演函数的 Fourier 变换相同，反演公式给出其连续代表
相同，因而两个 $L^1$ 等价类相同。
\end{proof}

\begin{proposition}[测度唯一性]\label{proposition:7-2-2-4}
若 $\mu,\nu$ 是 $\R^d$ 上有限复 Borel 测度，且
$\widehat\mu(\xi)=\widehat\nu(\xi)$ 对每个 $\xi\in\R^d$ 成立，则 $\mu=\nu$。
\end{proposition}

\begin{proof}
令 $\lambda=\mu-\nu$，则 $\widehat\lambda=0$。用反演证明中的 Gaussian 定义
\[
 h_\varepsilon(x)=\int_{\R^d}\phi_\varepsilon(x-y)\dd\lambda(y).
\]
Tonelli 定理给出
$\|h_\varepsilon\|_1\leq|\lambda|(\R^d)$。Gaussian 的一致连续性与全变差估计说明
$h_\varepsilon$ 连续。再用 Fubini 及命题~\ref{proposition:7-2-2-2} 的同一计算，
\[
 \widehat{h_\varepsilon}(\xi)
 =e^{-\pi\varepsilon|\xi|^2}\widehat\lambda(\xi)=0.
\]
于是 $h_\varepsilon,\widehat{h_\varepsilon}\in L^1$；定理~\ref{theorem:7-2-2-3} 作用于连续函数
$h_\varepsilon$，给出 $h_\varepsilon(x)=0$ 对每个 $x$ 成立。

任取 $\psi\in C_c(\R^d)$。Fubini 定理和 $\phi_\varepsilon(-z)=\phi_\varepsilon(z)$ 给出
\[
 0=\int\psi(x)h_\varepsilon(x)\dd x
 =\int(\psi*\phi_\varepsilon)(y)\dd\lambda(y).
\]
Gaussian 是近似恒等族，而 $\psi\in C_0$，故
$\psi*\phi_\varepsilon\to\psi$ 一致。由
\[
 \left|\int(\psi*\phi_\varepsilon-\psi)\dd\lambda\right|
 \leq\|\psi*\phi_\varepsilon-\psi\|_\infty|\lambda|(\R^d)
\]
可知 $\int\psi\dd\lambda=0$。最后写
$\dd\lambda=u\dd|\lambda|$，并用 $C_c(\R^d)$ 在
$L^1(|\lambda|)$ 中的稠密性逼近 $\overline u$，得到
$|\lambda|(\R^d)=0$。所以 $\lambda=0$，即 $\mu=\nu$。
\end{proof}

图~\ref{fig:7-2-convolution-inversion} 同时记录两条路径。上方交换方块把空间卷积送到频率乘法；
下方则说明反演并非把振荡积分直接“算回去”，而是先乘 Gaussian、回到空间侧的近似恒等，再令
$\varepsilon\downarrow0$。积分交换在固定 $\varepsilon>0$ 时绝对可积，这正是正则化的作用。

\begin{figure}[htbp]
\centering
\begin{tikzpicture}[node distance=1.25cm and 2.0cm]
  \node[book node] (fg) {$f,\ g\in L^1$};
  \node[book warm node,right=of fg] (conv) {空间卷积\\$f*g$};
  \node[book node,below=of fg] (hats) {$\widehat f,\ \widehat g$};
  \node[book strong node,right=of hats] (prod) {频率乘法\\$\widehat f\,\widehat g$};
  \draw[book arrow] (fg)--node[above,book label]{卷积} (conv);
  \draw[book arrow] (fg)--node[left,book label]{Fourier} (hats);
  \draw[book arrow] (conv)--node[right,book label]{Fourier} (prod);
  \draw[book accent arrow] (hats)--node[below,book label]{逐点乘} (prod);

  \node[book node,below=1.8cm of hats] (obs) {观测\\$\widehat f$};
  \node[book warm node,right=1.8cm of obs] (reg)
    {Gaussian 正则化\\$e^{-\pi\varepsilon|\xi|^2}\widehat f$};
  \node[book strong node,below=1.5cm of reg] (ai)
    {空间近似恒等\\$f*\phi_\varepsilon$};
  \node[book strong node,right=1.8cm of ai] (limit) {Lebesgue 点\\$f(x)$};
  \draw[book arrow] (obs)--(reg);
  \draw[book accent arrow] (reg)--node[right,book label,align=center]{绝对可积\\反变换} (ai);
  \draw[book accent arrow] (ai)--node[above,book label]{$\varepsilon\downarrow0$} (limit);
\end{tikzpicture}
\caption{卷积--乘法交换方块与 Gaussian 正则化反演路径。}
\label{fig:7-2-convolution-inversion}
\end{figure}

这项测度唯一性是第五章“特征函数决定概率分布”的有限复测度版本。它只回答相同变换能否来自两个
不同对象；若一列概率变换逐点收敛，要进一步推出概率弱收敛，仍需紧性以及极限在零点连续，
不能从唯一性单独推出存在性。

\begin{exercise}
\begin{enumerate}
  \item 从一维函数 $g(x)=e^{-\pi x^2}$ 的微分方程证明 $\widehat g=g$，再用 Fubini 与伸缩换元
        推出
        $\widehat{e^{-\pi a|x|^2}}(\xi)=a^{-d/2}e^{-\pi|\xi|^2/a}$；逐项注明支配函数与
        雅可比行列式因子。
  \item 在本章的 $2\pi$ 约定下，对 $f\in L^1(\R^d)$ 从定义证明平移、调制与可逆线性伸缩公式：
        \[
        \begin{aligned}
        \widehat{f(\,\cdot-a)}(\xi)
          &=e^{-2\pi ia\cdot\xi}\widehat f(\xi),\\
        \widehat{e^{2\pi i\eta\cdot x}f(x)}(\xi)
          &=\widehat f(\xi-\eta),\\
        \widehat{f(A\,\cdot)}(\xi)
          &=|\det A|^{-1}\widehat f(A^{-T}\xi).
        \end{aligned}
        \]
  \item 对 $f,g\in L^1(\R^d)$ 从头验证定义 $f*g$ 所需的 a.e. 可积性，并写出允许交换积分的
        绝对可积上界，由此证明 $\widehat{f*g}=\widehat f\,\widehat g$。
  \item 令 $\phi_\varepsilon(x)=\varepsilon^{-d/2}e^{-\pi|x|^2/\varepsilon}$。在
        $f,\widehat f\in L^1$ 的假设下，完整证明
        \[
        \int\widehat f(\xi)e^{2\pi ix\cdot\xi}e^{-\pi\varepsilon|\xi|^2}\dd\xi
        =(f*\phi_\varepsilon)(x),
        \]
        并分别说明 Lebesgue 点近似恒等定理与支配收敛在令 $\varepsilon\downarrow0$ 时负责哪一边。
\end{enumerate}
\end{exercise}


\subsection{Schwartz 空间、Plancherel 与热/Poisson 半群}\label{subsec:7-2-3}
\index{Schwartz 空间}\index{Fourier 乘子}\index{热半群}\index{Poisson 半群}

$L^1$ 反演的附加条件 $\widehat f\in L^1$ 并不自动成立。为了得到一个对微分、乘多项式和
Fourier 变换均封闭的测试函数空间，需要同时控制无穷远处的衰减与所有阶导数。Schwartz 空间用一族
半范数同时表达这些要求，而不将它们归结为单个范数。随后，$L^2$ 变换将由这个稠密测试函数空间完备化得到；
这与圆周上从三角多项式完备化到 $L^2(\T)$ 是同一结构。

若 $\alpha=(\alpha_1,\ldots,\alpha_d)\in(\Z_{\geq0})^d$，沿用
$|\alpha|=\sum_j\alpha_j$、$x^\alpha=\prod_jx_j^{\alpha_j}$ 与
$\partial^\alpha=\prod_j\partial_j^{\alpha_j}$ 的多重指标记号。

\begin{definition}[Schwartz 空间]\label{def:7-2-3-1}
Schwartz 空间 $\mathcal S(\R^d)$ 由所有 $f\in C^\infty(\R^d)$ 组成，并要求对每对多重指标
$\alpha,\beta$，
\[
 p_{\alpha,\beta}(f)
 :=\sup_{x\in\R^d}|x^\alpha\partial^\beta f(x)|<\infty.
\]
取遍所有 $p_{\alpha,\beta}$ 所生成的局部凸拓扑；亦即
$f_n\to f$ 当且仅当每个 $p_{\alpha,\beta}(f_n-f)\to0$。这族可数半范数给出
$\mathcal S$ 的 \Frechet{} 拓扑。
\end{definition}

为说明最后一句，枚举半范数为 $(p_k)_{k\geq1}$，则
\[
 d_{\mathcal S}(f,g)=\sum_{k=1}^\infty2^{-k}
   \min\{1,p_k(f-g)\}
\]
诱导上述拓扑。若 $(f_n)$ 关于该度量 Cauchy，则固定一个半范数 $p_k$ 和
$0<\varepsilon<1$ 后，当
$d_{\mathcal S}(f_n,f_m)<2^{-k}\varepsilon$ 时必有
$p_k(f_n-f_m)<\varepsilon$。因此每个 $(\partial^\beta f_n)$ 关于全空间上确界范数 Cauchy，存在
有界连续函数 $g_\beta$ 使
\[
 \|\partial^\beta f_n-g_\beta\|_\infty\longrightarrow0.
\]
置 $f=g_0$。对每个多重指标 $\gamma$、坐标 $j$ 与实数 $t$，逐坐标使用
\[
 \partial^\gamma f_n(x+te_j)-\partial^\gamma f_n(x)
 =\int_0^t\partial^{\gamma+e_j}f_n(x+se_j)\dd s.
\]
被积函数在连接 $x$ 与 $x+te_j$ 的线段上一致收敛；令 $n\to\infty$ 得
\[
 g_\gamma(x+te_j)-g_\gamma(x)
 =\int_0^tg_{\gamma+e_j}(x+se_j)\dd s.
\]
微积分基本定理于是给 $\partial_jg_\gamma=g_{\gamma+e_j}$。对 $|\gamma|$ 归纳可知
$f\in C^\infty$ 且 $\partial^\beta f=g_\beta$。

最后，$(x^\alpha\partial^\beta f_n)$ 在上确界范数中也是 Cauchy，故一致收敛到某个
$h_{\alpha,\beta}$；逐点取极限给
$h_{\alpha,\beta}=x^\alpha\partial^\beta f$。因此
\[
 p_{\alpha,\beta}(f_n-f)
 =\|x^\alpha\partial^\beta f_n-h_{\alpha,\beta}\|_\infty
 \longrightarrow0.
\]
所以极限仍在 $\mathcal S$ 中，且该度量完备。

Schwartz 函数属于每个 $L^p$。例如
\begin{equation}\label{eq:7-2-schwartz-l1-bound}
 \|h\|_1
 \leq C_d\sup_x(1+|x|)^{d+1}|h(x)|
 \leq C'_d\sum_{|\gamma|\leq d+1}p_{\gamma,0}(h).
\end{equation}
微分或乘以多项式只会移动半范数的指标；Leibniz 公式还说明两个 Schwartz 函数的乘积仍在
$\mathcal S$。卷积也保持该空间。事实上，对 $f,g\in\mathcal S$，微分可移入卷积，而
$x=(x-y)+y$ 给出
\[
 x^\alpha\partial^\beta(f*g)(x)
 =\sum_{\gamma\leq\alpha}\binom{\alpha}{\gamma}
 \bigl[(z^{\alpha-\gamma}\partial^\beta f(z))
       *(z^\gamma g(z))\bigr](x).
\]
每一项都是一个 $L^\infty$ 函数与一个 $L^1$ 函数的卷积，故一致有界；这证明
$p_{\alpha,\beta}(f*g)<\infty$。

\begin{definition}[缓增分布]\label{def:7-2-3-2}
Schwartz 空间上的连续线性泛函称为\emph{缓增分布}（tempered distribution），全体记作
$\mathcal S'(\R^d)$。
\end{definition}

\begin{proposition}[缓增分布的基本例子]\label{proposition:7-2-3-1}
下列对象都以自然方式给出 $\mathcal S'(\R^d)$ 中的元素。
\begin{enumerate}
  \item 若 $f\in L^1_{\mathrm{loc}}(\R^d)$，并且对某些 $C>0$ 与整数 $N\geq0$ 有
  $|f(x)|\leq C(1+|x|)^N$ a.e.，则
  \[
    \langle T_f,\varphi\rangle=\int_{\R^d}f(x)\varphi(x)\dd x.
  \]
  \item 每个 $f\in L^p(\R^d)$（$1\leq p\leq\infty$）按同一公式定义缓增分布。
  \item 每个多项式按第一项的公式定义缓增分布。
  \item 若复 Radon 测度 $\mu$ 对某些 $C,M>0$ 满足
  \[
    |\mu|(B(0,R))\leq C(1+R)^M\qquad(R\geq1),
  \]
  则 $\langle T_\mu,\varphi\rangle=\int\varphi\dd\mu$ 定义缓增分布。
\end{enumerate}
\end{proposition}

\begin{proof}
在第一种情形，取整数 $K>N+d$。则
\[
 |\langle T_f,\varphi\rangle|
 \leq C\sup_x(1+|x|)^K|\varphi(x)|
       \int_{\R^d}(1+|x|)^{N-K}\dd x.
\]
加权上确界由有限个 $p_{\alpha,0}$ 控制，所以 $T_f$ 连续；多项式是这一情形的特例。

若 $f\in L^p$，\Holder{} 不等式给出
\[
 |\langle T_f,\varphi\rangle|
 \leq \|f\|_p\|\varphi\|_{p'},
\]
而当 $1\leq q<\infty$ 时，
\[
 \|\varphi\|_q
 \leq\left(\int_{\R^d}(1+|x|)^{-q(d+1)}\dd x\right)^{1/q}
       \sup_x(1+|x|)^{d+1}|\varphi(x)|;
\]
$q=\infty$ 时则有 $\|\varphi\|_\infty=p_{0,0}(\varphi)$。因此
$\mathcal S$ 连续嵌入每个 $L^q$，所以上述泛函连续。

最后考虑测度。把空间分成 $B(0,2)$ 与环带
$2^k\leq|x|<2^{k+1}$。取整数 $K>M+1$，则
\begin{align*}
 \int|\varphi|\dd|\mu|
 &\leq C_{M,K}\sup_x(1+|x|)^K|\varphi(x)|
   \left(1+\sum_{k=1}^\infty2^{-k(K-M)}\right).
\end{align*}
右端仍由有限个 Schwartz 半范数控制，因而 $T_\mu$ 连续。
\end{proof}

\begin{definition}[Fourier 乘子]\label{def:7-2-3-3}
给定频率函数 $m(\xi)$，在使右端有意义的函数上定义形式算子
\[
 T_mf=\mathcal F^{-1}(m\widehat f).
\]
称 $m$ 为该算子的 Fourier 乘子。
\end{definition}

本节的两个基本乘子族是
\[
 m_t^{\mathrm{heat}}(\xi)=e^{-4\pi^2t|\xi|^2},
 \qquad
 m_t^{\mathrm{Pois}}(\xi)=e^{-2\pi t|\xi|},
 \qquad t>0.
\]
指数恒等式 $m_s m_t=m_{s+t}$ 预示相应算子族具有半群律；要把形式计算变成定理，还须识别其空间核
并证明在所用函数空间中连续。

\begin{theorem}[Schwartz 自同构]\label{theorem:7-2-3-1}
Fourier 变换是 $\mathcal S(\R^d)$ 上的连续线性双射，逆映射为反 Fourier 变换，而且逆映射也连续。
换言之，$\mathcal F:\mathcal S\to\mathcal S$ 是 \Frechet{} 空间自同构。
\end{theorem}

\begin{proof}
对 $f\in\mathcal S$，命题~\ref{proposition:7-2-2-2} 的微分与乘法公式可以任意次迭代，得到
\[
 \partial_\xi^\beta\widehat f
 =(-2\pi i)^{|\beta|}\widehat{x^\beta f},
\]
以及
\begin{equation}\label{eq:7-2-schwartz-seminorm-identity}
 \xi^\alpha\partial_\xi^\beta\widehat f(\xi)
 =(2\pi i)^{-|\alpha|}(-2\pi i)^{|\beta|}
   \widehat{\partial^\alpha(x^\beta f)}(\xi).
\end{equation}
取上确界，使用 $\|\widehat h\|_\infty\leq\|h\|_1$ 与
式~\eqref{eq:7-2-schwartz-l1-bound}，再展开 Leibniz 公式，得到
\begin{equation}\label{eq:7-2-schwartz-continuity}
 p_{\alpha,\beta}(\widehat f)
 \leq C_{\alpha,\beta,d}
 \sum_{(\gamma,\delta)\in E_{\alpha,\beta}}
 p_{\gamma,\delta}(f),
\end{equation}
其中 $E_{\alpha,\beta}$ 是一个有限指标集。因此 $\widehat f\in\mathcal S$，而且每个输出半范数由
有限个输入半范数控制；这正是 $\mathcal F$ 的连续性。

令 $Rg(x)=g(-x)$。由定义 $\mathcal Gg:=\check g=R\widehat g$，而
$p_{\alpha,\beta}(Rg)=p_{\alpha,\beta}(g)$；所以刚证的估计说明反 Fourier 算子
$\mathcal G=R\mathcal F$ 也连续地把 $\mathcal S$ 映入自身。又因
$f,\widehat f\in L^1$ 且 $f$ 处处连续，反演定理~\ref{theorem:7-2-2-3} 给出
\[
 \mathcal G\mathcal Ff=f\qquad(f\in\mathcal S).
\]
由 $\mathcal G=R\mathcal F$，故
$\mathcal G\mathcal F=I$ 等价于 $R\mathcal F^2=I$，即
$\mathcal F^2=R$。线性换元又给出 $\mathcal F R=R\mathcal F$。因此
\[
 \mathcal F\mathcal G
 =\mathcal F R\mathcal F=R\mathcal F^2=R^2=I.
\]
所以 $\mathcal F$ 与 $\mathcal G$ 互为连续逆映射。
\end{proof}

\begin{definition}[缓增分布的 Fourier 变换]\label{def:7-2-3-4}
对 $T\in\mathcal S'(\R^d)$，定义其 Fourier 变换 $\widehat T$ 为
\[
 \langle\widehat T,\varphi\rangle
 :=\langle T,\widehat\varphi\rangle
 \qquad(\varphi\in\mathcal S(\R^d)).
\]
\end{definition}

Schwartz 自同构定理保证 $\varphi\mapsto\widehat\varphi$ 连续，因此上式确实定义
$\mathcal S'$ 中的元素。本章只用这一转置定义计算 Dirac 测度和常数函数的 Fourier 变换；一般的
分布运算留待后续课程。

\begin{theorem}[Plancherel 定理]\label{theorem:7-2-3-2}
对 $f,g\in\mathcal S(\R^d)$，
\begin{equation}\label{eq:7-2-euclidean-parseval}
 \langle\widehat f,\widehat g\rangle_{L^2}
 =\langle f,g\rangle_{L^2},
 \qquad \|\widehat f\|_2=\|f\|_2.
\end{equation}
Fourier 变换唯一延拓为 $L^2(\R^d)$ 上的满射线性等距算子；其逆是反 Fourier 变换的
$L^2$ 延拓。因此该延拓是酉算子。
\end{theorem}

\begin{proof}
若 $f,g\in\mathcal S$，则 $f,g,\widehat f,\widehat g$ 都属于 $L^1\cap L^2$。特别地，
\[
 \int\!\!\int|\widehat f(\xi)g(x)|\dd x\dd\xi
 =\|\widehat f\|_1\|g\|_1<\infty.
\]
故 Fubini 定理与反演公式给出
\begin{align*}
 \langle\widehat f,\widehat g\rangle
 &=\int\widehat f(\xi)
   \left(\int\overline{g(x)}e^{2\pi ix\cdot\xi}\dd x\right)\dd\xi\\
 &=\int\overline{g(x)}
   \left(\int\widehat f(\xi)e^{2\pi ix\cdot\xi}\dd\xi\right)\dd x\\
 &=\int f(x)\overline{g(x)}\dd x
 =\langle f,g\rangle.
\end{align*}
取 $g=f$ 得等距式。

$C_c^\infty(\R^d)\subset\mathcal S$，而第四章的光滑紧支撑稠密性定理
\ref{theorem:4-3-1-3} 说明 $\mathcal S$ 在 $L^2$ 中稠密。给 $f\in L^2$，取
$f_n\in\mathcal S$ 使 $f_n\to f$ 于 $L^2$。等距式给出
\[
 \|\widehat f_n-\widehat f_m\|_2=\|f_n-f_m\|_2,
\]
所以 $(\widehat f_n)$ 在完备的 $L^2$ 中有极限。定义
$\mathcal F_2f=\lim_n\widehat f_n$。同一等式说明定义与逼近列无关、保持线性，并且
$\|\mathcal F_2f\|_2=\|f\|_2$；稠密性还给出这种连续延拓的唯一性。

反 Fourier 变换在 $\mathcal S$ 上也是等距，故以相同步骤延拓为
$\mathcal G_2:L^2\to L^2$。在稠密子空间 $\mathcal S$ 上已有
$\mathcal G_2\mathcal F_2=\mathcal F_2\mathcal G_2=I$。两边都是有界算子；对任意 $L^2$
函数取 Schwartz 逼近并令极限通过算子，得到这些恒等式在整个 $L^2$ 上成立。因此
$\mathcal F_2$ 满射，逆为 $\mathcal G_2$。
\end{proof}

定理中的 $L^2$ Fourier 变换一般不是绝对收敛积分。例如当
$d/2<\alpha\leq d$ 时，$(1+|x|)^{-\alpha}\in L^2(\R^d)\setminus L^1(\R^d)$。
对此类函数，$\mathcal F_2f$ 的含义是任意 Schwartz 逼近列的变换在 $L^2$ 中的极限；换用另一逼近列
不会改变所得等价类，但定理不承诺每个频率点都有绝对积分值。

\begin{proposition}[热半群]\label{proposition:7-2-3-3}
令
\[
 p_t(x)=(4\pi t)^{-d/2}e^{-|x|^2/(4t)},
 \qquad P_tf=p_t*f,\qquad t>0,
\]
并令 $P_0=I$。对每个固定的 $1\leq p<\infty$，$(P_t)_{t\geq0}$ 是
$L^p(\R^d)$ 上强连续、保正的收缩半群。对 $f\in\mathcal S$ 与 $t>0$，
\[
 \partial_tP_tf=\Delta P_tf=P_t\Delta f.
\]
并且在 $L^p$ 中有右导数
\[
 \lim_{t\downarrow0}\frac{P_tf-f}{t}=\Delta f.
\]
若 $G_p$ 表示该 $C_0$ 半群按定义~\ref{def:4-4-3-2} 得到的生成元，则
\begin{equation}\label{eq:7-2-heat-generator-closure}
 G_p=A_p:=\overline{\Delta|_{\mathcal S}}^{\,L^p\times L^p}.
\end{equation}
换言之，
\[
 D(A_p)=\left\{f\in L^p:\begin{array}{l}
   \text{存在 }\varphi_n\in\mathcal S,
   \ \varphi_n\to f\text{ 于 }L^p,\\[-2pt]
   \Delta\varphi_n\text{ 于 }L^p\text{ 收敛}
 \end{array}\right\},
 \quad A_pf=\lim_n\Delta\varphi_n.
\]
特别地，$\mathcal S$ 是 $A_p$ 的核心。此陈述在 $p=1$ 时不把 $D(A_1)$ 改写成
$W^{2,1}$。
\end{proposition}

\begin{proof}
Gaussian 积分给出 $p_t\geq0$ 与 $\int p_t=1$，故 Young 不等式说明
$\|P_tf\|_p\leq\|f\|_p$，且 $P_t$ 保正。由例~\ref{ex:7-2-2-1}，
\[
 \widehat{p_t}(\xi)=e^{-4\pi^2t|\xi|^2}.
\]
于是卷积公式给出
$\widehat{p_s*p_t}=e^{-4\pi^2(s+t)|\xi|^2}=\widehat{p_{s+t}}$。
定理~\ref{theorem:7-2-2-3} 的函数唯一性应用于二者之差，得到
$p_s*p_t=p_{s+t}$，从而 $P_sP_t=P_{s+t}$。又因
$p_t(x)=t^{-d/2}p_1(x/\sqrt t)$ 是近似恒等族，定理~\ref{theorem:4-4-2-1} 给出
$P_tf\to f$ 于 $L^p$，$1\leq p<\infty$。半群律把零点处强连续性移到每个时刻。

若 $f\in\mathcal S$，则
$e^{-4\pi^2t|\xi|^2}\widehat f(\xi)\in\mathcal S$。在任意闭的正时间区间上可逐个 Schwartz
半范数对乘子求导；定理~\ref{theorem:7-2-3-1} 的反变换给出
\begin{align*}
 \mathcal F(\partial_tP_tf)
 &=-4\pi^2|\xi|^2e^{-4\pi^2t|\xi|^2}\widehat f(\xi),\\
 \mathcal F(\Delta P_tf)
 &=-4\pi^2|\xi|^2e^{-4\pi^2t|\xi|^2}\widehat f(\xi).
\end{align*}
Fourier 的单射性给出热方程。并且
\[
 P_tf-f=\int_0^tP_s\Delta f\dd s
\]
作为 $L^p$ 值积分成立。具体地，对 $0<\eta<t$，正时间上的可微性与微积分基本定理给
\[
 P_tf-P_\eta f=\int_\eta^t\partial_sP_sf\dd s
 =\int_\eta^tP_s\Delta f\dd s.
\]
当 $\eta\downarrow0$ 时，$P_\eta f\to f$；又因 $\Delta f\in\mathcal S\subset L^p$，
$P_s\Delta f\to\Delta f$，故右端在 $L^p$ 中收敛到从 $0$ 开始的 Bochner 积分。于是
\[
 \left\|\frac{P_tf-f}{t}-\Delta f\right\|_p
 \leq\frac1t\int_0^t\|P_s\Delta f-\Delta f\|_p\dd s
 \longrightarrow0.
\]
因此
$\mathcal S\subset D(G_p)$ 且 $G_pf=\Delta f$。

生成元 $G_p$ 是闭算子，故 $\Delta|_{\mathcal S}$ 的图闭包 $A_p$ 满足
$A_p\subset G_p$。反过来取 $f\in D(G_p)$。一般 $C_0$ 半群的生成元交换性质
\ref{proposition:4-4-3-1} 给出
\[
 f_\varepsilon:=P_\varepsilon f\in D(G_p),\qquad
 G_pf_\varepsilon=P_\varepsilon G_pf.
\]
当 $\varepsilon\downarrow0$ 时，$f_\varepsilon\to f$ 且
$G_pf_\varepsilon\to G_pf$ 于 $L^p$。

固定 $\varepsilon>0$。取 $g_n\in\mathcal S$ 使 $g_n\to f$ 于 $L^p$，并令
$\varphi_n=P_\varepsilon g_n$。由 Fourier 乘子表示及 Schwartz 自同构，
$\varphi_n\in\mathcal S$。收缩性给出
\[
 \|\varphi_n-f_\varepsilon\|_p\leq\|g_n-f\|_p\longrightarrow0.
\]
此外
\[
 \partial_t p_t(x)
 =\left(\frac{|x|^2}{4t^2}-\frac d{2t}\right)p_t(x)
 =\Delta p_t(x).
\]
当 $t$ 限于 $[\varepsilon/2,3\varepsilon/2]$ 时，右端及时间差商由一个可积的
“多项式乘 Gaussian”控制；支配收敛给出
$\|(p_{\varepsilon+h}-p_\varepsilon)/h-\Delta p_\varepsilon\|_1\to0$。因此
\[
 G_pf_\varepsilon=(\Delta p_\varepsilon)*f,
 \qquad
 \Delta\varphi_n=(\Delta p_\varepsilon)*g_n.
\]
这两个等式也可由 Gaussian 的显式导数和 Fubini 得到。因此
\[
 \|\Delta\varphi_n-G_pf_\varepsilon\|_p
 \leq\|\Delta p_\varepsilon\|_1\|g_n-f\|_p\longrightarrow0.
\]
取 $\varepsilon_k\downarrow0$，并在第 $k$ 步选 $n(k)$，使相应的
$\varphi_k:=P_{\varepsilon_k}g_{n(k)}$ 满足
\[
 \|\varphi_k-f_{\varepsilon_k}\|_p<\frac1k,
 \qquad
 \|\Delta\varphi_k-G_pf_{\varepsilon_k}\|_p<\frac1k.
\]
于是
\begin{align*}
 \|\varphi_k-f\|_p
 &\leq k^{-1}+\|f_{\varepsilon_k}-f\|_p\longrightarrow0,\\
 \|\Delta\varphi_k-G_pf\|_p
 &\leq k^{-1}+\|G_pf_{\varepsilon_k}-G_pf\|_p\longrightarrow0.
\end{align*}
这给出了图闭包所需的同一条逼近列，并证明
$f\in D(A_p)$、$A_pf=G_pf$，故 $G_p\subset A_p$，并完成式
\eqref{eq:7-2-heat-generator-closure} 的证明。上述稠密、收缩与
$\|\Delta p_\varepsilon\|_1$ 估计在 $p=1$ 仍全部成立，所以端点 $p=1$ 已包含在图闭包论证中。
\end{proof}

生成元结论逐个固定 $p<\infty$ 解释。它既不声称不同 $p$ 的定义域相同，也不声称热半群在
$L^\infty$ 上强连续。具体地，令 $f=\1_{\{x_1>0\}}$，并记 $\Phi$ 为标准一维正态分布函数。
热核的坐标乘积结构给出
\[
 P_tf(x)=\Phi\!\left(\frac{x_1}{\sqrt{2t}}\right).
\]
当 $x_1\downarrow0$ 时右侧趋于 $1/2$，而 $f(x)=1$；每个靠近跳面的正侧薄层都有正测度，故
$\|P_tf-f\|_\infty\geq1/2$ 对每个 $t>0$ 成立。于是该半群在 $L^\infty$ 上不强连续。

\begin{proposition}[Poisson 半群与调和延拓]\label{proposition:7-2-3-4}
对 $t>0$，令
\[
 q_t(x)=c_d\frac{t}{(t^2+|x|^2)^{(d+1)/2}},
 \qquad
 c_d=\frac{\Gamma((d+1)/2)}{\pi^{(d+1)/2}},
 \qquad Q_tf=q_t*f.
\]
则对 $s,t>0$，
\[
 \widehat{q_t}(\xi)=e^{-2\pi t|\xi|},
 \qquad Q_sQ_t=Q_{s+t}.
\]
若 $f\in L^p(\R^d)$，$1\leq p<\infty$，则
$Q_tf\to f$ 于 $L^p$；若 $f\in C_0(\R^d)$，则一致收敛。对这些数据，卷积代表
$u(x,t)=Q_tf(x)$ 在上半空间 $\R^d\times(0,\infty)$ 调和：
\[
 (\Delta_x+\partial_t^2)u=0.
\]
\end{proposition}

\begin{proof}
先证明 Poisson 核与热核的从属关系（subordination）。对 $t>0$、$\lambda\geq0$，
\begin{equation}\label{eq:7-2-subordination-scalar}
 \frac{t}{2\sqrt\pi}\int_0^\infty
 s^{-3/2}e^{-t^2/(4s)}e^{-\lambda s}\dd s=e^{-t\sqrt\lambda}.
\end{equation}
为核对该积分，令 $r=t/(2\sqrt s)$，并记
$a=t\sqrt\lambda/2$。左边化为
\[
 \frac2{\sqrt\pi}J(a),
 \qquad J(a)=\int_0^\infty e^{-r^2-a^2/r^2}\dd r.
\]
当 $a>0$ 时可在积分号下求导；再在所得积分中换元 $u=a/r$，有
\[
 J'(a)=-2a\int_0^\infty r^{-2}e^{-r^2-a^2/r^2}\dd r=-2J(a).
\]
而 $J(0)=\sqrt\pi/2$，故 $J(a)=(\sqrt\pi/2)e^{-2a}$；由单调收敛也覆盖 $a=0$。
这证明式~\eqref{eq:7-2-subordination-scalar}。

将热核代入并使用 Tonelli 定理，得到
\begin{align*}
 &\frac{t}{2\sqrt\pi}\int_0^\infty
  s^{-3/2}e^{-t^2/(4s)}p_s(x)\dd s\\
 &\quad=\frac{t}{2\sqrt\pi}(4\pi)^{-d/2}
  \int_0^\infty s^{-(d+3)/2}
  e^{-(t^2+|x|^2)/(4s)}\dd s\\
 &\quad=\frac{\Gamma((d+1)/2)}{\pi^{(d+1)/2}}
       \frac{t}{(t^2+|x|^2)^{(d+1)/2}}=q_t(x),
\end{align*}
其中最后一步使用 $r=(t^2+|x|^2)/(4s)$。再用 Fubini、热核 Fourier 公式与
式~\eqref{eq:7-2-subordination-scalar}（取 $\lambda=4\pi^2|\xi|^2$），得到
\[
 \widehat{q_t}(\xi)
 =\frac{t}{2\sqrt\pi}\int_0^\infty
 s^{-3/2}e^{-t^2/(4s)}e^{-4\pi^2s|\xi|^2}\dd s
 =e^{-2\pi t|\xi|}.
\]
取 $\xi=0$ 可见 $\int q_t=1$。乘子相乘并使用 $L^1$ Fourier 唯一性，得到
$q_s*q_t=q_{s+t}$，即 Poisson 半群律。

直接核对调和性。置 $\rho=t^2+|x|^2$、$\alpha=(d+1)/2$，暂略常数 $c_d$。对
$q=t\rho^{-\alpha}$，
\begin{align*}
 \Delta_xq
 &=-2\alpha dt\rho^{-\alpha-1}
   +4\alpha(\alpha+1)t|x|^2\rho^{-\alpha-2},\\
 \partial_t^2q
 &=-6\alpha t\rho^{-\alpha-1}
   +4\alpha(\alpha+1)t^3\rho^{-\alpha-2}.
\end{align*}
两式相加后，$2(\alpha+1)=d+3$ 使系数抵消，故
$(\Delta_x+\partial_t^2)q_t=0$。在任意正时间带
$0<a\leq t\leq b$ 上，对实际用到的 $|\gamma|+j\leq3$，直接对
$t(t^2+|x|^2)^{-(d+1)/2}$ 求导并把有界的 $t$ 次幂吸收进常数，可得
\[
 \sup_{a\leq t\leq b}
 \bigl|\partial_t^j\partial_x^\gamma q_t(x)\bigr|
 \leq C_{a,b,\gamma,j}(1+|x|)^{-d-1}.
\]
右端属于每个 $L^r(\R^d)$，$1\leq r\leq\infty$。下面把端点所需的导数交换写清。若
$1<p<\infty$，令 $p'=p/(p-1)$；对任一这样的核导数 $k$，其平移在 $L^{p'}$ 中连续，故差商恒等式
\[
 \frac{k(\,\cdot+h e_j)-k}{h}
 =\int_0^1\partial_jk(\,\cdot+\theta h e_j)\dd\theta
\]
给出差商在 $L^{p'}$ 中趋于 $\partial_jk$，再由 \Holder{} 不等式把极限送入与 $f$ 的卷积。
当 $p=1$ 时，核的下一阶导数连续且在无穷远消失，因而一致连续；同一恒等式给出差商在
$L^\infty$ 中收敛，再与 $f\in L^1$ 配对。对时间差商及二阶导数重复此论证，得到
\[
 \partial_t^2(q_t*f)=(\partial_t^2q_t)*f,
 \qquad \Delta_x(q_t*f)=(\Delta_xq_t)*f.
\]
这些卷积随 $(x,t)$ 连续，因此 $u=Q_tf$ 是经典调和函数。若 $f\in C_0\subset L^\infty$，则改用
核导数在 $L^1$ 中的平移连续性与 $L^1$ 差商收敛，结论相同。

最后，$q_t(x)=t^{-d}q_1(x/t)$，且 $q_1\geq0$、$\int q_1=1$。这是近似恒等族。
定理~\ref{theorem:4-4-2-1} 分别给出 $L^p$（$p<\infty$）强收敛与 $C_0$ 一致收敛。
\end{proof}

这里没有断言 $Q_t$ 在 $L^\infty$ 上强连续。以
$f=\1_{\{x_1>0\}}$ 为例，对其余坐标积分后得到一维 Poisson 核，因而
\[
 Q_tf(x)=\frac12+\frac1\pi\arctan\frac{x_1}{t}.
\]
当 $x_1\downarrow0$ 时右侧趋于 $1/2$，而 $f(x)=1$；正侧任意薄层都有正测度，所以
$\|Q_tf-f\|_\infty\geq1/2$ 对每个 $t>0$ 成立，强连续性确实失败。

\begin{example}[Gaussian 与 Hermite]\label{ex:7-2-3-1}
函数 $g(x)=e^{-\pi|x|^2}$ 属于 $\mathcal S$，且
$\widehat g=g$。命题~\ref{proposition:7-2-2-2} 给出具体的第一阶关系
\[
 \widehat{x_jg}(\xi)=-i\xi_jg(\xi),
 \qquad
 \widehat{\partial_jg}(\xi)=2\pi i\xi_jg(\xi).
\]
反复乘 $x_j$ 与求导，会得到“多项式乘 Gaussian”的函数族；Leibniz 公式说明它们全在
$\mathcal S$。令
\[
 V_N=\{P(x)g(x):\deg P\leq N\}.
\]
对 $|\alpha|\leq N$，乘法--微分公式逐项给出
\[
 \widehat{x^\alpha g}(\xi)
 =\left(\frac{i}{2\pi}\right)^{|\alpha|}\partial^\alpha g(\xi)
 =Q_\alpha(\xi)g(\xi),
 \qquad \deg Q_\alpha\leq|\alpha|.
\]
故 $\mathcal F(V_N)\subset V_N$。反 Fourier 变换等于 Fourier 变换后再作反射，而反射保持
$V_N$，所以实际有 $\mathcal F(V_N)=V_N$。这些 $V_N$ 是按次数递增的有限维子空间。按总次数对
“多项式乘 Gaussian”作正交化可得到 Hermite 函数；又因 $\mathcal F$ 在 $L^2$ 上为酉算子，它保持
$V_N\ominus V_{N-1}$（约定 $V_{-1}=\{0\}$），故每一总次数层内的 Fourier 作用都归结为有限维酉
矩阵。因而 Hermite 函数提供了可显式计算的 Fourier 测试族，并将在谐振子谱问题中再次出现。
\end{example}

\begin{example}[Dirac 分布]\label{ex:7-2-3-2}
定义 $\langle\delta_0,\varphi\rangle=\varphi(0)$。由于
$|\varphi(0)|\leq p_{0,0}(\varphi)$，有 $\delta_0\in\mathcal S'$。按定义
\ref{def:7-2-3-4}，
\[
 \langle\widehat{\delta_0},\varphi\rangle
 =\langle\delta_0,\widehat\varphi\rangle
 =\widehat\varphi(0)=\int\varphi(x)\dd x
 =\langle1,\varphi\rangle.
\]
所以 $\widehat{\delta_0}=1$。反过来，Schwartz 反演给出
\[
 \langle\widehat1,\varphi\rangle
 =\int\widehat\varphi(x)\dd x=\varphi(0)
 =\langle\delta_0,\varphi\rangle,
\]
故 $\widehat1=\delta_0$。这些等式只表示两个缓增分布对每个 Schwartz 测试函数的作用相同，并未把
$\delta_0$ 当作普通点值函数。
\end{example}

\begin{example}[热核]\label{ex:7-2-3-3}
把例~\ref{ex:7-2-2-1} 中的参数取为 $a=(4\pi t)^{-1}$，有
\begin{align*}
 \widehat{p_t}(\xi)
 &=(4\pi t)^{-d/2}
   (4\pi t)^{d/2}e^{-4\pi^2t|\xi|^2}\\
 &=e^{-4\pi^2t|\xi|^2}.
\end{align*}
因此 $p_s*p_t=p_{s+t}$ 在频率侧只是
$e^{-4\pi^2s|\xi|^2}e^{-4\pi^2t|\xi|^2}
=e^{-4\pi^2(s+t)|\xi|^2}$。同一个乘子求时间导数给出
$-4\pi^2|\xi|^2$，正好是 $\Delta$ 的 Fourier 符号。
\end{example}

\begin{example}[Brownian 与热半群的同一计算]\label{ex:7-2-3-4}
若第六章的标准 Brownian motion 满足 $B_t\sim N(0,tI)$，则其密度是
\[
 (2\pi t)^{-d/2}e^{-|y|^2/(2t)}=p_{t/2}(y).
\]
所以对有界 Borel 函数 $f$，以及对任意 $f\in L^p(\R^d)$、$1\le p\le\infty$
（此时取卷积所定义的代表元），
\[
 \E f(x+B_t)=\int f(x+y)p_{t/2}(y)\dd y=P_{t/2}f(x),
\]
其中最后一步使用 $p_{t/2}(y)=p_{t/2}(-y)$。若 $f\in\mathcal S$，写
$B_t=\sqrt t\,Z$，其中 $Z\sim N(0,I_d)$。二阶 Taylor 公式给
\[
 f(x+y)-f(x)
 =\nabla f(x)\cdot y+\frac12y^{\mathsf T}D^2f(x)y+r(x,y),
 \qquad
 |r(x,y)|\leq\frac12|y|^2\omega_{D^2f}(|y|),
\]
其中 $\omega_{D^2f}(s)\to0$ 当 $s\downarrow0$。利用
$\E Z_j=0$、$\E Z_jZ_k=\delta_{jk}$，得到
\[
 \sup_x\left|
   \frac{\E f(x+B_t)-f(x)}t-\frac12\Delta f(x)
 \right|
 \leq\frac12\E\!\left[|Z|^2
        \omega_{D^2f}(\sqrt t|Z|)\right]\longrightarrow0.
\]
最后一步使用支配收敛，因为
$\omega_{D^2f}\leq2\|D^2f\|_\infty$ 且 $\E|Z|^2<\infty$。特别地，
\[
 \lim_{t\downarrow0}\frac{\E f(x+B_t)-f(x)}t
 =\frac12\Delta f(x).
\]
因此概率时间 $t$ 对应本节 PDE 热时间 $t/2$：Brownian 生成元是
$\tfrac12\Delta$，而 $P_t$ 的生成元是 $\Delta$。方差参数中的这个 $1/2$ 不能由记号省略。
\end{example}

图~\ref{fig:7-2-schwartz-plancherel-semigroups} 汇总本小节的三层结构。内层的 Schwartz 方块允许
逐点积分、微分与反演；垂直的稠密完备化把它唯一扩张为 $L^2$ 上的酉算子；右侧两个指数乘子则把连续
时间演化化成乘法。生成元属于无界算子层，必须连同图闭包定义域理解。

\begin{figure}[htbp]
\centering
\begin{tikzpicture}[node distance=1.25cm and 2.0cm]
  \node[book strong node] (sx) {$\mathcal S_x$};
  \node[book strong node,right=of sx] (sxi) {$\mathcal S_\xi$};
  \draw[book accent arrow,bend left=14] (sx) to node[above,book label]{$\mathcal F$} (sxi);
  \draw[book accent arrow,bend left=14] (sxi) to node[below,book label]{$\mathcal F^{-1}$} (sx);

  \node[book node,below=1.55cm of sx] (l2x) {$L^2_x$};
  \node[book node,below=1.55cm of sxi] (l2xi) {$L^2_\xi$};
  \draw[book dashed arrow] (sx)--node[left,book label]{稠密完备化} (l2x);
  \draw[book dashed arrow] (sxi)--node[right,book label]{稠密完备化} (l2xi);
  \draw[book arrow,bend left=12] (l2x) to node[above,book label]{酉} (l2xi);
  \draw[book arrow,bend left=12] (l2xi) to (l2x);

  \node[book warm node,right=2.0cm of sxi] (heat)
    {热乘子\\$e^{-4\pi^2t|\xi|^2}$};
  \node[book warm node,below=1.2cm of heat] (pois)
    {Poisson 乘子\\$e^{-2\pi t|\xi|}$};
  \draw[book arrow] (sxi)--(heat);
  \draw[book arrow] (sxi)--(pois);
  \node[book warning node,below=1.0cm of pois] (gen)
    {生成元\\图闭包与定义域};
  \draw[book dashed arrow] (heat)--(gen);
\end{tikzpicture}
\caption{Schwartz 自同构、$L^2$ 完备化与热/Poisson 乘子。}
\label{fig:7-2-schwartz-plancherel-semigroups}
\end{figure}

\begin{exercise}
\begin{enumerate}
  \item 用 Leibniz 公式逐项证明 $\mathcal S(\R^d)$ 对微分和乘多项式闭合；再对
        $f,g\in\mathcal S$ 把积分分成 $|y|\le |x|/2$ 与其补集，证明 $f*g\in\mathcal S$。
  \item 从恒等式
        $\xi^\alpha\partial_\xi^\beta\widehat f
        =(2\pi i)^{-|\alpha|}(-2\pi i)^{|\beta|}
        \widehat{\partial^\alpha(x^\beta f)}$
        出发，对每个输出半范数写出有限个输入半范数的显式控制，从而证明
        $\mathcal F:\mathcal S\to\mathcal S$ 连续。
  \item 只用 $\mathcal S$ 在 $L^2$ 中稠密、Parseval 等式与 $L^2$ 完备性，构造 Fourier 的
        $L^2$ 延拓；逐步证明定义与逼近列无关，并用反 Fourier 延拓证明其像为整个 $L^2$。
  \item 从乘子 $e^{-4\pi^2t|\xi|^2}$ 推出热算子的半群律。对 $f\in\mathcal S$ 证明乘子可在
        $t$ 上求导并得到 $\partial_tP_tf=\Delta P_tf$；再指出近似恒等性质为何只给
        $1\le p<\infty$ 上的强连续性，而不自动给 $L^\infty$ 情形。
\end{enumerate}
\end{exercise}

圆周与欧氏空间的 Fourier 理论都从所有指数测试函数上的积分值唯一恢复对象：圆周用 \Fejer{} 正核和离散
$\ell^2$ 完备化，欧氏空间用 Gaussian 正则化和连续 $L^2$ 完备化。下一节将不再预先给出测度，
而从一列候选矩出发，研究正定性与紧性何时足以保证代表测度存在。

\section{矩、弱收敛与测度识别}\label{sec:07-03}

一个有限测度可以在许多不同的测试函数族上积分：多项式给出矩，复指数给出 Fourier 系数，半直线指标给出
分布函数。测试函数族太小时，不同测度可能无法区分；即使测试函数族足够大，还要回答这些数值是否真的来自
某个测度。本节把这两个问题分开处理。第一小节研究矩数据的正性与表示，第二小节研究测度列的选择
与弱收敛，最后一小节比较几类变换的识别能力，并用 Carleman 条件修补无界支集上的矩不唯一性。

Hamburger、Stieltjes、Hausdorff 与 Herglotz 的矩问题分别把支集限制在
$\R$、$[0,\infty)$、$[0,1]$ 与圆周 $\T$。它们都可写成“正线性泛函是否来自积分”，但支集不同，
可检验的正性条件也不同。尤其不能把无界支集问题中的存在性、质量控制与唯一性压缩成同一句表示
定理。

\subsection{矩问题、正性与 Hankel/Toeplitz 结构}\label{subsec:7-3-1}
\index{矩问题}\index{Hankel 矩阵}\index{Toeplitz 矩阵}\index{Herglotz 表示}

先从最小的可计算模型开始。若 $\mu$ 是 $\R$ 上的有限正测度并且所需矩有限，则任意多项式
$p$ 都满足 $\int |p|^2\,\dd\mu\geq0$。展开平方后，这个不等式只涉及有限多个矩；因此每次检验
都是有限维矩阵的半正定性。圆周上的三角多项式给出同样的机制，只是矩阵由 Hankel 型变成
Toeplitz 型。

\begin{example}[Dirac 矩]\label{ex:7-3-1-1}
对 $\mu=\delta_a$，
\[
 c_n=\int_\R x^n\,\delta_a(\dd x)=a^n.
\]
令 $v_N(a)=(1,a,\ldots,a^N)^{\mathsf T}$，则
\[
 H_N=(c_{j+k})_{0\leq j,k\leq N}
     =v_N(a)v_N(a)^{\mathsf T}.
\]
所以 $H_N$ 半正定且秩为一。更一般地，若
$\mu=\sum_{\ell=1}^m w_\ell\delta_{a_\ell}$，其中 $w_\ell>0$ 且 $a_\ell$ 两两不同，则
\[
 H_N=V_N\operatorname{diag}(w_1,\ldots,w_m)V_N^{\mathsf T},
 \qquad (V_N)_{j\ell}=a_\ell^j.
\]
Vandermonde 行列式说明 $N\geq m-1$ 时 $V_N$ 的列线性无关，因而
$\operatorname{rank}H_N=m$；对任意 $N$，秩至多为 $m$。矩阵的有限秩正是有限原子数留下的代数
痕迹。
\end{example}

\begin{example}[圆周正定序列]\label{ex:7-3-1-2}
设 $\mu$ 是 $\T=\R/(2\pi\Z)$ 上的有限正测度，并令
\[
 c_n=\int_\T e^{-int}\,\mu(\dd t),\qquad n\in\Z.
\]
于是 $c_{-n}=\overline{c_n}$。对任意 $a_0,\ldots,a_N\in\C$，直接展开得
\begin{align*}
 \sum_{j,k=0}^N\overline{a_j}c_{j-k}a_k
 &=\int_\T\sum_{j,k=0}^N
   \overline{a_j}a_k e^{-i(j-k)t}\,\mu(\dd t)\\
 &=\int_\T\left|\sum_{k=0}^N a_ke^{ikt}\right|^2\mu(\dd t)\geq0.
\end{align*}
下文的 Herglotz 定理说明，这些有限二次型的正性不仅必要，而且足以构造出唯一的 $\mu$。
\end{example}

\begin{example}[对数正态不定矩]\label{ex:7-3-1-3}
令 $Y\sim N(0,1)$、$X=e^Y$。$X$ 的密度为
\[
 f(x)=\frac{1}{x\sqrt{2\pi}}
      \exp\!\left(-\frac{(\log x)^2}{2}\right),\qquad x>0,
\]
且
\[
 \int_0^\infty x^n f(x)\,\dd x=\E e^{nY}=e^{n^2/2},
 \qquad n=0,1,2,\ldots.
\]
对 $|\varepsilon|\leq1$ 定义 Stieltjes 族
\[
 f_\varepsilon(x)=f(x)\bigl[1+\varepsilon\sin(2\pi\log x)\bigr].
\]
括号内非负。对每个整数 $n\geq0$，扰动项的 $n$ 阶矩为
\begin{align*}
 \int_0^\infty x^nf(x)\sin(2\pi\log x)\,\dd x
 &=\Im\E e^{(n+2\pi i)Y}\\
 &=\Im\exp\!\left(\frac{(n+2\pi i)^2}{2}\right)\\
 &=e^{(n^2-4\pi^2)/2}\sin(2\pi n)=0.
\end{align*}
取 $n=0$ 还说明 $f_\varepsilon$ 的积分为一。因此所有 $f_\varepsilon$ 都是概率密度，并共享矩
$e^{n^2/2}$；当 $\varepsilon\neq\eta$ 时，两密度在一个正 Lebesgue 测度集合上不同。这个计算表明：
即使每一阶矩都存在且全部相同，分布律仍可能不同。
\end{example}

\begin{definition}[矩泛函]\label{def:7-3-1-1}
给定实数列 $(c_n)_{n\geq0}$，在 $\C[x]$ 上定义
\[
 L\!\left(\sum_{k=0}^N a_kx^k\right)=\sum_{k=0}^Na_kc_k.
\]
对多项式 $p(x)=\sum a_kx^k$，记
$\overline p(x)=\sum\overline{a_k}x^k$。若
\[
 L(\overline p\,p)\geq0\qquad(p\in\C[x]),
\]
则称 $L$ 为\emph{正矩泛函}。
\end{definition}

\begin{definition}[Hankel 与 Toeplitz 矩阵]\label{def:7-3-1-2}
单边矩序列 $(c_n)_{n\geq0}$ 的第 $N$ 个 \emph{Hankel 矩阵}是
\[
 H_N=(c_{j+k})_{0\leq j,k\leq N}.
\]
满足 $c_{-n}=\overline{c_n}$ 的双边序列的第 $N$ 个 \emph{Toeplitz 矩阵}是
\[
 T_N=(c_{j-k})_{0\leq j,k\leq N}.
\]
所谓所有截断半正定，是指 $a^*H_Na\geq0$ 或 $a^*T_Na\geq0$ 对每个 $N$ 与每个相应有限向量
$a$ 成立。
\end{definition}

\begin{definition}[确定与不确定矩问题]\label{def:7-3-1-3}
若一个矩序列存在满足指定支集条件的正表示测度，并且该测度唯一，则称相应矩问题是
\emph{确定的}（determinate）；若存在至少两个不同的表示测度，则称为
\emph{不确定的}（indeterminate）。没有表示测度的序列不归入这两个名称。存在性与确定性是两个独立
问题：先要证明表示存在，随后才谈得上它是否唯一。
\end{definition}

\begin{proposition}[Hankel 正性必要性]\label{proposition:7-3-1-1}
设 $\mu$ 是 $\R$ 上的正测度，且 $c_n=\int x^n\,\dd\mu$ 对所有 $n\geq0$ 有限。则每个
$H_N$ 半正定。若进一步 $\supp\mu\subset[0,\infty)$，则每个移位 Hankel 矩阵
\[
 H_N^{(1)}=(c_{j+k+1})_{0\leq j,k\leq N}
\]
也半正定。
\end{proposition}

\begin{proof}
对 $p(x)=\sum_{k=0}^Na_kx^k$，有限求和可直接移入积分，故
\begin{align*}
 a^*H_Na
 &=\sum_{j,k=0}^N\overline{a_j}a_kc_{j+k}
   =\int_\R\sum_{j,k=0}^N\overline{a_j}a_kx^{j+k}\,\mu(\dd x)\\
 &=\int_\R|p(x)|^2\,\mu(\dd x)\geq0.
\end{align*}
若 $\mu$ 支持在 $[0,\infty)$，同样的展开给出
\[
 a^*H_N^{(1)}a=\int_{[0,\infty)}x|p(x)|^2\,\mu(\dd x)\geq0.
\]
两项结论分别成立。
\end{proof}

这里必须保留支集条件。经典 Hamburger 判据以所有 $H_N\succeq0$ 刻画 $\R$ 上的表示；经典
Stieltjes 判据还要求所有 $H_N^{(1)}\succeq0$，以刻画 $[0,\infty)$ 上的表示。完整的充分性证明
需要无界对称算子的方法。圆周与闭区间情形中支集是紧的，因而可以直接用有限正测度和紧性完成构造。

\begin{theorem}[Herglotz 表示]\label{theorem:7-3-1-2}
设双边序列 $(c_n)_{n\in\Z}$ 满足 $c_{-n}=\overline{c_n}$。存在唯一的 $\T$ 上有限正测度
$\mu$ 使
\[
 c_n=\int_\T e^{-int}\,\mu(\dd t),\qquad n\in\Z,
\]
当且仅当所有 Toeplitz 截断 $T_N$ 半正定。
\end{theorem}

\begin{proof}
必要性已在例~\ref{ex:7-3-1-2} 的有限平方展开中证明。以下证明充分性。

先分离零质量情形。对 $n\geq1$，$T_n$ 的第 $0,n$ 行列所成主子式是
\[
 \begin{pmatrix}c_0&c_{-n}\\c_n&c_0\end{pmatrix}.
\]
半正定性给 $c_0\geq0$ 以及
$c_0^2-|c_n|^2\geq0$。若 $c_0=0$，则 $c_n=0$ 对所有 $n\geq1$，再由共轭对称性可知整个序列
为零；零测度给出表示。以下设 $c_0>0$。

对 $N\geq0$ 定义三角多项式
\begin{equation}\label{eq:7-3-herglotz-fejer-density}
 P_N(t)=\frac{1}{N+1}\sum_{j,k=0}^Nc_{j-k}e^{i(j-k)t}.
\end{equation}
把 $a_k=e^{-ikt}$ 代入 $T_N$ 的二次型，得到
$P_N(t)=(N+1)^{-1}a^*T_Na\geq0$。于是
\[
 \mu_N(\dd t)=P_N(t)\frac{\dd t}{2\pi}
\]
是有限正测度。按差 $m=j-k$ 收集式~\eqref{eq:7-3-herglotz-fejer-density} 中的项，有
\[
 P_N(t)=\sum_{|m|\leq N}
 \left(1-\frac{|m|}{N+1}\right)c_me^{imt}.
\]
圆周指数函数的正交性因而给出
\begin{equation}\label{eq:7-3-herglotz-fejer-coefficients}
 \int_\T e^{-int}\,\mu_N(\dd t)
 =
 \begin{cases}
 \left(1-\dfrac{|n|}{N+1}\right)c_n,& |n|\leq N,\\[4pt]
 0,& |n|>N.
 \end{cases}
\end{equation}
特别地，$\mu_N(\T)=c_0$。概率测度 $c_0^{-1}\mu_N$ 全部支持在紧空间 $\T$ 上，故由
Prokhorov 定理~\ref{theorem:5-3-3-2} 可取子列 $N_r\to\infty$，使
$c_0^{-1}\mu_{N_r}$ 弱收敛于某个概率 $\lambda$。令 $\mu=c_0\lambda$。对每个固定 $n\in\Z$，
函数 $t\mapsto e^{-int}$ 有界连续；弱收敛与
式~\eqref{eq:7-3-herglotz-fejer-coefficients} 给出
\[
 \int_\T e^{-int}\,\mu(\dd t)
 =\lim_{r\to\infty}\int_\T e^{-int}\,\mu_{N_r}(\dd t)=c_n.
\]
这就得到表示。

若 $\mu$ 与 $\nu$ 都表示 $(c_n)$，则有限复测度 $\mu-\nu$ 的全部圆周 Fourier 系数为零。
圆周 Fourier 唯一性命题~\ref{proposition:7-2-1-4} 给出 $\mu-\nu=0$，所以表示唯一。
\end{proof}

\begin{theorem}[Hausdorff 矩判据]\label{theorem:7-3-1-3}
实数列 $(c_n)_{n\geq0}$ 是 $[0,1]$ 上某个有限正测度的矩序列，当且仅当
\begin{equation}\label{eq:7-3-hausdorff-complete-monotonicity}
 (-1)^k\Delta^kc_n\geq0\qquad(n,k\geq0),
 \qquad \Delta c_n=c_{n+1}-c_n.
\end{equation}
表示测度唯一，其总质量为 $c_0$。
\end{theorem}

\begin{proof}
若 $c_n=\int_0^1x^n\,\mu(\dd x)$，则有限差分可以逐项移入积分。由
$\Delta^kx^n=x^n(x-1)^k$ 得
\[
 (-1)^k\Delta^kc_n
 =\int_0^1x^n(1-x)^k\,\mu(\dd x)\geq0.
\]
这证明必要性。

反过来，假设式~\eqref{eq:7-3-hausdorff-complete-monotonicity} 成立。若 $c_0=0$，则
$c_n\geq0$，而 $c_n-c_{n+1}=-\Delta c_n\geq0$；因此
$0\leq c_n\leq c_0=0$，零测度即为所求。以下设 $c_0>0$。

对每个 $N\geq1$ 及 $0\leq k\leq N$ 置
\begin{equation}\label{eq:7-3-hausdorff-weights}
 p_{N,k}=\binom Nk(-1)^{N-k}\Delta^{N-k}c_k,
 \qquad
 \mu_N=\sum_{k=0}^Np_{N,k}\delta_{k/N}.
\end{equation}
假设保证 $p_{N,k}\geq0$。为计算这些权重，不妨仍以 $L(x^j)=c_j$ 记系数泛函，并令
\[
 B_{N,k}(x)=\binom Nkx^k(1-x)^{N-k}
\]
为 Bernstein 基函数。展开 $(1-x)^{N-k}$ 可得
\begin{align*}
 L(B_{N,k})
 &=\binom Nk\sum_{j=0}^{N-k}(-1)^j\binom{N-k}{j}c_{k+j}\\
 &=\binom Nk(-1)^{N-k}\Delta^{N-k}c_k=p_{N,k}.
\end{align*}
Bernstein 恒等式 $\sum_{k=0}^NB_{N,k}=1$ 因而给
\[
 \sum_{k=0}^Np_{N,k}=L(1)=c_0.
\]
更精确地，记下降阶乘 $(k)_r=k(k-1)\cdots(k-r+1)$。二项分布的有限代数恒等式
\begin{equation}\label{eq:7-3-bernstein-factorial-identity}
 \sum_{k=0}^N(k)_rB_{N,k}(x)=(N)_rx^r
\end{equation}
可由
$(k)_r\binom Nk=(N)_r\binom{N-r}{k-r}$ 及二项式定理直接得到。对
式~\eqref{eq:7-3-bernstein-factorial-identity} 施加 $L$，便有
\begin{equation}\label{eq:7-3-hausdorff-factorial-moments}
 \sum_{k=0}^N(k)_rp_{N,k}=(N)_rc_r.
\end{equation}

令 $S(r,j)$ 为第二类 Stirling 数，即
$k^r=\sum_{j=0}^rS(r,j)(k)_j$。当 $N\geq r$ 时，
式~\eqref{eq:7-3-hausdorff-factorial-moments} 给出
\begin{align*}
 \int_0^1x^r\,\mu_N(\dd x)
 &=\frac1{N^r}\sum_{k=0}^Nk^rp_{N,k}\\
 &=\sum_{j=0}^rS(r,j)\frac{(N)_j}{N^r}c_j
 \longrightarrow c_r.
\end{align*}
最后一步使用 $S(r,r)=1$、$(N)_r/N^r\to1$，而 $j<r$ 时
$(N)_j/N^r\to0$。这一步说明普通幂矩不是被假定出来的，而是由下降阶乘矩逐阶恢复。

归一化测度 $c_0^{-1}\mu_N$ 全部支持在共同紧集 $[0,1]$ 上。再次应用
定理~\ref{theorem:5-3-3-2}，可取子列使
$c_0^{-1}\mu_{N_\ell}\Rightarrow\lambda$，其中 $\lambda$ 是概率测度。令 $\mu=c_0\lambda$，则
$\mu_{N_\ell}$ 对所有有界连续测试积分收敛于 $\mu$。每个 $x^r$ 在 $[0,1]$ 上有界连续，所以
\[
 \int_0^1x^r\,\mu(\dd x)
 =\lim_{\ell\to\infty}\int_0^1x^r\,\mu_{N_\ell}(\dd x)=c_r.
\]
存在性得证。

若 $\mu,\nu$ 是两个表示测度，则它们对每个多项式的积分相同。多项式由
Stone--Weierstrass 定理~\ref{theorem:1-3-3-2} 在 $C([0,1])$ 中一致稠密；又因两测度有限，
积分关于一致范数连续，故它们对每个连续函数的积分都相同。Radon 开集公式
\ref{theorem:3-1-1-4} 随即给出 $\mu=\nu$。唯一性得证。
\end{proof}

\begin{proposition}[GNS 原型]\label{proposition:7-3-1-4}
设 $L$ 是正矩泛函。公式
\[
 \langle p,q\rangle_L=L(\overline p\,q)
\]
在 $\C[x]$ 上定义半内积。其零空间
$\mathcal N=\{p:L(\overline p\,p)=0\}$ 在乘以 $x$ 后不变。因此
\[
 M_x[p]=[xp]
\]
在内积空间 $\C[x]/\mathcal N$ 上良定义；把该商空间完备化为 $\mathcal H_L$ 后，$M_x$ 是定义
在稠密子空间 $\C[x]/\mathcal N$ 上的对称算子。
\end{proposition}

\begin{proof}
因为矩 $c_n$ 为实数，$\langle p,q\rangle_L$ 是 Hermitian 半正定型。正性先给出
Cauchy--Schwarz 不等式。记
$a=L(\overline p p)$、$b=L(\overline q q)$、
$c=L(\overline p q)$。若 $b>0$，在
\[
 0\leq L(\overline{(p+\lambda q)}(p+\lambda q))
 =a+\lambda c+\overline\lambda\,\overline c+|\lambda|^2b
\]
中取 $\lambda=-\overline c/b$，便得
$0\leq a-|c|^2/b$，即
\[
 |L(\overline p q)|^2
 \leq L(\overline p p)L(\overline q q).
\]
若 $b=0$，则对 $t\in\R$ 分别取 $\lambda=t$ 与 $\lambda=it$，得到
\[
 0\leq a+2t\Re c,
 \qquad
 0\leq a-2t\Im c
 \qquad(t\in\R).
\]
让 $t$ 分别取任意大的正、负值，迫使 $\Re c=\Im c=0$。因此 $c=0$，同一不等式仍成立。

现在取 $p\in\mathcal N$。Cauchy--Schwarz 给
\[
 L(\overline p\,r)=0\qquad(r\in\C[x]).
\]
关键是不能只在这里声称零空间自动为理想；令 $r=x^2p$，才得到
\[
 L(\overline{xp}\,xp)=L(x^2\overline p\,p)
 =L(\overline p\,x^2p)=0.
\]
它还说明 $\mathcal N$ 是该半内积型的根空间，特别是线性子空间。上式给出
$xp\in\mathcal N$。若 $[p]=[q]$，则 $p-q\in\mathcal N$，从而
$x(p-q)\in\mathcal N$；所以 $M_x[p]=[xp]$ 与代表元无关。

最后，由 $x=\overline x$ 以及多项式乘法可交换，
\[
 \langle M_x[p],[q]\rangle_L
 =L(x\overline p\,q)
 =L(\overline p\,xq)
 =\langle[p],M_x[q]\rangle_L.
\]
这证明 $M_x$ 对称。商空间在其完备化中稠密，故定义域稠密。这里没有断言 $M_x$ 有界、闭或自伴；
这些性质以及自伴扩张所产生的谱测度属于后续泛函分析。
\end{proof}

\begin{figure}[htbp]
\centering
\begin{tikzpicture}[node distance=10mm and 18mm]
  \node[book strong node] (mom) {矩数据\\$(c_n)$};
  \node[book node,right=of mom] (mat) {有限截断\\$H_N$ 或 $T_N$};
  \node[book warm node,right=of mat] (pos) {二次型正性\\有限维检验};
  \node[book node,below=of mat] (ip) {半内积\\$\langle p,q\rangle_L$};
  \node[book node,below=of pos] (approx) {有限正测度\\$\mu_N$};
  \node[book strong node,below=of ip] (gns) {GNS 原型\\商空间与乘法算子};
  \node[book strong node,below=of approx] (meas) {表示测度\\$\mu$};
  \node[book warning node,below=of meas] (boundary) {支集与唯一性\\需另行检验};
  \draw[book arrow] (mom)--(mat);
  \draw[book arrow] (mat)--(pos);
  \draw[book arrow] (pos)--(ip);
  \draw[book arrow] (pos)--(approx);
  \draw[book dashed arrow] (ip)--(gns);
  \draw[book accent arrow] (approx)--node[right,book label]{紧性极限}(meas);
  \draw[book dashed arrow] (meas)--(boundary);
\end{tikzpicture}
\caption{矩问题的两条有限维构造途径。矩阵正性一方面产生 GNS 半内积，另一方面在 Herglotz 与 Hausdorff
问题中直接产生有限正测度；紧性只负责取极限，支集和唯一性仍需各自的检验。}
\label{fig:7-3-moment-positive-routes}
\end{figure}

这张图也解释了正定核与时间序列协方差的表示入口：每个有限样本只给出半正定矩阵；若这些矩阵彼此
相容，表示定理才把它们组织成一个测度或过程。正交多项式、\Pade{} 逼近与连分数则从 GNS 商空间中
对 $1,x,x^2,\ldots$ 做正交化开始。这些联系由上述有限维正性与紧性构造已经得到解释；
一般谱定理将在更完整的算子理论中给出它们的统一表述。

\begin{exercise}
设 $\mu=\sum_{\ell=1}^mw_\ell\delta_{a_\ell}$，其中权重为正、原子互异。证明
$\operatorname{rank}H_N=\min\{m,N+1\}$，并指出若允许权重有正有负，证明中的哪一步失效。
\end{exercise}
\begin{exercise}
验证 $[0,\infty)$ 支集给出移位 Hankel 正性，并构造一个 Hankel 半正定但移位 Hankel 不半正定的
矩序列。
\end{exercise}
\begin{exercise}
在 Herglotz 定理中直接计算 $P_N$ 的第 $n$ 个 Fourier 系数，并解释为何必须先把 $c_0=0$ 与
$c_0>0$ 分开。
\end{exercise}

矩表示回答“数据来自什么测度”。接下来转向另一个问题：给定一列已经存在的测度，哪些有限观测足以
保证它们有极限，并且极限没有在无穷远损失质量？

\subsection{Helly 选择、分布函数与弱收敛}\label{subsec:7-3-2}
\index{Helly 选择定理}\index{分布函数}\index{弱收敛}\index{模糊收敛}

实线上的概率测度可以压缩成一个递增函数。这个编码的优势是：在每个有理点取值都是一个有界数，
因此可以用可数对角法抽取子列；它的代价是跳点处不能要求普通点态收敛。移动原子把这一边界展示得
最清楚。

\begin{example}[移动 Dirac]\label{ex:7-3-2-1}
令 $\mu_n=\delta_{1/n}$。对每个 $f\in C_b(\R)$，
\[
 \int f\,\dd\mu_n=f(1/n)\longrightarrow f(0)=\int f\,\dd\delta_0,
\]
故 $\delta_{1/n}\Rightarrow\delta_0$。相应分布函数为
\[
 F_n(x)=\1_{[1/n,\infty)}(x),
 \qquad F(x)=\1_{[0,\infty)}(x).
\]
若 $x\neq0$，则对充分大的 $n$ 有 $F_n(x)=F(x)$；但
$F_n(0)=0$，而 $F(0)=1$。失败恰好发生在极限分布函数的跳点。
\end{example}

\begin{example}[质量逃逸]\label{ex:7-3-2-2}
令 $\mu_n=\delta_n$。若 $f\in C_c(\R)$ 且 $\supp f\subset[-R,R]$，则 $n>R$ 时
\[
 \int f\,\dd\mu_n=f(n)=0.
\]
所以 $\delta_n$ 模糊收敛到零测度。然而它不可能弱收敛到零测度，因为常数函数 $1$ 属于
$C_b(\R)$，而
\[
 \int1\,\dd\delta_n=1\not\longrightarrow0.
\]
事实上其分布函数在每个固定实数处都趋于零，候选极限的右端质量仍为零，不是概率分布函数。
\end{example}

\begin{example}[经验分布函数]\label{ex:7-3-2-3}
给定样本 $X_1,\ldots,X_n$，经验测度与经验分布函数分别是
\[
 \widehat\mu_n=\frac1n\sum_{k=1}^n\delta_{X_k},
 \qquad
 F_n(x)=\widehat\mu_n(( -\infty,x])
       =\frac1n\sum_{k=1}^n\1_{\{X_k\leq x\}}.
\]
若 $X_k$ 独立同分布且公共分布函数为 $F$，则对每个固定 $x$，随机变量
$\1_{\{X_k\leq x\}}$ 是均值 $F(x)$ 的 Bernoulli 变量，强大数律给
$F_n(x)\to F(x)$ a.s.。这里的“每个固定 $x$”不能直接换成“同时对所有 $x$”；后者是
Glivenko--Cantelli 定理所提供的更强一致结论。
\end{example}

\begin{definition}[分布函数]\label{def:7-3-2-1}
$\R$ 上有限正 Borel 测度 $\mu$ 的\emph{分布函数}是
\[
 F_\mu(x)=\mu(( -\infty,x]).
\]
\end{definition}

\begin{proposition}[分布函数的刻画]\label{proposition:7-3-2-distribution-functions}
有限正 Borel 测度的分布函数递增且右连续，并且
\[
 F_\mu(x-):=\lim_{y\uparrow x}F_\mu(y)=\mu(( -\infty,x)),
 \qquad
 F_\mu(x)-F_\mu(x-)=\mu(\{x\}).
\]
若 $\mu$ 是概率测度，则
\[
 \lim_{x\to-\infty}F_\mu(x)=0,
 \qquad
 \lim_{x\to+\infty}F_\mu(x)=1.
\]
反之，具有上述两个端点极限的右连续非减函数 $F$ 是唯一一个概率测度的分布函数。
\end{proposition}

\begin{proof}
单调性直接来自集合包含关系。若 $y_n\downarrow x$，则
$(-\infty,y_n]\downarrow(-\infty,x]$；有限测度从上连续，故 $F_\mu(y_n)\to F_\mu(x)$。
若 $y_n\uparrow x$，则
$(-\infty,y_n]\uparrow(-\infty,x)$；测度从下连续给出
$F_\mu(y_n)\to\mu((-\infty,x))$，继而由
$(-\infty,x]=(-\infty,x)\mathbin{\dot\cup}\{x\}$ 得到跳跃公式。令半直线
分别向空集与全空间单调收敛，便得概率情形的两个端点极限。

反过来，对这样的 $F$ 应用 Lebesgue--Stieltjes 构造
（定理~\ref{theorem:2-3-3-2}），得到唯一测度 $\mu_F$，满足
$\mu_F((a,b])=F(b)-F(a)$ 以及 $\mu_F(( -\infty,b])=F(b)$。再令 $b\to+\infty$ 可知
$\mu_F(\R)=1$，所以 $F$ 正是概率测度 $\mu_F$ 的分布函数。
\end{proof}

\begin{definition}[Helly 收敛]\label{def:7-3-2-2}
设 $F_n,F$ 都递增。若
\[
 F_n(x)\longrightarrow F(x)
\]
对 $F$ 的每个连续点 $x$ 成立，则称 $F_n$ 按连续点 \emph{Helly 收敛}于 $F$。
当这些函数编码有限测度时，还要另行核对 $\pm\infty$ 处的极限，以判断极限是否保留总质量。
\end{definition}

\begin{definition}[模糊收敛与弱收敛]\label{def:7-3-2-3}
设 $\mu_n,\mu$ 是 $\R$ 上局部有限正 Radon 测度。若
\[
 \int f\,\dd\mu_n\longrightarrow\int f\,\dd\mu
 \qquad(f\in C_c(\R)),
\]
则称 $\mu_n$ \emph{模糊收敛}（vague convergence）于 $\mu$。若 $\mu_n,\mu$ 是概率，并且同一收敛
对每个 $f\in C_b(\R)$ 成立，则称 $\mu_n$ \emph{弱收敛}（weak convergence）于 $\mu$，记作
$\mu_n\Rightarrow\mu$。从 $C_c$ 升到 $C_b$ 必须控制无穷远尾部；紧性正是这一统一控制。
\end{definition}

\begin{theorem}[Helly 选择定理]\label{theorem:7-3-2-1}
设 $(F_n)$ 是 $\R$ 上的递增函数序列，并存在有限常数 $A<B$ 使
$A\leq F_n(x)\leq B$ 对所有 $n,x$ 成立。则存在子列 $(F_{n_j})$ 与一个取值于 $[A,B]$ 的递增
右连续函数 $F$，使
\[
 F_{n_j}(x)\longrightarrow F(x)
\]
对 $F$ 的每个连续点成立。
\end{theorem}

\begin{proof}
枚举有理数 $\Q=\{q_1,q_2,\ldots\}$。有界数列 $(F_n(q_1))_n$ 有收敛子列；在该子列中再使
$F_n(q_2)$ 收敛，如此递推。取第 $j$ 层子列的第 $j$ 项所成对角子列，仍记为 $(F_{n_j})$，则
\[
 \phi(q)=\lim_{j\to\infty}F_{n_j}(q)
\]
对每个 $q\in\Q$ 存在。由于每个 $F_n$ 递增，$\phi$ 在 $\Q$ 上递增。

定义右连续规范
\begin{equation}\label{eq:7-3-helly-limit-definition}
 F(x)=\inf\{\phi(q):q\in\Q,\ q>x\}.
\end{equation}
$F$ 递增且仍取值于 $[A,B]$。若 $x_m\downarrow x$，则 $F(x_m)\geq F(x)$。给定
$\varepsilon>0$，由下确界定义可取有理数 $q>x$ 使
$\phi(q)<F(x)+\varepsilon$；当 $m$ 足够大时 $x_m<q$，从而
$F(x_m)\leq\phi(q)<F(x)+\varepsilon$。所以 $F(x_m)\downarrow F(x)$，即 $F$ 右连续。

对有理数 $a<x<b$，单调性给
\[
 F_{n_j}(a)\leq F_{n_j}(x)\leq F_{n_j}(b).
\]
令 $j\to\infty$，再分别对 $a\uparrow x$、$b\downarrow x$ 取上确界和下确界，得到
\begin{equation}\label{eq:7-3-helly-squeeze}
 F(x-)\leq\liminf_jF_{n_j}(x)
 \leq\limsup_jF_{n_j}(x)\leq F(x).
\end{equation}
这里
$F(x-)=\sup_{a<x,\,a\in\Q}\phi(a)$：一边由 $\phi$ 的单调性得到，另一边则在任意 $y<x$ 与
$x$ 之间插入一个有理数得到。若 $x$ 是 $F$ 的连续点，则 $F(x-)=F(x)$，
式~\eqref{eq:7-3-helly-squeeze} 迫使 $F_{n_j}(x)\to F(x)$。

最后说明这样的连续点足够多。对正整数 $m,r$，在 $[-m,m]$ 中跳跃大于 $1/r$ 的点只能有有限个：
若依次取有限多个这样的点，所有跳跃之和不超过 $B-A$，故点数至多 $r(B-A)+1$。任一不连续点都有
正跳跃，因而属于这些有限集的可数并。故递增函数的不连续点至多可数，证明完成。
\end{proof}

\begin{theorem}[分布函数判据]\label{theorem:7-3-2-2}
对 $\R$ 上概率测度 $\mu_n,\mu$，下列两件事等价：
\begin{enumerate}
  \item $\mu_n\Rightarrow\mu$；
  \item 对 $F_\mu$ 的每个连续点 $x$，有
  $F_{\mu_n}(x)\to F_\mu(x)$。
\end{enumerate}
\end{theorem}

\begin{proof}
先设 $\mu_n\Rightarrow\mu$。若 $x$ 是 $F_\mu$ 的连续点，则
$\mu(\{x\})=0$，所以集合 $(-\infty,x]$ 的边界对 $\mu$ 为零。Portmanteau 定理
\ref{theorem:5-3-3-1} 的连续集结论给
\[
 F_{\mu_n}(x)=\mu_n(( -\infty,x])\longrightarrow
 \mu(( -\infty,x])=F_\mu(x).
\]

反过来，设连续点收敛成立。先证明 $(\mu_n)$ 具有紧性。给定
$0<\varepsilon<1/4$，可选取
$F_\mu$ 的连续点 $a<b$，使
\[
 F_\mu(a)<\varepsilon,
 \qquad F_\mu(b)>1-\varepsilon;
\]
这是因为不连续点至多可数，而分布函数在两端趋于 $0,1$。对充分大的 $n$，
$F_{\mu_n}(a)<2\varepsilon$ 且 $F_{\mu_n}(b)>1-2\varepsilon$，故
\[
 \mu_n([a,b])\geq\mu_n((a,b])
 =F_{\mu_n}(b)-F_{\mu_n}(a)>1-4\varepsilon.
\]
有限多个尚未包括的 $\mu_n$ 各自可由一个紧区间捕获大于
$1-4\varepsilon$ 的质量；把这些区间与 $[a,b]$ 取有限并，便得到整列的统一紧区间。
由于 $4\varepsilon$ 可任意小，这正是 $(\mu_n)$ 的紧性。

 现取 $f\in C_b(\R)$。由刚证的紧性以及 $\mu$ 本身的紧性，可对任意
 $\eta>0$ 选 $R$，使
 \[
 \sup_n\mu_n([-R,R]^c)<\eta,
 \qquad \mu([-R,R]^c)<\eta.
 \]
 取 $\chi\in C_c(\R)$，满足 $0\leq\chi\leq1$ 且在 $[-R,R]$ 上等于一，并记
 $g=f\chi$。先直接由分布函数的区间增量证明
 $\int g\,\dd\mu_n\to\int g\,\dd\mu$。选取 $F_\mu$ 的连续点 $a'<b'$，使
 $\supp g\subset(a',b')$。给定
 $\delta>0$，利用 $g$ 在 $[a',b']$ 上的一致连续性，并利用 $F_\mu$ 的连续点在每个非空开区间中
 稠密，可取分割
 \[
  a'=x_0<x_1<\cdots<x_m=b'
 \]
 使所有 $x_j$ 都是 $F_\mu$ 的连续点，且 $g$ 在每个
 $[x_{j-1},x_j]$ 上的振幅小于 $\delta$。在每段选一点 $\xi_j$ 并令
 \[
  s=\sum_{j=1}^m g(\xi_j)\1_{(x_{j-1},x_j]}.
 \]
 因为 $\mu_n,\mu$ 都是概率测度，
 \[
  \left|\int(g-s)\,\dd\mu_n\right|\leq\delta,
  \qquad
  \left|\int(g-s)\,\dd\mu\right|\leq\delta.
 \]
 另一方面，
 \[
  \int s\,\dd\mu_n
  =\sum_{j=1}^m g(\xi_j)
    \bigl(F_{\mu_n}(x_j)-F_{\mu_n}(x_{j-1})\bigr)
  \longrightarrow\int s\,\dd\mu.
 \]
 因而
 \[
  \limsup_{n\to\infty}
  \left|\int g\,\dd\mu_n-\int g\,\dd\mu\right|\leq2\delta.
 \]
 令 $\delta\downarrow0$，得到所需的紧支测试收敛。另一方面，
 \begin{align*}
 \left|\int f(1-\chi)\,\dd\mu_n\right|
 &\leq\|f\|_\infty\mu_n([-R,R]^c),\\
 \left|\int f(1-\chi)\,\dd\mu\right|
 &\leq\|f\|_\infty\mu([-R,R]^c).
\end{align*}
先令 $n\to\infty$，再令 $\eta\downarrow0$，得到
$\int f\,\dd\mu_n\to\int f\,\dd\mu$。这正是弱收敛。
\end{proof}

\begin{proposition}[Stieltjes 积分收敛]\label{proposition:7-3-2-3}
设 $\mu_n,\mu$ 是 $\R$ 上有限正测度，分布函数分别为 $F_n,F$。假设
$F_n(x)\to F(x)$ 对 $F$ 的每个连续点成立。若 $f\in C_c(\R)$，并且在某个包含
$\supp f$ 的有界开区间 $(a,b)$ 上有
\[
 \sup_n\mu_n((a,b])<\infty,
\]
则
\[
 \int_\R f\,\dd F_n=\int_\R f\,\dd\mu_n
 \longrightarrow
 \int_\R f\,\dd\mu=\int_\R f\,\dd F.
\]
\end{proposition}

\begin{proof}
可把 $a,b$ 略微向区间内部移动，使它们都是 $F$ 的连续点，仍有
$\supp f\subset(a,b)$。记
$M=\sup_n\mu_n((a,b])$。给定 $\varepsilon>0$，由 $f$ 在 $[a,b]$ 上一致连续，可选取分割
\[
 a=x_0<x_1<\cdots<x_m=b,
\]
使每个 $x_j$ 都是 $F$ 的连续点，并且 $f$ 在每个 $[x_{j-1},x_j]$ 上的振幅小于
$\varepsilon$。这可以做到，因为 $F$ 的不连续点至多可数，连续点在每个区间中稠密。

在每段选 $\xi_j\in(x_{j-1},x_j)$，置
\[
 s(x)=\sum_{j=1}^mf(\xi_j)\1_{(x_{j-1},x_j]}(x).
\]
由于 $f$ 在 $(a,b]$ 外为零且 $f(a)=0$，
\begin{align*}
 \left|\int f\,\dd\mu_n-\int s\,\dd\mu_n\right|
 &\leq\varepsilon\mu_n((a,b])\leq\varepsilon M,\\
 \left|\int f\,\dd\mu-\int s\,\dd\mu\right|
 &\leq\varepsilon\mu((a,b]).
\end{align*}
而阶梯积分是有限和
\[
 \int s\,\dd\mu_n
 =\sum_{j=1}^mf(\xi_j)
   \bigl(F_n(x_j)-F_n(x_{j-1})\bigr),
\]
每个端点都是 $F$ 的连续点，所以此有限和趋于
$\int s\,\dd\mu$。于是
\[
 \limsup_{n\to\infty}
 \left|\int f\,\dd\mu_n-\int f\,\dd\mu\right|
 \leq\varepsilon\bigl(M+\mu((a,b])\bigr).
\]
令 $\varepsilon\downarrow0$ 即得结论。
\end{proof}

\begin{corollary}[有限正测度的分布函数判据]\label{corollary:7-3-2-finite-measures}
设 $\mu_n,\mu$ 是 $\R$ 上有限正 Borel 测度，且
$F_{\mu_n}(x)\to F_\mu(x)$ 对 $F_\mu$ 的每个连续点成立。则
\[
 \int f\dd\mu_n\longrightarrow\int f\dd\mu
 \qquad(f\in C_c(\R)).
\]
若还有 $\mu_n(\R)\to\mu(\R)$，则同一结论对每个 $f\in C_b(\R)$ 成立。
\end{corollary}

\begin{proof}
对 $f\in C_c(\R)$，可选 $F_\mu$ 的连续点 $a<b$ 使 $\supp f\subset(a,b)$。于是
\[
 \mu_n((a,b])=F_{\mu_n}(b)-F_{\mu_n}(a)
 \longrightarrow F_\mu(b)-F_\mu(a),
\]
故这些区间质量一致有界；命题~\ref{proposition:7-3-2-3} 给出
$\int f\dd\mu_n\to\int f\dd\mu$，不需要总质量收敛。

再设 $M_n=\mu_n(\R)\to M=\mu(\R)$。若 $M=0$，则 $\mu=0$，并且对每个
$f\in C_b(\R)$，
\[
 \left|\int f\dd\mu_n\right|\leq\norm f_\infty M_n\longrightarrow0.
\]
若 $M>0$，则 $M_n>0$ 最终成立。对这些 $n$ 令
$\bar\mu_n=M_n^{-1}\mu_n$、$\bar\mu=M^{-1}\mu$。它们是概率测度，而且在
$F_\mu$ 的每个连续点（也就是 $F_{\bar\mu}$ 的连续点）都有
\[
 F_{\bar\mu_n}(x)=\frac{F_{\mu_n}(x)}{M_n}\longrightarrow
 \frac{F_\mu(x)}M=F_{\bar\mu}(x).
\]
定理~\ref{theorem:7-3-2-2} 因而给出 $\bar\mu_n\Rightarrow\bar\mu$。对任意
$f\in C_b(\R)$，乘回总质量便得
\[
 \int f\dd\mu_n=M_n\int f\dd\bar\mu_n
 \longrightarrow M\int f\dd\bar\mu=\int f\dd\mu.
\]
这也给出了命题~\ref{proposition:2-3-3-4} 在有限正测度上的分布函数版本。
\end{proof}

\begin{corollary}[Helly--Prokhorov 对照]\label{corollary:7-3-2-4}
每个具有紧性的 $\R$ 上概率测度序列都有弱收敛子列。具体地，Helly 选择从其分布函数中提取连续点
收敛子列，紧性保证所得极限函数在 $-\infty,+\infty$ 的极限分别为 $0,1$。
\end{corollary}

\begin{proof}
设 $F_n$ 是 $\mu_n$ 的分布函数。它们递增且取值于 $[0,1]$，故定理
\ref{theorem:7-3-2-1} 给出子列 $F_{n_j}$ 与递增右连续函数 $F$，使子列在 $F$ 的连续点收敛。

给定 $\varepsilon>0$，紧性提供 $R>0$，使
$\mu_n([-R,R])\geq1-\varepsilon$ 对所有 $n$ 成立。若 $x<-R$ 是 $F$ 的连续点，则
\[
 F(x)=\lim_jF_{n_j}(x)\leq\varepsilon;
\]
若 $x>R$ 是连续点，则 $F(x)\geq1-\varepsilon$。连续点在实线两端都可任意选取，故
\[
 \lim_{x\to-\infty}F(x)\leq\varepsilon,
 \qquad
 \lim_{x\to+\infty}F(x)\geq1-\varepsilon.
\]
让 $\varepsilon\downarrow0$，再用 $0\leq F\leq1$，得到两端极限 $0,1$。因此由 $F$ 的
Lebesgue--Stieltjes 构造得到一个概率测度 $\mu$，其分布函数为 $F$。分布函数判据
\ref{theorem:7-3-2-2} 给 $\mu_{n_j}\Rightarrow\mu$。
\end{proof}

\begin{figure}[htbp]
\centering
\begin{tikzpicture}[x=1cm,y=1cm]
  \begin{scope}
    \draw[book line,-{Stealth[length=2mm]}] (-2.7,0)--(2.8,0) node[right]{$x$};
    \draw[book line,-{Stealth[length=2mm]}] (0,-.15)--(0,2.35) node[above]{$F$};
    \draw[BookAccent,very thick] (-2.5,0)--(0,0)--(0,2)--(2.5,2);
    \draw[BookWarm,semithick,dashed] (-2.5,0)--(.7,0)--(.7,2)--(2.5,2);
    \draw[BookWarning,semithick,dash pattern=on 1.2pt off 1.2pt]
      (-2.5,0)--(1.3,0)--(1.3,2)--(2.5,2);
    \fill[BookWarning] (0,0) circle (1.8pt);
    \fill[BookAccent] (0,2) circle (1.8pt);
    \node[book label,align=center] at (0,-.65) {跳点 $x=0$：\\不要求点态收敛};
    \node[book label] at (-1.45,2.5) {移动 Dirac};
  \end{scope}
  \begin{scope}[xshift=8cm]
    \draw[book line,-{Stealth[length=2mm]}] (-2.8,0)--(3,0) node[right]{$x$};
    \fill[BookAccent!10] (-1.4,-.35) rectangle (1.4,.35);
    \draw[BookAccent,semithick] (-1.4,-.35) rectangle (1.4,.35);
    \node[book label] at (0,-.72) {固定紧区间};
    \foreach \x/\lab in {1.8/$\delta_n$,2.25/$\delta_{n+1}$,2.7/$\delta_{n+2}$}{
      \fill[BookWarning] (\x,0) circle (2.2pt);
      \node[book label,rotate=55] at (\x,.55) {\lab};
    }
    \draw[book dashed arrow] (1.65,1.15)--(2.8,1.15)
      node[midway,above,book label]{质量逃向 $+\infty$};
    \node[book label] at (0,2.15) {$C_c$ 看不见图外质量};
  \end{scope}
\end{tikzpicture}
\caption{左：移动原子的分布函数只在极限跳点失配。右：$\delta_n$ 离开每个固定紧区间，因而可
模糊收敛到零，却不能作为概率弱收敛。}
\label{fig:7-3-helly-jump-escape}
\end{figure}

Helly 选择是一般 Prokhorov 定理在实线上的坐标化版本：有理点对角抽取负责选出候选，
紧性负责保证候选仍有全部质量，分布函数判据负责把坐标收敛还原为弱收敛。经验分布、
统计量的分布与 Stieltjes 变换的极限都可沿这条路线处理。一般 Polish 空间没有一条天然的实数坐标，
因而要使用第~5~章的可数有界 Lipschitz 测试族；在更一般的拓扑空间中，序列甚至可能不够，必须回到
第~1~章的网。

\begin{exercise}
证明递增函数在有界区间内跳跃大于 $\varepsilon$ 的点只有有限个，并由此重新证明其不连续点至多
可数。
\end{exercise}
\begin{exercise}
逐点计算 $\delta_{1/n}$ 的分布函数，并分别在 $x<0$、$x=0$、$x>0$ 检验
定理~\ref{theorem:7-3-2-2}。
\end{exercise}
\begin{exercise}
设有限正测度 $\mu_n$ 模糊收敛于 $\mu$，且 $\mu_n(\R)\to\mu(\R)$。用紧支撑截断证明
$\int f\,\dd\mu_n\to\int f\,\dd\mu$ 对每个 $f\in C_b(\R)$ 成立。
\end{exercise}

分布函数以半直线指标观测测度。下面改用复指数、指数衰减核、Cauchy 核与多项式，并逐一说明哪些
测试族总能识别测度，哪些测试族需要额外增长条件。

\subsection{变换决定测度、矩确定性与 Carleman 条件}\label{subsec:7-3-3}
\index{Carleman 条件}\index{矩方法}\index{Laplace 变换}\index{准解析性}

特征函数对每个概率都存在，并且总能决定分布；Laplace 与 Cauchy--Stieltjes 变换只在各自的自然域
定义，却把支集或预解式结构直接编码进去；多项式矩最便于代数计算，但在无界支集上可能无法
识别测度。差别不在于这些变换是否“形式上含有全部系数”，而在于局部数据能否通过增长控制传播到
整个定义域。

\begin{example}[紧支集确定性]\label{ex:7-3-3-1}
设有限正测度 $\mu,\nu$ 都支持在共同紧区间 $K=[-R,R]$，并满足
\[
 \int_Kx^n\,\mu(\dd x)=\int_Kx^n\,\nu(\dd x)
 \qquad(n=0,1,2,\ldots).
\]
线性组合说明两测度对每个多项式积分相同。给定 $f\in C(K)$，由 Stone--Weierstrass 定理可取
多项式 $p_j$ 使 $\|p_j-f\|_{\infty,K}\to0$，于是
\begin{align*}
 \left|\int_K f\,\dd(\mu-\nu)\right|
 &\leq\left|\int_K(f-p_j)\,\dd\mu\right|
      +\left|\int_K(f-p_j)\,\dd\nu\right|\\
 &\leq\|f-p_j\|_{\infty,K}\bigl(\mu(K)+\nu(K)\bigr)\longrightarrow0.
\end{align*}
故两测度对全部连续测试相同，从而 $\mu=\nu$。紧支集同时保证多项式有界并允许一致逼近；这两点在
$\R$ 上都不再自动成立。
\end{example}

\begin{example}[Gaussian Carleman]\label{ex:7-3-3-2}
标准 Gaussian 的奇矩为零，偶矩满足
\[
 m_{2n}=\E X^{2n}=(2n-1)!!=\frac{(2n)!}{2^nn!}.
\]
Stirling 公式给
\[
 m_{2n}^{1/(2n)}\sim\sqrt{\frac{2n}{e}},
 \qquad
 m_{2n}^{-1/(2n)}\sim\sqrt{\frac{e}{2n}}.
\]
因此 $\sum_nm_{2n}^{-1/(2n)}$ 与 $\sum_nn^{-1/2}$ 一样发散。即使不用 Stirling，也有
$m_{2n}=\prod_{j=1}^n(2j-1)\leq(2n)^n$，从而
$m_{2n}^{-1/(2n)}\geq(2n)^{-1/2}$。下文的 Carleman 定理于是说明 Gaussian 分布由全部矩唯一
决定。
\end{example}

\begin{example}[矩收敛但质量问题]\label{ex:7-3-3-3}
固定整数 $r\geq1$，对 $R>1$ 令
\[
 \mu_R=(1-R^{-(r+1)})\delta_0+R^{-(r+1)}\delta_R.
\]
其总质量为一。对 $1\leq k\leq r$，
\[
 \int x^k\,\dd\mu_R=R^{k-r-1}\longrightarrow0
 =\int x^k\,\dd\delta_0,
\]
而
\[
 \int x^{r+1}\,\dd\mu_R=1.
\]
事实上 $\mu_R\Rightarrow\delta_0$，因为对 $f\in C_b(\R)$，
\[
 \left|\int f\,\dd\mu_R-f(0)\right|
 =R^{-(r+1)}|f(R)-f(0)|
 \leq2\|f\|_\infty R^{-(r+1)}\to0.
\]
所以这个例子揭示的是“任意固定的前 $r$ 阶矩不能控制下一阶矩”，不是紧性失败。若改取
$\delta_R$，远端质量不衰减，才会连紧性也失败。两种尾部机制必须分开。
\end{example}

\begin{definition}[Laplace 变换]\label{def:7-3-3-1}
对 $[0,\infty)$ 上有限正 Borel 测度 $\mu$，定义
\[
 L_\mu(s)=\int_{[0,\infty)}e^{-sx}\,\mu(\dd x),\qquad s\geq0.
\]
同一公式在复半平面 $\Re z>0$ 给出全纯函数。若
$\int e^{a x}\,\dd\mu<\infty$ 对某个 $a>0$ 成立，则定义域可越过零延伸到
$\Re z>-a$，并可在零附近逐阶求导。
\end{definition}

\begin{definition}[Carleman 条件]\label{def:7-3-3-2}
设概率测度 $\mu$ 的全部偶矩
$m_{2n}=\int_\R x^{2n}\,\mu(\dd x)$ 有限。若
\begin{equation}\label{eq:7-3-carleman-condition}
 \sum_{n=1}^\infty m_{2n}^{-1/(2n)}=\infty,
\end{equation}
则称 $\mu$ 满足 \emph{Carleman 条件}。若某个 $m_{2n}=0$，相应项按 $+\infty$ 解释；此时
$\mu=\delta_0$。
\end{definition}

\begin{definition}[矩方法的标准假设]\label{def:7-3-3-3}
设 $X_n,X$ 为实随机变量。称它们满足\emph{矩方法的标准假设}，若：
\begin{enumerate}
  \item 对每个整数 $k\geq0$，$\E|X_n|^k$ 与 $\E|X|^k$ 都有限；
  \item 对每个 $k\geq0$，有 $\E X_n^k\to m_k=\E X^k$；
  \item 矩序列 $(m_k)$ 在 $\R$ 上确定 $X$ 的概率律。
\end{enumerate}
这里必须给出所有阶矩，而不是一个固定的有限截断。证明紧性使用二阶矩，传递第 $k$ 阶矩则
使用一个严格更高的偶数阶矩。
\end{definition}

\begin{theorem}[变换唯一性链]\label{theorem:7-3-3-1}
下列三类变换都决定相应测度。
\begin{enumerate}
  \item $\R$ 上两个有限复 Borel 测度若具有相同的 Fourier--Stieltjes 变换
  $\widehat\mu(t)=\int e^{itx}\,\mu(\dd x)$，则两测度相同。
  \item $\R$ 上两个有限实符号测度若其 Cauchy 变换
  $G_\mu(z)=\int(z-x)^{-1}\,\mu(\dd x)$ 在上半平面相同，则两测度相同。对一般有限复测度，给定
  上、下两个半平面的相同数据仍有同一结论。
  \item $[0,\infty)$ 上两个有限实符号测度若其 Laplace 变换对所有 $s\geq0$ 相同，则两测度
  相同。
\end{enumerate}
\end{theorem}

\begin{proof}
每一项都对两测度之差 $\sigma$ 证明：相应变换为零蕴含 $\sigma=0$。

先证 Fourier 情形。令
\[
 \gamma_\varepsilon(x)=\frac1{\sqrt{4\pi\varepsilon}}
 \exp\!\left(-\frac{x^2}{4\varepsilon}\right),
 \qquad \varepsilon>0.
\]
Gaussian 反演公式是
\[
 \gamma_\varepsilon(u)=\frac1{2\pi}
 \int_\R e^{itu}e^{-\varepsilon t^2}\,\dd t.
\]
由于 $|\sigma|(\R)<\infty$ 且 $e^{-\varepsilon t^2}\in L^1(\R)$，Fubini 定理给卷积测度一个连续
密度
\begin{align*}
 (\gamma_\varepsilon*\sigma)(x)
 &=\int_\R\gamma_\varepsilon(x-y)\,\sigma(\dd y)\\
 &=\frac1{2\pi}\int_\R e^{itx}e^{-\varepsilon t^2}
   \widehat\sigma(-t)\,\dd t=0.
\end{align*}
对 $f\in C_c(\R)$ 再用 Fubini，
\[
 \int_\R(f*\gamma_\varepsilon)(y)\,\sigma(\dd y)
 =\int_\R f(x)(\gamma_\varepsilon*\sigma)(x)\,\dd x=0.
\]
Gaussian 是近似恒等核，所以 $f*\gamma_\varepsilon\to f$ 一致；乘以有限全变差并令
$\varepsilon\downarrow0$，得到 $\int f\,\dd\sigma=0$。有限复 Radon 测度由 $C_c$ 测试唯一决定，
故 $\sigma=0$。

再证 Cauchy 情形。记 Poisson 核
\[
 P_v(u)=\frac1\pi\frac{v}{u^2+v^2},\qquad v>0.
\]
若 $\sigma$ 为实符号测度，则
\begin{equation}\label{eq:7-3-cauchy-poisson-boundary}
 -\frac1\pi\Im G_\sigma(u+iv)
 =\int_\R P_v(u-x)\,\sigma(\dd x)=(P_v*\sigma)(u).
\end{equation}
所以上半平面中 $G_\sigma=0$ 蕴含 $P_v*\sigma=0$ 对每个 $v>0$ 成立。Poisson 核总积分为一，
并且对每个 $\delta>0$，
$\int_{|u|>\delta}P_v(u)\,\dd u\to0$。把积分分成 $|u|<\delta$ 与其补集，利用
$f\in C_c\subset C_0$ 的一致连续性，可得 $f*P_v\to f$ 一致。因此
\[
 \int f\,\dd\sigma
 =\lim_{v\downarrow0}\int(f*P_v)\,\dd\sigma
 =\lim_{v\downarrow0}\int f(u)(P_v*\sigma)(u)\,\dd u=0.
\]
故 $\sigma=0$。若 $\sigma$ 为复测度，则不能把虚部移入复测度积分；但两个半平面的数据给出
\[
 G_\sigma(u+iv)-G_\sigma(u-iv)
 =-2\pi i(P_v*\sigma)(u).
\]
若两边的 Cauchy 数据都为零，仍有 $P_v*\sigma=0$，上述 $C_c$ 逼近对复测度同样成立。

最后证 Laplace 情形。令 $T:[0,\infty)\to(0,1]$，$T(x)=e^{-x}$，并把推前符号测度
$\tau=T_*\sigma$ 视为 $[0,1]$ 上在点 $0$ 无质量的测度。对每个整数 $n\geq0$，
\[
 \int_{[0,1]}y^n\,\tau(\dd y)
 =\int_{[0,\infty)}e^{-nx}\,\sigma(\dd x)=L_\sigma(n)=0.
\]
故 $\tau$ 对所有多项式积分为零。多项式在 $C([0,1])$ 中一致稠密，有限全变差使积分关于一致范数
连续，所以 $\tau$ 对所有连续函数积分为零，进而 $\tau=0$。映射 $T$ 是
$[0,\infty)$ 与 $(0,1]$ 间的 Borel 双射，故 $\sigma=0$。三项均证毕。
\end{proof}

Carleman 条件比指数矩条件弱，不能把 Fourier 变换当成零点附近的全纯函数。所需传播原理是一个
纯实变量的准解析引理。下面只证明本应用需要的对数凸特例；矩产生的序列本身就对数凸，不需要把任意
序列替换成某个不保导数上界的“正则化”。

\begin{lemma}[Carleman--Bang 实变量准解析引理]\label{lemma:7-3-3-carleman-bang}
设 $M_0=1$、$M_n>0$，并且
\[
 M_n^2\leq M_{n-1}M_{n+1}\qquad(n\geq1).
\]
设 $F\in C^\infty(\R;\C)$，且存在常数 $C,A>0$ 使
\[
 \|F^{(n)}\|_\infty\leq CA^nM_n\qquad(n\geq0).
\]
若
$\sum_{n\geq1}M_n^{-1/n}=\infty$，并且在某个 $t_0\in\R$ 有
$F^{(n)}(t_0)=0$ 对所有 $n$ 成立，则 $F\equiv0$。
\end{lemma}

\begin{proof}
平移自变量，并以
\[
 \widetilde F(s)=C^{-1}F(t_0+s/A)
\]
替换 $F$，可把 $t_0=0$ 以及 $C=A=1$ 作为假设。以下省略波浪号。令
\[
 \mu_n=\frac{M_n}{M_{n-1}},\qquad n\geq1.
\]
对数凸性等价于 $\mu_n$ 递增。

第一步由题设级数推出所需的比值级数发散。置 $a_n=\mu_n^{-1}$，则因 $M_0=1$，
\[
 M_n^{-1/n}=(a_1a_2\cdots a_n)^{1/n}=:G_n.
\]
对 $a_1,2a_2,\ldots,na_n$ 用算术--几何平均不等式，得到
\[
 G_n\leq\frac{\sum_{j=1}^nj a_j}{n(n!)^{1/n}}.
\]
又由
$\log(n!)=\sum_{j=1}^n\log j\geq\int_1^n\log x\,\dd x\geq n\log n-n$
可得 $(n!)^{1/n}\geq n/e$，所以
\begin{equation}\label{eq:7-3-carleman-amgm}
 G_n\leq\frac{e}{n^2}\sum_{j=1}^nj a_j.
\end{equation}
对任意 $N$ 把有限和换序，并使用
$j\sum_{n=j}^\infty n^{-2}\leq2$，便有
\[
 \sum_{n=1}^NG_n
 \leq e\sum_{j=1}^Nj a_j\sum_{n=j}^N\frac1{n^2}
 \leq2e\sum_{j=1}^Na_j.
\]
因此题设中 $\sum G_n=\infty$ 蕴含
\begin{equation}\label{eq:7-3-carleman-ratio-divergence}
 \sum_{n=1}^\infty\frac1{\mu_n}=\infty.
\end{equation}

第二步定义 Bang 层级函数
\begin{equation}\label{eq:7-3-bang-function}
 B(x)=\sup_{n\geq0}\frac{|F^{(n)}(x)|}{e^nM_n}.
\end{equation}
归一化导数界给 $0\leq B\leq1$。第 $n$ 项统一不超过 $e^{-n}$，故
$B$ 是有限个连续函数之最大值的统一极限：只取 $0\leq n<N$ 所产生的最大值与 $B$ 相差至多
$e^{-N}$。因此 $B$ 连续。所有导数在零点消失又给 $B(0)=0$。

第三步建立传播估计。固定整数 $q\geq1$ 以及 $x,x+h\in\R$。若 $0\leq n<q$，对
$F^{(n)}$ 写到 $q-n-1$ 阶并使用积分余项：
\begin{align*}
 F^{(n)}(x+h)
 &=\sum_{\ell=0}^{q-n-1}
   \frac{F^{(n+\ell)}(x)}{\ell!}h^\ell\\
 &\quad+\frac{h^{q-n}}{(q-n-1)!}
   \int_0^1(1-s)^{q-n-1}F^{(q)}(x+sh)\,\dd s.
\end{align*}
由于 $\mu_j$ 递增，$n+\ell\leq q$ 时
\[
 \frac{M_{n+\ell}}{M_n}\leq\mu_q^\ell,
 \qquad
 \frac{M_q}{M_n}\leq\mu_q^{q-n}.
\]
结合式~\eqref{eq:7-3-bang-function} 与
$\|F^{(q)}\|_\infty\leq M_q$，得到
\begin{align*}
 \frac{|F^{(n)}(x+h)|}{e^nM_n}
 &\leq B(x)\sum_{\ell=0}^{q-n-1}
   \frac{(e\mu_q|h|)^\ell}{\ell!}\\
 &\quad+e^{-q}\frac{(e\mu_q|h|)^{q-n}}{(q-n)!}\\
 &\leq\max\{B(x),e^{-q}\}\exp(e\mu_q|h|).
\end{align*}
若 $n\geq q$，左端不超过 $e^{-n}\leq e^{-q}$。对 $n$ 取上确界，得到 Bang 基本不等式
\begin{equation}\label{eq:7-3-bang-propagation}
 B(x+h)\leq\max\{B(x),e^{-q}\}\exp(e\mu_q|h|).
\end{equation}

最后用层级穿越反证。若某个 $t>0$ 满足 $B(t)>0$，可选整数 $K$ 使
$e^{-K}<B(t)$。对每个 $k\geq K$ 定义
\[
 x_k=\inf\{x\in[0,t]:B(x)\geq e^{-k}\}.
\]
集合非空，因为 $B(t)>e^{-K}\geq e^{-k}$。又因 $B(0)=0<e^{-k}$ 及 $B$ 连续，有
$B(x_k)=e^{-k}$。层级 $e^{-(k+1)}$ 小于 $e^{-k}$，所以
$0\leq x_{k+1}\leq x_k\leq t$。把
式~\eqref{eq:7-3-bang-propagation} 用于
$x=x_{k+1}$、$h=x_k-x_{k+1}$、$q=k+1$，得到
\[
 e^{-k}\leq e^{-(k+1)}
 \exp\!\bigl(e\mu_{k+1}(x_k-x_{k+1})\bigr),
\]
因而
\[
 x_k-x_{k+1}\geq\frac1{e\mu_{k+1}}.
\]
对 $k=K,\ldots,N$ 求和可得
\[
 t\geq x_K-x_{N+1}
 \geq\frac1e\sum_{k=K}^N\frac1{\mu_{k+1}},
\]
这与式~\eqref{eq:7-3-carleman-ratio-divergence} 矛盾。故 $B(t)=0$ 对所有 $t>0$。对函数
$s\mapsto F(-s)$ 重复同一论证，得到负半轴上的结论。于是 $B\equiv0$，而
$|F(x)|/M_0\leq B(x)$，所以 $F\equiv0$。整个论证只用了实变量 Taylor 公式、AM--GM 与连续
层级穿越，没有使用复分析恒等定理。
\end{proof}

\begin{theorem}[Carleman 定理]\label{theorem:7-3-3-2}
若 $\R$ 上概率测度 $\mu$ 满足 Carleman 条件
\eqref{eq:7-3-carleman-condition}，则 $\mu$ 的 Hamburger 矩问题确定：任何与 $\mu$ 具有全部相同
矩的概率测度都等于 $\mu$。
\end{theorem}

\begin{proof}
先证明一个核心稠密性结论。设 $\rho$ 是满足 Carleman 条件的概率测度。若某个
$m_{2n}=\int|x|^{2n}\,\dd\rho$ 为零，则 $\rho$ 集中在 $\{0\}$，即 $\rho=\delta_0$，多项式在
$L^2(\rho)$ 中当然稠密。以下设所有偶矩为正，并置
\[
 M_n=m_{2n}^{1/2},\qquad M_0=1.
\]
Cauchy--Schwarz 不等式给
\[
 m_{2n}^2
 =\left(\int |x|^{n-1}|x|^{n+1}\,\rho(\dd x)\right)^2
 \leq m_{2n-2}m_{2n+2}.
\]
因此
\[
 M_n^2=m_{2n}
 \leq\sqrt{m_{2n-2}m_{2n+2}}=M_{n-1}M_{n+1};
\]
也就是说，所需的序列自身对数凸。Carleman 级数恰为
$\sum M_n^{-1/n}=\infty$。

设 $f\in L^2(\rho)$ 与所有多项式正交。若 $\norm f_{L^2(\rho)}=0$，则 $f=0$ 已经成立；以下设该
范数为正，并定义
\[
 F(t)=\int_\R e^{itx}f(x)\,\rho(\dd x).
\]
对每个 $n$，Cauchy--Schwarz 给
\[
 \int|x|^n|f(x)|\,\rho(\dd x)
 \leq\|f\|_{L^2(\rho)}m_{2n}^{1/2}<\infty.
\]
由 $|(e^{ihx}-1)/h|\leq|x|$ 反复应用支配收敛，可逐阶求导：
\[
 F^{(n)}(t)=i^n\int_\R x^ne^{itx}f(x)\,\rho(\dd x).
\]
于是
\[
 \|F^{(n)}\|_\infty
 \leq\|f\|_{L^2(\rho)}M_n,
 \qquad
 F^{(n)}(0)=i^n\int x^nf\,\dd\rho=0.
\]
引理~\ref{lemma:7-3-3-carleman-bang} 给出 $F\equiv0$。复测度
$\sigma=f\rho$ 有限，因为
$\int|f|\,\dd\rho\leq\|f\|_2$；其 Fourier--Stieltjes 变换正是 $F$。由
定理~\ref{theorem:7-3-3-1} 的第一项，$\sigma=0$，故 $f=0$ $\rho$-a.e.。这证明多项式在
$L^2(\rho)$ 中稠密。

现在设概率 $\nu$ 与 $\mu$ 的全部矩相同，并令 $\rho=(\mu+\nu)/2$。$\rho$ 与 $\mu$ 有同样的
偶矩，所以也满足 Carleman 条件。两测度都绝对连续于 $\rho$；令
\[
 h=\frac{\dd(\mu-\nu)}{\dd\rho}.
\]
由 $|\mu-\nu|\leq\mu+\nu=2\rho$ 可知 $|h|\leq2$ $\rho$-a.e.，故 $h\in L^2(\rho)$。对每个
$n\geq0$，
\[
 \int x^nh\,\dd\rho=\int x^n\,\dd\mu-\int x^n\,\dd\nu=0.
\]
刚证的多项式稠密性迫使 $h=0$，所以 $\mu-\nu=0$。这就证明了矩确定性。
\end{proof}

例~\ref{ex:7-3-3-2} 验证 Gaussian 满足条件；例~\ref{ex:7-3-1-3} 的对数正态 Stieltjes 族则
展示了条件缺失时的真实失败。准解析引理并未预设 $F$ 具有全纯延拓；它仅凭导数上界与一个
发散级数，把一点的无限阶平坦性传播到整条实线。

\begin{proposition}[指数矩蕴含确定性]\label{proposition:7-3-3-3}
若 $\R$ 上概率测度 $\mu$ 满足
\[
 \int_\R e^{a|x|}\,\mu(\dd x)<\infty
\]
对某个 $a>0$ 成立，则 $\mu$ 由全部矩唯一决定。
\end{proposition}

\begin{proof}
设概率 $\nu$ 与 $\mu$ 具有相同矩。由非负级数的单调收敛，
\begin{align*}
 \int_\R\cosh(ax)\,\nu(\dd x)
 &=\sum_{n=0}^\infty\frac{a^{2n}}{(2n)!}
   \int x^{2n}\,\nu(\dd x)\\
 &=\sum_{n=0}^\infty\frac{a^{2n}}{(2n)!}
   \int x^{2n}\,\mu(\dd x)
 =\int_\R\cosh(ax)\,\mu(\dd x)<\infty.
\end{align*}
由于 $e^{a|x|}\leq2\cosh(ax)$，$\nu$ 也有 $a$ 阶双边指数矩。

因此
\[
 M_\mu(z)=\int e^{zx}\,\mu(\dd x),
 \qquad
 M_\nu(z)=\int e^{zx}\,\nu(\dd x)
\]
都在带形域 $|\Re z|<a$ 全纯。对 $|z|<a$，$e^{zx}$ 的幂级数由 $e^{|z||x|}$ 支配，可以逐项
积分，故
\[
 M_\mu(z)=\sum_{n=0}^\infty\frac{z^n}{n!}\int x^n\,\dd\mu
 =\sum_{n=0}^\infty\frac{z^n}{n!}\int x^n\,\dd\nu=M_\nu(z).
\]
全纯恒等定理把相等性扩展到整个带形域；特别地，在 $z=it$ 上两者的特征函数相同。
定理~\ref{theorem:7-3-3-1} 的第一项给 $\mu=\nu$。这里全纯性由指数矩真正保证；Carleman 定理
没有这一保证，因而不能用这段证明代替准解析论证。这与评注
\ref{rem:5-4-1-complex-mgf} 所强调的复参数矩母函数边界一致。
\end{proof}

\begin{theorem}[矩方法的收敛定理]\label{theorem:7-3-3-4}
若实随机变量 $X_n,X$ 满足定义~\ref{def:7-3-3-3} 中的矩方法标准假设，则
\[
 X_n\Rightarrow X.
\]
\end{theorem}

\begin{proof}
记 $\mu_n=\Law(X_n)$、$\mu=\Law(X)$。二阶矩收敛给出
\[
 C_2:=\sup_n\int x^2\,\mu_n(\dd x)<\infty.
\]
Markov 不等式于是给
\[
 \sup_n\mu_n([-R,R]^c)
 \leq\frac{C_2}{R^2}\longrightarrow0.
\]
所以测度族 $(\mu_n)$ 具有紧性。

任取一个子列。由 Prokhorov 定理~\ref{theorem:5-3-3-2}，可再取子列
$(\mu_{n_j})$ 弱收敛于某个概率测度 $\nu$。下面逐阶证明 $\nu$ 保留全部矩。固定 $k\geq0$，取偶数
$q>k$。$q$ 阶矩收敛且 $q$ 为偶数，故
\[
 C_q:=\sup_n\int|x|^q\,\mu_n(\dd x)<\infty.
\]
对 $R>0$，
\begin{equation}\label{eq:7-3-method-ui-original}
 \sup_n\int_{|x|>R}|x|^k\,\mu_n(\dd x)
 \leq R^{k-q}C_q\longrightarrow0.
\end{equation}
这正是第 $k$ 阶矩所需的一致可积尾估计。

还必须先给弱极限建立同样的高阶界。对
$g_M(x)=\min\{|x|^q,M\}$，函数 $g_M$ 有界连续，所以
\[
 \int g_M\,\dd\nu
 =\lim_{j\to\infty}\int g_M\,\dd\mu_{n_j}\leq C_q.
\]
令 $M\uparrow\infty$ 并用单调收敛，得到
$\int|x|^q\,\dd\nu\leq C_q$，从而
\begin{equation}\label{eq:7-3-method-ui-limit}
 \int_{|x|>R}|x|^k\,\nu(\dd x)
 \leq R^{k-q}C_q\longrightarrow0.
\end{equation}

取连续截断 $0\leq\chi_R\leq1$，使 $\chi_R=1$ 于 $[-R,R]$，并在某个更大的紧区间外为零。
函数 $x^k\chi_R(x)$ 有界连续，故
\[
 \int x^k\chi_R(x)\,\mu_{n_j}(\dd x)
 \longrightarrow
 \int x^k\chi_R(x)\,\nu(\dd x).
\]
式~\eqref{eq:7-3-method-ui-original} 与
式~\eqref{eq:7-3-method-ui-limit} 控制等式两边去掉截断的误差。先令 $j\to\infty$，再令
$R\to\infty$，便得到
\[
 \int x^k\,\nu(\dd x)
 =\lim_{j\to\infty}\int x^k\,\mu_{n_j}(\dd x)
 =m_k=\int x^k\,\mu(\dd x).
\]
由于 $k$ 任意，$\nu$ 与 $\mu$ 的全部矩相同；标准假设中的矩确定性给出 $\nu=\mu$。

我们已经证明：原序列的每个子列都有进一步子列弱收敛到 $\mu$。若 $\mu_n$ 本身不弱收敛到
$\mu$，则由弱拓扑的有界 Lipschitz 度量化命题~\ref{proposition:5-3-3-weak-metric}，存在
$\varepsilon>0$ 与子列使
$d_{\mathrm{BL}}(\mu_{n_j},\mu)\geq\varepsilon$；但该子列又有进一步子列弱收敛于 $\mu$，与
$d_{\mathrm{BL}}\to0$ 矛盾。因此 $\mu_n\Rightarrow\mu$，即 $X_n\Rightarrow X$。
\end{proof}

\begin{figure}[htbp]
\centering
\begin{tikzpicture}[node distance=5.5mm and 15mm]
  \node[book strong node,minimum width=3.2cm] (h1) {测试族};
  \node[book strong node,right=of h1,minimum width=3.7cm] (h2) {自然定义域／代价};
  \node[book strong node,right=of h2,minimum width=3.7cm] (h3) {识别结论};

  \node[book node,below=of h1,minimum width=3.2cm] (cb) {$C_b$};
  \node[book node,right=of cb,minimum width=3.7cm] (cbd) {概率上总可测试};
  \node[book warm node,right=of cbd,minimum width=3.7cm] (cbr) {定义弱收敛并决定测度};

  \node[book node,below=of cb,minimum width=3.2cm] (cf) {$e^{itx}$};
  \node[book node,right=of cf,minimum width=3.7cm] (cfd) {对有限测度总存在};
  \node[book warm node,right=of cfd,minimum width=3.7cm] (cfr) {无条件唯一};

  \node[book node,below=of cf,minimum width=3.2cm] (cs) {$(z-x)^{-1}$};
  \node[book node,right=of cs,minimum width=3.7cm] (csd) {$z$ 离开实支集};
  \node[book warm node,right=of csd,minimum width=3.7cm] (csr) {Poisson 边界唯一};

  \node[book node,below=of cs,minimum width=3.2cm] (la) {$e^{-sx}$};
  \node[book node,right=of la,minimum width=3.7cm] (lad) {$[0,\infty)$，$s\geq0$};
  \node[book warm node,right=of lad,minimum width=3.7cm] (lar) {Stone--Weierstrass 唯一};

  \node[book node,below=of la,minimum width=3.2cm] (mo) {$1,x,x^2,\ldots$};
  \node[book warning node,right=of mo,minimum width=3.7cm] (mod) {矩可能不存在或增长过快};
  \node[book warning node,right=of mod,minimum width=3.7cm] (mor) {紧支集／Carleman 时唯一};

  \foreach \a/\b in {cb/cbd,cbd/cbr,cf/cfd,cfd/cfr,cs/csd,csd/csr,la/lad,lad/lar,mo/mod}{
    \draw[book arrow] (\a)--(\b);
  }
  \draw[book dashed arrow,draw=BookWarning] (mod)--node[above=1mm,book label,font=\scriptsize]{附加条件}(mor);
\end{tikzpicture}
\caption{不同测试族的存在域与识别能力。复指数、Cauchy 核和半轴 Laplace 核各有直接唯一性机制；
多项式矩在无界支集上需要紧支集、指数矩或 Carleman 一类增长条件。}
\label{fig:7-3-transform-determining-families}
\end{figure}

这张比较表给出选工具的原则。组合概率中的极限定理常能逐阶计算矩，适合使用
定理~\ref{theorem:7-3-3-4}；谱与预解式问题通常直接产生 Cauchy--Stieltjes 变换；卷积与独立
和则优先使用总存在的特征函数。无论选择哪一族测试，证明都要分成“先由紧性产生候选极限，
再由决定类识别候选”两步。弱收敛本身不能直接测试无界的 $x^k$，而全部矩存在也不能代替矩确定性。

\begin{remark}[三条不能省略的边界]
\begin{enumerate}
  \item 矩母函数可能除零点外处处为无穷，特征函数却对每个概率、每个实参数都存在；
  \item 弱收敛只控制有界连续测试。要传递第 $k$ 阶矩，必须另有更高阶矩给出一致可积尾界；
  \item 若极限矩问题不确定，可在例~\ref{ex:7-3-1-3} 的两个同矩分布之间交替。此时每一阶矩都
  恒定，分布却不收敛。
\end{enumerate}
\end{remark}

\begin{exercise}
设 $\mu,\nu$ 支持在同一紧集 $K\subset\R$ 且矩相同。把
例~\ref{ex:7-3-3-1} 的证明改写为对符号测度 $\mu-\nu$ 的全变差估计。
\end{exercise}
\begin{exercise}
不用完整 Stirling 公式，只用
$m_{2n}\leq(2n)^n$ 验证标准 Gaussian 的 Carleman 条件；再计算方差为 $\sigma^2$、均值为
$m$ 时条件仍成立。
\end{exercise}
\begin{exercise}
在定理~\ref{theorem:7-3-3-4} 中固定 $k$，写出取偶数 $q>k$ 后
$\{|X_n|^k\}$ 一致可积的完整估计，并说明只知道 $k$ 阶矩有界为什么不够。
\end{exercise}

本节从矩阵正性出发构造并识别测度，又用 Helly 与紧性控制测度列的极限。下一节将把
$f\mapsto\int f\,\dd\mu$ 反过来看成正线性泛函，证明 Daniell--Stone 与 Riesz--Markov 表示；
那时本节的有限构造会获得一个更统一、但逻辑上更晚的抽象解释。

\section{表示、原子逼近与进入泛函分析}\label{sec:07-04}

\subsection{Daniell--Stone 与 Riesz--Markov：从正泛函到 Radon 测度}\label{subsec:7-4-1}
\index{Daniell--Stone 定理}\index{Riesz--Markov 表示定理}\index{Radon 测度!正泛函表示}

测度通常先在集合上给出，再由简单函数逼近定义积分。Daniell 的路线反过来：先规定一族可观测函数的
“积分”，再追问这些数值是否来自一个可数可加测度。有限可加并不困难；真正把有限信息升级为测度的，
是单调列下降到零时积分也下降到零。Riesz--Markov 定理则把拓扑加入这条路线：紧支连续函数不仅生成
Borel 信息，还能借助局部 Urysohn 引理把紧集和开集夹在同一个测试函数中，因而恢复 Radon 正则性。

\begin{definition}[Daniell 泛函]\label{def:7-4-1-1}
设 $X$ 为集合，$\mathcal L$ 是 $X$ 上含常数的有界实函数向量格，即
$f,g\in\mathcal L$ 蕴含 $af+bg,f\vee g,f\wedge g\in\mathcal L$。正线性泛函
$I:\mathcal L\to\R$ 若满足
\[
 f_n\downarrow0\quad\Longrightarrow\quad I(f_n)\downarrow0,
\]
就称为 \emph{Daniell 泛函}。这里的下降和极限均为逐点意义；这个条件称为 Daniell 连续性。
\end{definition}

向量格条件还蕴含 $|f|,f^+,f^-$ 以及对任意连续分段线性函数 $q:\R\to\R$ 的复合
$q\circ f$（只需在 $f(X)$ 所在的有界区间上分段线性）仍属于 $\mathcal L$。这些截断正是外测度
构造中切分覆盖的工具。

\begin{definition}[正 Radon 泛函]\label{def:7-4-1-2}
设 $X$ 为局部紧 Hausdorff 空间。线性泛函 $\Lambda:C_c(X;\R)\to\R$ 若满足
\[
 f\ge0\quad\Longrightarrow\quad\Lambda(f)\ge0,
\]
就称为正 Radon 泛函。它不要求 $C_c(X)$ 带一个全局 Banach 范数；但若 $K$ 紧，取
$\chi\in C_c(X,[0,1])$ 在 $K$ 上等于 $1$，则对 $\supp f\subset K$ 有
\begin{equation}\label{eq:7-4-1-local-bound}
 |\Lambda(f)|\le \Lambda(\chi)\norm f_\infty.
\end{equation}
\end{definition}

确实，$-\norm f_\infty\chi\le f\le\norm f_\infty\chi$，分别对两边应用正性便得到
\eqref{eq:7-4-1-local-bound}。所以正性自动给出每个固定紧支撑层上的一致范数有界性，而没有虚构
$C_c(X)$ 的单一全局上确界范数结构。

\begin{definition}[表示测度]\label{def:7-4-1-3}
若 Borel 测度 $\mu$ 满足
\[
 \Lambda(f)=\int_X f\dd\mu\qquad(f\in C_c(X)),
\]
则称 $\mu$ 表示 $\Lambda$。本节所称 Radon 测度均满足：紧集有限、开集内正则、Borel 集外正则。
\end{definition}

\begin{theorem}[Daniell--Stone]\label{theorem:7-4-1-1}
设 $\mathcal L$ 是集合 $X$ 上含常数的有界实函数向量格，$I:\mathcal L\to\R$ 正线性，
$I(1)<\infty$，并具有 Daniell 连续性。则在 $\sigma(\mathcal L)$ 上存在唯一有限测度 $\mu$，使
\[
 I(f)=\int_X f\dd\mu\qquad(f\in\mathcal L).
\]
\end{theorem}

\begin{proof}
\textbf{第一步：可数函数覆盖成本。}
对任意函数 $h:X\to[0,\infty]$ 定义 Daniell 上积分
\begin{equation}\label{eq:7-4-1-daniell-upper-integral}
 J(h)=\inf\left\{
   \sum_{n=1}^\infty I(g_n):
   g_n\in\mathcal L_+,\quad
   h\le\sum_{n=1}^\infty g_n
 \right\}.
\end{equation}
若 $g\in\mathcal L_+$ 且 $g\le\sum_ng_n$，置
\[
 r_N=\left(g-\sum_{n=1}^Ng_n\right)^+.
\]
则 $r_N\in\mathcal L_+$ 且 $r_N\downarrow0$。又有
$g\le\sum_{n=1}^Ng_n+r_N$，所以正性与 Daniell 连续性给
\begin{equation}\label{eq:7-4-1-cover-dominates}
 I(g)\le\sum_{n=1}^NI(g_n)+I(r_N)
 \longrightarrow\sum_{n=1}^\infty I(g_n).
\end{equation}
对所有覆盖取下确界得到 $I(g)\le J(g)$；以 $g$ 自身作单项覆盖又给 $J(g)\le I(g)$，故
\begin{equation}\label{eq:7-4-1-J-agrees-I}
 J(g)=I(g)\qquad(g\in\mathcal L_+).
\end{equation}

现在定义
\begin{equation}\label{eq:7-4-1-daniell-outer}
 \mu^*(A)=J(\1_A)\qquad(A\subset X).
\end{equation}
空覆盖给 $\mu^*(\varnothing)=0$，单调性来自可用覆盖的包含关系。又因常数函数 $1$ 覆盖每个
$A\subset X$，有 $\mu^*(A)\leq I(1)<\infty$。若
$A\subset\bigcup_kA_k$，为每个 $A_k$ 取成本在
$\mu^*(A_k)+\varepsilon2^{-k}$ 内的可数函数覆盖，再把双重列枚举成一列，便得到
\[
 \mu^*(A)\le\sum_k\mu^*(A_k)+\varepsilon.
\]
令 $\varepsilon\downarrow0$，得到可数次可加性。因此 $\mu^*$ 是外测度；并且由
\eqref{eq:7-4-1-J-agrees-I}，
$\mu^*(X)=J(1)=I(1)<\infty$。这一步说明为什么不能用单个主控函数的成本下确界代替
\eqref{eq:7-4-1-daniell-upper-integral}：可数次可加性来自把各集合的可数覆盖合并，而不是来自一次
有限格运算。

\textbf{第二步：半直线逆像满足 \Caratheodory{} 条件。}
固定 $f\in\mathcal L$、$c\in\R$，令 $E=\{f>c\}$。取任意 $B\subset X$，并取一个有限成本覆盖
\[
 \1_B\le\sum_{n=1}^\infty g_n,\qquad
 S=\sum_{n=1}^\infty I(g_n)<\infty.
\]
这种覆盖总能取到，因为常数 $1$ 覆盖 $X$。令
\[
 u_k=\bigl(k(f-c)^+\bigr)\wedge1,\qquad
 M_n=\norm{g_n}_\infty,\qquad
 a_{n,k}=g_n\wedge(M_n\,u_k),\qquad a_{n,0}=0.
\]
则 $a_{n,k}\uparrow g_n\1_E$ 逐点成立。增量
$d_{n,k}=a_{n,k}-a_{n,k-1}$ 属于 $\mathcal L_+$，而双重族 $(d_{n,k})_{n,k}$ 覆盖 $B\cap E$，因为
\[
 \sum_{n,k}d_{n,k}(x)=\sum_ng_n(x)\quad(x\in E).
\]
记
\[
 D=\sum_n\lim_{k\to\infty}I(a_{n,k}).
\]
由于 $0\le a_{n,k}\le g_n$，有 $0\le D\le S$；非负级数的单调收敛还给
\[
 \mu^*(B\cap E)\le\sum_{n,k}I(d_{n,k})=D.
\]
另一方面，对固定 $k$ 令 $b_{n,k}=g_n-a_{n,k}$。在 $E^c$ 上 $u_k=0$，故 $(b_{n,k})_n$ 覆盖
$B\setminus E$，于是
\[
 \mu^*(B\setminus E)
 \le\sum_nI(b_{n,k})
 =S-\sum_nI(a_{n,k}).
\]
令 $k\to\infty$，非负级数的单调收敛给
\[
 \mu^*(B\cap E)+\mu^*(B\setminus E)\le D+(S-D)=S.
\]
最后对 $B$ 的覆盖成本取下确界，得到 \Caratheodory{} 反向不等式；正向由外测度次可加性成立。
所以每个 $\{f>c\}$ 都可测。这样的半直线逆像生成 $\sigma(\mathcal L)$，故
$\sigma(\mathcal L)$ 包含在 \Caratheodory{} 可测类中。由定理~\ref{theorem:2-1-2-2}，
$\mu=\mu^*|_{\sigma(\mathcal L)}$ 是有限测度。

\textbf{第三步：层级逼近恢复 $I$。}
仍令 $E_t=\{f>t\}$，其中 $f\in\mathcal L$。函数
\[
 u_{t,k}=\bigl(k(f-t)^+\bigr)\wedge1
\]
递增到 $\1_{E_t}$，并约定 $u_{t,0}=0$。由 \eqref{eq:7-4-1-J-agrees-I} 和单调性，
$I(u_{t,k})=J(u_{t,k})\le J(\1_{E_t})=\mu(E_t)$。反向把增量
$u_{t,k}-u_{t,k-1}$ 枚举成一个可数 $\mathcal L_+$ 覆盖；其和在 $E_t$ 上等于 $1$，故
\begin{equation}\label{eq:7-4-1-level-functional}
 \mu(E_t)\le\sum_{k=1}^\infty I(u_{t,k}-u_{t,k-1})
 =\lim_{k\to\infty}I(u_{t,k})\le\mu(E_t).
\end{equation}
为完整写出闭水平集一侧，令
\[
 w_{t,k}=\bigl(k(t-f)^+\bigr)\wedge1.
\]
把刚证的开水平集公式用于函数 $-f$ 和阈值 $-t$，得到
$I(w_{t,k})\uparrow\mu(\{f<t\})$。因此，对
\[
 v_{t,k}=1-w_{t,k}
\]
有
\begin{equation}\label{eq:7-4-1-closed-level-functional}
 I(v_{t,k})
 =I(1)-I(w_{t,k})
 \downarrow\mu(X)-\mu(\{f<t\})
 =\mu(\{f\ge t\}).
\end{equation}

现在设 $0\le f\le M$。给定 $\delta>0$，只把求和写到 $j\delta\le M+\delta$，并置
\[
 s_\delta=\delta\sum_{j\ge1}\1_{\{f>j\delta\}},
 \qquad
 t_\delta=\delta\1_X+
 \delta\sum_{j\ge1}\1_{\{f\ge j\delta\}}.
\]
这是有限和。固定 $x$ 并记 $y=f(x)/\delta$，则
\[
 \#\{j\geq1:j<y\}\leq y
 \leq1+\#\{j\geq1:j<y\},
\]
以及
\[
 y\leq1+\#\{j\geq1:j\leq y\}\leq y+1.
\]
分别乘以 $\delta$，便逐点得到
\[
 0\le f-s_\delta\le\delta,\qquad
 0\le t_\delta-f\le\delta.
\]
对下方每个水平使用 $u_{j\delta,k}$，有限求和后有
$\delta\sum_ju_{j\delta,k}\le s_\delta\le f$。令 $k\to\infty$，由
\eqref{eq:7-4-1-level-functional} 得
\[
 \int_Xs_\delta\dd\mu
 =\delta\sum_j\mu(f>j\delta)\le I(f).
\]
对上方水平使用 $v_{j\delta,k}$，有
$f\le\delta1+\delta\sum_jv_{j\delta,k}$；令 $k\to\infty$ 并用
\eqref{eq:7-4-1-closed-level-functional}，得到
\[
 I(f)\le\delta I(1)+\delta\sum_j\mu(f\ge j\delta)
 =\int_Xt_\delta\dd\mu.
\]
因为 $\mu(X)=I(1)<\infty$，上下阶梯与 $f$ 的一致误差至多 $\delta$。令 $\delta\downarrow0$，夹逼给
$I(f)=\int f\dd\mu$。一般实值 $f$ 写成 $f^+-f^-$ 即得表示。

\textbf{第四步：唯一性。}
若有限测度 $\nu$ 也表示 $I$，则对 $A\in\sigma(\mathcal L)$ 的任意函数覆盖
$\1_A\le\sum_ng_n$，积分的单调收敛性给
\[
 \nu(A)\le\sum_n\int g_n\dd\nu=\sum_nI(g_n).
\]
对覆盖取下确界，得 $\nu(A)\le\mu^*(A)=\mu(A)$。又
$\nu(X)=I(1)=\mu(X)$；把同一不等式用于 $A^c$，便得到反向不等式，故 $\nu(A)=\mu(A)$。
\end{proof}

若函数格不含常数或 $I(1)=\infty$，应先在可由某个 $h\in\mathcal L_+$ 主控的局部层上构造，再核对
这些局部测度的相容性。不能把有限版定理中的 $1$ 直接删去而保留原证明。

下面进入拓扑版本。为避免把正则性藏在一句“标准构造”中，我们从开集上的测试函数上确界开始。

\begin{theorem}[Riesz--Markov 表示]\label{theorem:7-4-1-2}
设 $X$ 为局部紧 Hausdorff 空间。每个正线性泛函
$\Lambda:C_c(X;\R)\to\R$ 都由唯一 Radon 测度 $\mu$ 表示。反之，每个 Radon 测度都在
$C_c(X;\R)$ 上定义正线性泛函。
\end{theorem}

\begin{proof}
\textbf{第一步：开集上的函数上确界及其可加性。}
对开集 $U\subset X$ 定义
\begin{equation}\label{eq:7-4-1-open-mass}
 p(U)=\sup\{\Lambda(f):f\in C_c(X),\ 0\le f\le1,\ \supp f\subset U\}.
\end{equation}
零函数是 $U=\varnothing$ 时的唯一候选，故 $p(\varnothing)=0$；若 $U\subset V$，$U$ 的每个候选也是
$V$ 的候选，故 $p(U)\le p(V)$。若 $U$ 相对紧，取
$\chi\in C_c(X,[0,1])$ 在 $\overline U$ 上等于 $1$，则每个候选 $f$ 满足 $f\le\chi$，从而
$p(U)\le\Lambda(\chi)<\infty$。

若 $U\subset\bigcup_{j\ge1}U_j$ 且 $f$ 是 \eqref{eq:7-4-1-open-mass} 中的候选函数，紧集
$\supp f$ 被有限多个 $U_j$ 覆盖。命题~\ref{proposition:1-3-1-finite-partition-of-unity}
给出相应的有限单位分解 $(\varphi_j)$；于是
\[
 f=\sum_jf\varphi_j,\qquad
 0\le f\varphi_j\le1,\qquad
 \supp(f\varphi_j)\subset U_j.
\]
正性给 $\Lambda(f)\le\sum_jp(U_j)$，取上确界得到
\begin{equation}\label{eq:7-4-1-open-subadd}
 p(U)\le\sum_{j=1}^\infty p(U_j).
\end{equation}
若 $U,V$ 不交，分别取支撑在其中的候选函数 $f,g$，则 $0\le f+g\le1$ 且支撑在 $U\cup V$，故
$p(U\cup V)\ge p(U)+p(V)$。对两两不交的开集 $(U_j)$，归纳应用这项有限超可加性可得
\[
 p\left(\bigcup_{j\geq1}U_j\right)
 \geq p\left(\bigcup_{j=1}^NU_j\right)
 =\sum_{j=1}^Np(U_j)
\]
对每个 $N$ 成立；令 $N\to\infty$，再结合 \eqref{eq:7-4-1-open-subadd}，便得到可数可加性。

还需要开集的内部耗尽公式
\begin{equation}\label{eq:7-4-1-open-inner-p}
 p(U)=\sup\{p(V):V\text{ 开},\ \overline V\text{ 紧且 }\overline V\subset U\}.
\end{equation}
事实上，给定候选 $f$，局部紧性允许取开集 $V$ 使
$\supp f\subset V\Subset U$，于是 $\Lambda(f)\le p(V)$；反向由单调性得到。

\textbf{第二步：开外逼近产生外测度和 Borel 测度。}
对任意 $A\subset X$ 置
\begin{equation}\label{eq:7-4-1-riesz-outer}
 \mu^*(A)=\inf\{p(U):A\subset U,\ U\text{ 开}\}.
\end{equation}
空集和单调性直接成立。为证可数次可加性，若某个
$\mu^*(A_n)=\infty$，则所需不等式右端为无穷，结论直接成立。以下设每个
$\mu^*(A_n)<\infty$。给定 $\varepsilon>0$，由下确界定义可逐个选开集
$A_n\subset U_n$，使
$p(U_n)<\mu^*(A_n)+\varepsilon2^{-n}$；于是
\eqref{eq:7-4-1-open-subadd} 给
\[
 \mu^*\Bigl(\bigcup_nA_n\Bigr)
 \le p\Bigl(\bigcup_nU_n\Bigr)
 \le\sum_np(U_n)
 <\sum_n\mu^*(A_n)+\varepsilon.
\]
故 $\mu^*$ 是外测度，而且对开集 $U$ 有 $\mu^*(U)=p(U)$。

在检验 Borel 可测性之前，先建立紧集公式。若 $K$ 紧，则
\begin{equation}\label{eq:7-4-1-riesz-compact-majorant}
 \mu^*(K)=\inf\{\Lambda(\phi):
 \phi\in C_c(X),\ \phi\ge0,\ \phi|_K\ge1\}.
\end{equation}
记右端为 $q(K)$。给定开集 $U\supset K$，局部 Urysohn 引理给支撑在 $U$ 中且在 $K$ 上等于 $1$
的帽函数 $\phi$，所以 $q(K)\le\Lambda(\phi)\le p(U)$；对 $U$ 取下确界得到
$q(K)\le\mu^*(K)$。反向，若 $\phi$ 是右端候选且 $0<\eta<1$，开集
$U_\eta=\{\phi>1-\eta\}$ 包含 $K$。对每个支撑在 $U_\eta$ 中的 $0\le g\le1$，逐点有
$(1-\eta)g\le\phi$，故
\[
 \mu^*(K)\le p(U_\eta)\le\frac{\Lambda(\phi)}{1-\eta}.
\]
令 $\eta\downarrow0$，再对 $\phi$ 取下确界，得到另一向不等式。特别地，局部帽函数保证每个紧集的
外测度有限。

若紧集 $K,L$ 不交，取任意开集 $W\supset K\cup L$。在开集 $W\setminus L$ 中对紧集 $K$ 使用
局部 Urysohn 引理，取 $\phi\in C_c(X,[0,1])$，使 $\phi=1$ 于 $K$ 且
$\supp\phi\subset W\setminus L$。于是
\[
 U=W\cap\{\phi>2/3\},\qquad V=W\cap\{\phi<1/3\}
\]
是不交开集，分别包含 $K,L$，并且都包含于 $W$。集合函数 $p$ 在不交开集上的
超可加性给
\[
 p(W)\ge p(U)+p(V)\ge\mu^*(K)+\mu^*(L).
\]
对 $W$ 取下确界，再结合外测度次可加性，得到紧集有限可加性
\begin{equation}\label{eq:7-4-1-riesz-compact-additivity}
 \mu^*(K\cup L)=\mu^*(K)+\mu^*(L).
\end{equation}
此外，\eqref{eq:7-4-1-open-inner-p} 可升级成
\begin{equation}\label{eq:7-4-1-riesz-inner-regular}
 p(U)=\sup\{\mu^*(K):K\subset U,\ K\text{ 紧}\}.
\end{equation}
一向来自 $\mu^*(K)\le p(U)$。反向取 $V\Subset U$；对每个开集 $O\supset\overline V$ 有
$p(O)\ge p(V)$，故 $\mu^*(\overline V)\ge p(V)$，再用
\eqref{eq:7-4-1-open-inner-p} 对 $V$ 取上确界。

下面专门证明开集满足 \Caratheodory{} 条件。先由
\eqref{eq:7-4-1-riesz-compact-additivity} 和
\eqref{eq:7-4-1-riesz-inner-regular} 得到：若 $K\subset V$、$K$ 紧而 $V$ 开，则
\begin{equation}\label{eq:7-4-1-riesz-remove-compact}
 p(V)\ge\mu^*(K)+p(V\setminus K).
\end{equation}
事实上，对任意紧集 $L\subset V\setminus K$，
$\mu^*(K)+\mu^*(L)=\mu^*(K\cup L)\le p(V)$；再对 $L$ 取上确界即可。

固定开集 $G$ 和 $A\subset X$，只需处理 $\mu^*(A)<\infty$。取开集 $V\supset A$ 使
$p(V)<\mu^*(A)+\varepsilon$，并令 $W=V\cap G$。由
\eqref{eq:7-4-1-riesz-inner-regular} 取紧集 $K\subset W$，使
$\mu^*(K)>p(W)-\varepsilon$。因为 $A\cap G\subset W$ 且
$A\setminus G\subset V\setminus K$，由 \eqref{eq:7-4-1-riesz-remove-compact}，
\begin{align*}
 \mu^*(A\cap G)+\mu^*(A\setminus G)
 &\le p(W)+p(V\setminus K)\\
 &\le p(W)+p(V)-\mu^*(K)\\
 &<\mu^*(A)+2\varepsilon.
\end{align*}
令 $\varepsilon\downarrow0$ 得 \Caratheodory{} 反向不等式；正向是外测度次可加性。因此所有开集可测，
\Caratheodory{} 可测类包含 $\Borel(X)$。以下把限制在 Borel 集上的测度记为 $\mu$。

\textbf{第三步：Radon 正则性。}
外正则性就是 \eqref{eq:7-4-1-riesz-outer}；紧集有限性来自
\eqref{eq:7-4-1-riesz-compact-majorant}；开集内正则性来自
\eqref{eq:7-4-1-riesz-inner-regular}，因为对开集 $U$ 有 $\mu(U)=\mu^*(U)=p(U)$。所以 $\mu$ 为
Radon 测度。这里所需的拓扑分离仅是紧集与开集之间的局部 Urysohn 引理，并不要求局部紧
Hausdorff 空间本身全局正规。

\textbf{第四步：从层集恢复 $\Lambda$。}
取 $0\le f\in C_c(X)$，令 $M=\norm f_\infty$、$K=\supp f$。若 $M=0$，则 $f=0$，积分表示由线性性成立。
对任意整数 $N\geq1$，令 $\delta=M/N$。对 $j=1,\ldots,N$，开集 $U_j=\{f>j\delta\}$ 的定义给出
$g_j\in C_c(U_j,[0,1])$，使
$\Lambda(g_j)>\mu(U_j)-\varepsilon/N$。这里
$\overline{U_j}\subset K\cap\{f\geq j\delta\}$ 是紧集，故
$\mu(U_j)\leq\mu(K)<\infty$，所以上确界的这一近似选择确实有意义。由于一个点至多落入
$\lfloor f(x)/\delta\rfloor$ 个 $U_j$，有 $\delta\sum_jg_j\le f$，故
\begin{equation}\label{eq:7-4-1-riesz-lower}
 \Lambda(f)\ge\delta\sum_{j=1}^N\mu(f>j\delta)-\delta\varepsilon.
\end{equation}

再令
\[
 h_j=\delta^{-1}\bigl((f-(j-1)\delta)^+\wedge\delta\bigr),
 \qquad f=\delta\sum_{j=1}^Nh_j.
\]
每个 $h_j$ 满足 $0\le h_j\le1$，其支撑包含于紧集
$K_j=K\cap\{f\ge(j-1)\delta\}$。对任意开集 $O\supset K_j$，$h_j$ 是 $p(O)$ 的候选函数，
故 $\Lambda(h_j)\le p(O)$；对 $O$ 取下确界得
$\Lambda(h_j)\le\mu(K_j)$。于是
\begin{equation}\label{eq:7-4-1-riesz-upper}
 \Lambda(f)\le\delta\sum_{j=1}^N\mu(K\cap\{f\ge(j-1)\delta\}).
\end{equation}
令 $G(t)=\mu(f>t)$，并记它在 $[0,M]$ 上的右、左端点和为
\[
 R_N^-:=\delta\sum_{j=1}^NG(j\delta),
 \qquad
 R_N^+:=\delta\sum_{j=0}^{N-1}G(j\delta).
\]
$G$ 有界且递减，故 Riemann 可积，并且
$R_N^-,R_N^+\to\int_0^MG(t)\dd t$。另一方面，令
\[
 U_N:=\delta\sum_{j=1}^N\mu\bigl(K\cap\{f\ge(j-1)\delta\}\bigr).
\]
由于正水平集均包含于 $K$，逐项相减得到
\begin{align*}
 0\le U_N-R_N^+
 &=\delta\bigl(\mu(K)-\mu(f>0)\bigr)
   +\delta\sum_{j=1}^{N-1}\mu(f=j\delta)\\
 &\le2\delta\mu(K).
\end{align*}
最后一个不等式使用了两两不交的水平集
$\{f=j\delta\}$。因此 $U_N\to\int_0^MG(t)\dd t$。把
\eqref{eq:7-4-1-riesz-lower}--\eqref{eq:7-4-1-riesz-upper} 写成
\[
 R_N^- -\delta\varepsilon\le\Lambda(f)\le U_N
\]
并令 $N\to\infty$（即 $\delta=M/N\downarrow0$），层蛋糕公式给
\[
 \Lambda(f)=\int_0^M\mu(f>t)\dd t=\int_Xf\dd\mu.
\]
实值 $f$ 由 $f^+-f^-$ 得到表示。

反过来，若 $\mu$ 为 Radon 测度，$f\in C_c(X)$ 的支撑紧，故
$\int|f|\dd\mu\le\norm f_\infty\mu(\supp f)<\infty$；积分的线性与正性给出正 Radon 泛函。

最后在本定理内完成唯一性。若另一个 Radon 测度 $\nu$ 也表示 $\Lambda$，则对开集 $U$，每个
$0\leq f\leq1$、$f\in C_c(U)$ 都满足 $\Lambda(f)\leq\nu(U)$，故 $p(U)\leq\nu(U)$。反过来，
对任意紧集 $K\subset U$ 取 $f\in C_c(U,[0,1])$ 且 $f=1$ 于 $K$，便有
$\nu(K)\leq\Lambda(f)\leq p(U)$；对 $K$ 使用开集内正则性得 $\nu(U)\leq p(U)$。因此
$\nu(U)=p(U)=\mu(U)$ 对所有开集成立。两测度的 Borel 外正则性随后给
$\nu(E)=\mu(E)$ 对每个 Borel 集 $E$ 成立。存在性、反向与唯一性均已证明；下一命题把刚用到的开集
恢复公式连同支撑刻画单独记录下来。
\end{proof}

\begin{proposition}[正则性与唯一性]\label{proposition:7-4-1-3}
若 $\mu$ 表示正泛函 $\Lambda$，则对每个开集 $U$，
\begin{equation}\label{eq:7-4-1-open-recovery}
 \mu(U)=\sup\{\Lambda(f):f\in C_c(U),\ 0\le f\le1\}.
\end{equation}
因此表示测度唯一。并且
\begin{equation}\label{eq:7-4-1-support-functional}
 x\notin\supp\mu
 \quad\Longleftrightarrow\quad
 \text{存在开邻域 }U\ni x\text{，使 }\Lambda(f)=0
 \text{ 对所有 }f\in C_c(U).
\end{equation}
\end{proposition}

\begin{proof}
若 $f$ 是 \eqref{eq:7-4-1-open-recovery} 中的候选函数，则
$\Lambda(f)=\int f\dd\mu\le\mu(U)$。反过来，给定紧集 $K\subset U$，局部 Urysohn 引理
\ref{theorem:1-3-1-2} 给 $f\in C_c(U,[0,1])$ 且 $f=1$ 于 $K$，所以
$\Lambda(f)\ge\mu(K)$。对 $K$ 取上确界并用开集内正则性，得到等号。

若 $\mu,\nu$ 都表示 $\Lambda$，开集公式给 $\mu(U)=\nu(U)$。对任意 Borel 集 $E$，外正则性给
\[
 \mu(E)=\inf_{U\supset E}\mu(U)
 =\inf_{U\supset E}\nu(U)=\nu(E),
\]
所以表示唯一。最后，若 $x\notin\supp\mu$，按支撑定义有开邻域 $U$ 满足 $\mu(U)=0$，支撑在
$U$ 中的函数积分均为零。反向若右侧成立，尤其所有非负候选函数的值为零；
\eqref{eq:7-4-1-open-recovery} 给 $\mu(U)=0$，故 $x\notin\supp\mu$。
\end{proof}

定理~\ref{theorem:7-3-1-2} 由 \Fejer{} 核和测度紧性证明 Herglotz 表示。现在给出另一种
正泛函证明。它所需的代数工具是 \Fejer{}--Riesz 分解：非负三角多项式可表为一个多项式的模平方。

\begin{lemma}[\Fejer{}--Riesz 分解]\label{lemma:7-4-1-fejer-riesz}
若三角多项式
\[
 q(t)=\sum_{k=-n}^na_ke^{ikt}
\]
在 $\T$ 上取实值且非负，则存在普通多项式 $P(z)=\sum_{k=0}^nb_kz^k$，使
\[
 q(t)=|P(e^{it})|^2.
\]
\end{lemma}

\begin{proof}
若 $q(t)\equiv a_0$，则 $a_0\ge0$；取常数多项式 $P\equiv\sqrt{a_0}$ 即得结论（包括
$a_0=0$）。以下设 $q$ 非常数，并删去两端的零系数，使 $a_n\ne0$；实值性给
$a_{-k}=\overline{a_k}$。令
\[
 Q(z)=z^n\sum_{k=-n}^na_kz^k.
\]
则 $Q$ 是次数 $2n$ 且常数项非零的多项式，并满足
\[
 Q(z)=z^{2n}\overline{Q(1/\overline z)}\qquad(z\ne0).
\]
所以不在单位圆上的零点按 $\alpha\leftrightarrow1/\overline\alpha$ 成对且重数相同。若
$e^{it_0}$ 是单位圆上的零点，则
$q(t)=e^{-int}Q(e^{it})$ 在 $t_0$ 的零点重数相同。设该重数为 $m$。由于 $q$ 是有限个指数函数的
线性组合，其 Taylor 展开在 $t_0$ 邻域成立；存在实数 $c\ne0$，使
\[
 q(t_0+h)=ch^m+o(|h|^m)\qquad(h\to0).
\]
若 $m$ 为奇数，右侧在 $h>0$ 与 $h<0$ 时异号，与 $q\geq0$ 矛盾；故 $m$ 为偶数（并且
$c>0$）。因此每个单位圆零点的重数均为偶数。

从每一对离圆零点中取一个，并从每个圆上零点取其一半重数，所得重数总和为 $n$。令 $P_0$ 以这些
零点为根，并置
$P_0^*(z)=z^n\overline{P_0(1/\overline z)}$；右式按 $P_0$ 的系数展开后在 $z=0$ 也延拓为
次数不超过 $n$ 的多项式。多项式 $Q$ 与
$P_0P_0^*$ 有完全相同的零点及重数，故 $Q=\lambda P_0P_0^*$。在 $|z|=1$ 上，
$P_0^*(z)=z^n\overline{P_0(z)}$，从而
$q(t)=\lambda|P_0(e^{it})|^2$。由 $q\ge0$ 且 $q\not\equiv0$ 得 $\lambda>0$；取
$P=\sqrt\lambda P_0$ 即可。
\end{proof}

\begin{proposition}[Herglotz 表示的正泛函证明]\label{proposition:7-4-1-herglotz-functional}
若复数列 $(c_n)_{n\in\Z}$ 满足 $c_{-n}=\overline{c_n}$，且每个 Toeplitz 矩阵
$T_N=(c_{j-k})_{0\le j,k\le N}$ 都半正定，则存在唯一的有限正 Borel 测度 $\mu$ 于 $\T$，使
\[
 c_n=\int_\T e^{-int}\dd\mu(t)\qquad(n\in\Z).
\]
\end{proposition}

\begin{proof}
在三角多项式代数 $\mathcal P$ 上定义
\[
 L\left(\sum_{k=-N}^Na_ke^{ikt}\right)=\sum_{k=-N}^Na_kc_{-k}.
\]
若 $r(t)=\sum_{k=0}^Nb_ke^{ikt}$，则
\[
 L(|r|^2)=\sum_{j,k=0}^N\overline{b_j}c_{j-k}b_k\ge0.
\]
由引理~\ref{lemma:7-4-1-fejer-riesz}，每个非负三角多项式都是某个 $|r|^2$，所以 $L$ 为正泛函。
关系 $c_{-n}=\overline{c_n}$ 先给出
$L(\overline h)=\overline{L(h)}$（$h\in\mathcal P$）。再把所用的 Cauchy--Schwarz 不等式写明。
给定 $p,q\in\mathcal P$，记
\[
 a=L(|p|^2),\qquad b=L(|q|^2),\qquad c=L(\overline q p).
\]
由 $L(|p+\lambda q|^2)\geq0$，对每个 $\lambda\in\C$ 有
\[
 0\leq a+\lambda\overline c+\overline\lambda c+|\lambda|^2b.
\]
若 $b>0$，取 $\lambda=-c/b$ 得 $|c|^2\leq ab$；若 $b=0$，先取
$\lambda=t\in\R$，再取 $\lambda=it$，并令 $t$ 在 $\R$ 中任意变化，分别迫使
$\Re c=0$ 与 $\Im c=0$。所以零分母情形也有 $c=0$，从而总有
\[
 |L(\overline q p)|^2\leq L(|p|^2)L(|q|^2).
\]
取 $q=1$，再由 $\norm p_\infty^2-|p|^2\ge0$，便有
\[
 |L(p)|^2\le L(1)L(|p|^2)
 \le c_0^2\norm p_\infty^2.
\]
故 $L$ 关于一致范数有界。由推论~\ref{corollary:7-2-1-2}，$\mathcal P$ 在 $C(\T)$ 中稠密，
于是 $L$ 唯一延拓为 $C(\T)$ 上的有界线性泛函 $\widetilde L$。该延拓仍为正：若
$f\in C(\T)$ 且 $f\ge0$，取 $p_m\in\mathcal P$ 一致逼近 $\sqrt f$，则
$|p_m|^2\to f$ 一致且 $L(|p_m|^2)\ge0$。

正性保证 $\widetilde L$ 在实值函数上取实值；又因 $\T$ 紧，
$C_c(\T;\R)=C(\T;\R)$。把 $\widetilde L$ 限制到这个实空间，并应用
定理~\ref{theorem:7-4-1-2}，得到唯一有限正测度 $\mu$。再由复线性性，
$\widetilde L(f)=\int f\dd\mu$ 对复值 $f\in C(\T)$ 也成立。特别地
\[
 \int_\T e^{-int}\dd\mu(t)=L(e^{-int})=c_n.
\]
若另一有限正测度具有相同的 Fourier 矩，则其积分泛函与 $L$ 在 $\mathcal P$ 上相同；三角多项式
稠密性和两个积分泛函的一致范数连续性说明它也表示 $\widetilde L$，故 Riesz--Markov 的唯一性迫使
它等于 $\mu$。
因此 Toeplitz 正性先在三角多项式上产生正泛函，一致范数连续性把它延到 $C(\T)$，
最后由 Riesz--Markov 定理得到表示测度。
\end{proof}

\begin{figure}[htbp]
\centering
\begin{tikzpicture}[node distance=7mm and 36mm]
  \node[book strong node] (fun) {正泛函\\$\Lambda:C_c(X)\to\R$};
  \node[book strong node,below=of fun] (open) {开集上确界\\$p(U)=\sup\Lambda(f)$};
  \node[book strong node,below=of open] (outer) {外测度\\$\mu^*(A)=\inf_{U\supset A}p(U)$};
  \node[book strong node,below=of outer] (radon) {Radon 测度\\内正则、外正则};
  \node[book warm node,below=of radon] (int) {积分表示\\$\Lambda(f)=\int f\dd\mu$};
  \node[book warm node,left=of outer] (ury) {Urysohn 帽函数\\$K\subset U$};
  \draw[book accent arrow] (fun)--(open);
  \draw[book accent arrow] (open)--(outer);
  \draw[book accent arrow] (outer)--(radon);
  \draw[book accent arrow] (radon)--(int);
  \draw[book dashed arrow,bend left=18] (ury.north east) to
    node[above,book label]{紧开夹逼}(open.west);
  \draw[book dashed arrow,bend right=18] (ury.south east) to
    node[below,book label]{层集阶梯和}(int.west);
\end{tikzpicture}
\caption{Riesz--Markov 构造链。开集公式产生外测度；Urysohn 函数同时负责局部有限、内正则与积分识别。}
\label{fig:7-4-1-riesz-route}
\end{figure}

图~\ref{fig:7-4-1-riesz-route} 中两条虚线是拓扑真正进入证明的位置。若只保留外测度箭头而删去
紧开夹逼，最多得到某个 Borel 测度，不能推出表示测度具有所需的 Radon 正则性。

\begin{theorem}[$C_0(X)^*$ 与复测度]\label{theorem:7-4-1-4}
设 $X$ 为局部紧 Hausdorff 空间，标量域为 $\R$ 或 $\C$。每个有界线性泛函
$\Lambda\in C_0(X)^*$ 唯一表示为
\[
 \Lambda(f)=\int_X f\dd\mu\qquad(f\in C_0(X)),
\]
其中 $\mu$ 是有限符号 Radon 测度（实情形）或有限复 Radon 测度（复情形），并且
\begin{equation}\label{eq:7-4-1-dual-variation-norm}
 \norm\Lambda=|\mu|(X).
\end{equation}
反之，每个这样的测度都定义一个有界线性泛函。
\end{theorem}

\begin{proof}
\textbf{第一步：先在泛函层作 Jordan 分解。}
先设标量为实数。对 $f\in C_0(X)_+$ 定义
\begin{equation}\label{eq:7-4-1-functional-positive-part}
 \Lambda^+(f)=\sup\{\Lambda(g):0\le g\le f\}.
\end{equation}
候选集非空，且
$|\Lambda(g)|\le\norm\Lambda\norm f_\infty$，所以上确界有限。正齐次性直接成立。若
$f_1,f_2\ge0$，候选 $g_i\le f_i$ 给
$\Lambda^+(f_1+f_2)\ge\Lambda^+(f_1)+\Lambda^+(f_2)$。反向取
$0\le h\le f_1+f_2$，令
\[
 h_1=h\wedge f_1,\qquad h_2=h-h_1.
\]
则 $0\le h_1\le f_1$，而
$h_2=(h-f_1)^+\le f_2$，故
$\Lambda(h)\le\Lambda^+(f_1)+\Lambda^+(f_2)$。对 $h$ 取上确界，得到可加性。

每个 $u\in C_0(X;\R)$ 都可写成两个正函数之差。若
$u=a-b=c-d$，其中 $a,b,c,d\ge0$，则 $a+d=b+c$；正锥上的可加性给
$\Lambda^+(a)-\Lambda^+(b)=\Lambda^+(c)-\Lambda^+(d)$。因此群完备化公式
\[
 \Lambda^+(a-b)=\Lambda^+(a)-\Lambda^+(b)
\]
良定义且给出实线性延拓；取 $a=u^+,b=u^-$ 就得到通常公式。若 $u\ge0$，
$0\le\Lambda^+(u)\le\norm\Lambda\norm u_\infty$，再对一般 $u$ 使用
$|\Lambda^+(u)|\le\Lambda^+(|u|)$，可知延拓有界且正。又因
\eqref{eq:7-4-1-functional-positive-part} 的候选中可取 $g=f$，
\[
 \Lambda^-(f):=\Lambda^+(f)-\Lambda(f)\ge0\qquad(f\ge0).
\]
所以 $\Lambda=\Lambda^+-\Lambda^-$ 先在泛函层分解为两个正泛函之差，随后可分别对
$\Lambda^+$ 和 $\Lambda^-$ 应用正 Riesz--Markov 定理。

为了把 $C_c(X)$ 上的表示延到 $C_0(X)$，先记录所需的稠密性。给 $u\in C_0(X)$ 和 $\varepsilon>0$，紧集
$K_\varepsilon=\{|u|\ge\varepsilon\}$ 有一个 $\chi\in C_c(X,[0,1])$ 满足
$\chi=1$ 于 $K_\varepsilon$；于是
$u\chi\in C_c(X)$ 且 $\norm{u-u\chi}_\infty<\varepsilon$。这正是
命题~\ref{proposition:1-3-1-4}。把
$\Lambda^\pm$ 限制到 $C_c(X)$，正 Riesz--Markov 定理给 Radon 测度 $\mu^\pm$。它们有限，因为
\[
 \mu^\pm(X)
 =\sup_{0\le f\le1,\ f\in C_c}\Lambda^\pm(f)
 \le\norm{\Lambda^\pm}<\infty.
\]
令 $\mu=\mu^+-\mu^-$。表示先在 $C_c$ 上成立；积分泛函与 $\Lambda$ 都关于上确界范数连续，故由
$C_c$ 稠密性延到全部 $C_0$。这里一般正泛函 $C_c(X)\to\R$ 的表示测度可以有无限总质量；正是
$C_0$ 泛函的全局有界性给出了上式中的 $\mu^\pm(X)<\infty$。

\textbf{第二步：复泛函由两个实泛函合成。}
设 $\Lambda$ 复线性。对实值 $f\in C_0(X;\R)$ 令
\[
 A(f)=\Re\Lambda(f),\qquad B(f)=\Im\Lambda(f).
\]
$A,B$ 是有界实线性泛函。由第一步，存在有限符号 Radon 测度 $\alpha,\beta$ 表示它们。
置 $\mu=\alpha+i\beta$。对实值 $u,v$，复线性给
\begin{align*}
 \Lambda(u+iv)
 &=\Lambda(u)+i\Lambda(v)\\
 &=A(u)-B(v)+i\bigl(B(u)+A(v)\bigr)\\
 &=\int_Xu\dd\alpha-\int_Xv\dd\beta
   +i\left(\int_Xu\dd\beta+\int_Xv\dd\alpha\right)\\
 &=\int_X(u+iv)\dd(\alpha+i\beta).
\end{align*}
故 $\mu$ 表示 $\Lambda$。这个顺序仍是“先分解泛函、再得到测度”。

\textbf{第三步：范数等于全变差。}
测度极分解给 $\dd\mu=h\dd|\mu|$，其中 $|h|=1$ $|\mu|$-a.e.。对
$\norm f_\infty\le1$，
\[
 |\Lambda(f)|\le\int|f|\dd|\mu|\le|\mu|(X),
\]
所以 $\norm\Lambda\le|\mu|(X)$。

反向不能把可测相位 $\overline h$ 直接当成 $C_0$ 测试。由
定理~\ref{theorem:3-3-3-1}，存在 $g_n\in C_c(X;\C)$ 使
$g_n\to\overline h$ 于 $L^1(|\mu|)$。令闭单位圆盘的径向投影为
\[
 \rho(z)=\begin{cases}z,&|z|\le1,\\ z/|z|,&|z|>1.
 \end{cases}
\]
则 $f_n=\rho\circ g_n\in C_c(X)$、$\norm{f_n}_\infty\le1$，并且因为径向投影固定单位圆周，
\[
 |f_n-\overline h|\le2|g_n-\overline h|.
\]
这个点态估计可分两种情况核对：$|g_n|\leq1$ 时 $f_n=g_n$，所需不等式立即成立；$|g_n|>1$ 时，由
$|\overline h|=1$ 与反三角不等式，
\[
 \left|\frac{g_n}{|g_n|}-\overline h\right|
 \leq\left|\frac{g_n}{|g_n|}-g_n\right|+|g_n-\overline h|
 =\bigl||g_n|-1\bigr|+|g_n-\overline h|
 \leq2|g_n-\overline h|.
\]
因此
\[
 \left|\Lambda(f_n)-|\mu|(X)\right|
 =\left|\int_X(f_nh-1)\dd|\mu|\right|
 \le\int_X|f_n-\overline h|\dd|\mu|\longrightarrow0.
\]
取上极限得到 $\norm\Lambda\ge|\mu|(X)$，与前一不等式合并即为
\eqref{eq:7-4-1-dual-variation-norm}。实符号情形可取 $h=\sgn\mu\in\{-1,1\}$，同一逼近完全适用。

若 $\mu,\nu$ 都表示同一泛函，则 $\mu-\nu$ 表示零泛函。把刚证明的范数公式
用于该差测度，得到 $|\mu-\nu|(X)=0$，所以 $\mu=\nu$。

最后，给定有限符号或复 Radon 测度，绝对值估计直接给
$|\int f\dd\mu|\le|\mu|(X)\norm f_\infty$，故它确实定义 $C_0(X)^*$ 中的泛函；把上面的
$L^1(|\mu|)$ 相位逼近原样应用于这个泛函，又给出反向范数不等式，所以其范数仍恰为 $|\mu|(X)$。
\end{proof}

\begin{example}[点值泛函]\label{ex:7-4-1-1}
固定 $x\in X$，令 $\Lambda_x(f)=f(x)$。对 Borel 集 $A$，Dirac 测度满足
\[
 \int_X f\dd\delta_x=f(x),\qquad
 \delta_x(A)=\1_A(x).
\]
故 $\delta_x$ 表示 $\Lambda_x$，且 $\norm{\Lambda_x}=1=|\delta_x|(X)$。更一般地，
$\sum_{j=1}^Nc_jf(x_j)$ 由 $\sum_jc_j\delta_{x_j}$ 表示；若 $c_j\ge0$，它就是正泛函。
\end{example}

\begin{example}[加权积分]\label{ex:7-4-1-2}
在 $X=\R^d$ 上设 $w\ge0$ 局部可积。对 $f\in C_c(\R^d)$，
\[
 \Lambda(f)=\int_{\R^d}f(x)w(x)\dd x
\]
有限且正，由 Radon 测度 $\mu=w\dd x$ 表示。更准确地，令
\[
 \operatorname{ess\,supp}w
 =\R^d\setminus\bigcup\{U:U\text{ 开且 }w=0\text{ a.e. 于 }U\}.
\]
这是一个闭集，并且
\[
 x\notin\supp\mu
 \Longleftrightarrow
 \text{存在开邻域 }U\ni x\text{ 使 }w=0\text{ a.e. 于 }U.
\]
所以 $\supp\mu=\operatorname{ess\,supp}w$。若另有 $w\in L^1$，则该泛函延到 $C_0$，且
$\norm\Lambda=\int|w|\dd x$；仅有局部可积时通常只在 $C_c$ 上定义。
\end{example}

\begin{example}[非正泛函]\label{ex:7-4-1-3}
取 $\phi\in C_c^1(\R)$ 且 $\phi'(0)\ne0$，令
$f_n(x)=n^{-1}\phi(nx)$。则
\[
 \norm{f_n}_\infty=n^{-1}\norm\phi_\infty\longrightarrow0,
 \qquad f_n'(0)=\phi'(0).
\]
因此 $f\mapsto f'(0)$ 在 $C_c^1$ 的
$\norm f_{C^1}=\norm f_\infty+\norm{f'}_\infty$ 拓扑下连续，却不可能连续延到带上确界范数的
$C_0(\R)$。若它由有限符号测度 $\nu$ 表示，就会有
$|f_n'(0)|\le|\nu|(\R)\norm{f_n}_\infty\to0$，与上式矛盾。测试空间的拓扑改变后，表示对象会从
测度升级为一阶分布。
\end{example}

两条表示路线的分工现在很清楚。Daniell--Stone 用单调连续性换取可数可加，适合从随机变量格或初等
积分出发；Riesz--Markov 用局部紧拓扑和紧支连续测试换取正则性，适合弱收敛和对偶空间。

\subsection{原子测度、经验测度与求积：有限对象逼近 Radon 测度}\label{subsec:7-4-2}
\index{原子测度!弱星稠密}\index{经验测度}\index{Tchakaloff 定理}\index{求积公式}

一个弱星基本邻域只询问有限多个积分。因此，为了进入这个邻域，并不需要逐点重建测度；只需把空间
分成有限片，使每个测试函数在每片上的振幅很小，再把整片质量搬到一个代表点。随机抽样把同一结构
变成经验测度，有限维凸性则把“近似匹配”加强为 Tchakaloff 的“精确匹配”。

\begin{example}[Riemann 求积]\label{ex:7-4-2-1}
设 $a=x_0<\cdots<x_N=b$，$\xi_j=(x_{j-1}+x_j)/2$，并令
\[
 \nu_P=\sum_{j=1}^N(x_j-x_{j-1})\delta_{\xi_j}.
\]
若 $|P|=\max_j(x_j-x_{j-1})$，连续函数 $f$ 的模连续性记为
$\omega_f(r)=\sup_{|x-y|\le r}|f(x)-f(y)|$，则
\begin{align*}
 \left|\int_a^bf(x)\dd x-\int f\dd\nu_P\right|
 &\le\sum_{j=1}^N\int_{x_{j-1}}^{x_j}|f(x)-f(\xi_j)|\dd x\\
 &\le(b-a)\omega_f(|P|/2)\longrightarrow0.
\end{align*}
这就是 Lebesgue 测度的有限原子逼近。若每段改用端点平均，得到梯形权重，完全相同的计算给误差
至多 $(b-a)\omega_f(|P|)$；Gauss 节点则不是只压低振幅误差，而是精确匹配一段有限维多项式空间。
\end{example}

\begin{example}[经验测度的测试积分]\label{ex:7-4-2-2}
若 $X_1,X_2,\ldots$ 独立同分布，共同分布为 $\mu$，则对任意有界实可测 $f$，
\[
 \int f\dd L_n=\frac1n\sum_{k=1}^nf(X_k),\qquad
 \E\int f\dd L_n=\int f\dd\mu,
\]
并且在 $f\in L^2(\mu)$ 时
\[
 \Var\!\left(\int f\dd L_n\right)=\frac{\Var(f(X_1))}{n}.
\]
单个 $f$ 的强大数律只给一个测试方向。Polish 假设的作用，是允许选取一个可数的
有界 Lipschitz 决定类，并对相应概率一事件取可数交；定理~\ref{theorem:7-4-2-3} 将完成这一步。
\end{example}

\begin{example}[矩匹配求积]\label{ex:7-4-2-3}
取 $K=[-1,1]$ 上的均匀概率 $\mu=\frac12\dd x$，并只观察
$V=\Span\{1,x,x^2\}$。两点测度
\[
 \nu=\frac12\delta_{-1/\sqrt3}+\frac12\delta_{1/\sqrt3}
\]
满足
\[
 \int1\dd\nu=1=\int1\dd\mu,\qquad
 \int x\dd\nu=0=\int x\dd\mu,\qquad
 \int x^2\dd\nu=\frac13=\int x^2\dd\mu.
\]
它没有逼近全部连续测试，却对指定的三维空间给出零误差。这正是 Gauss 两点求积的最低阶模型。
\end{example}

\begin{definition}[原子测度]\label{def:7-4-2-1}
形如
\[
 \nu=\sum_{j=1}^Nc_j\delta_{x_j}
\]
的符号或复测度称为有限原子测度。若 $c_j\ge0$ 且 $\sum_jc_j=1$，它是有限支持概率测度；对任意
可积测试函数，
$\int f\dd\nu=\sum_jc_jf(x_j)$。
\end{definition}

\begin{definition}[弱星测试邻域]\label{def:7-4-2-2}
把有限符号或复 Radon 测度依定理~\ref{theorem:7-4-1-4} 看作 $C_0(X)^*$ 的元素。
给定有限集 $G\subset C_0(X)$ 与 $\varepsilon>0$，定义
\[
 \mathcal N(\mu;G,\varepsilon)
 =\left\{\nu:
   \left|\int f\dd(\nu-\mu)\right|<\varepsilon
   \text{ 对每个 }f\in G
 \right\}.
\]
这些集合构成 $\mu$ 的弱星基本邻域。
\end{definition}

\begin{definition}[经验测度]\label{def:7-4-2-3}
若 $X_1,X_2,\ldots$ 为 $X$-值随机元，定义
\[
 L_n=\frac1n\sum_{k=1}^n\delta_{X_k}.
\]
它在每条样本路径上都是概率测度，并满足
$\int f\dd L_n=n^{-1}\sum_{k=1}^nf(X_k)$。
\end{definition}

\begin{theorem}[Dirac 组合弱星稠密]\label{theorem:7-4-2-1}
设 $X$ 为局部紧 Hausdorff 空间，$\mu$ 为有限符号或复 Radon 测度。存在有限原子测度网
$(\nu_\alpha)$，使
\[
 \nu_\alpha\weakstarto\mu,
 \qquad |\nu_\alpha|(X)\le|\mu|(X).
\]
网还可保持总质量 $\nu_\alpha(X)=\mu(X)$；若 $\mu\ge0$，可令每个 $\nu_\alpha\ge0$，从而同时保持
正性和总质量。若 $C_0(X)$ 可分，则可以用序列代替网。
\end{theorem}

\begin{proof}
令
\[
 \mathfrak D=\{(G,\varepsilon):G\subset C_0(X)\text{ 有限},\ \varepsilon>0\},
\]
并规定
$(G,\varepsilon)\preceq(H,\eta)$ 当且仅当 $G\subset H$ 且 $\eta\le\varepsilon$。它是有向集。

对每个 $\alpha=(G,\varepsilon)$，现在直接构造 $\nu_\alpha$。记 $M=|\mu|(X)$；若 $M=0$，取
$\nu_\alpha=0$。以下设 $M>0$，并选 $0<\delta<\varepsilon/(2M)$。由 $|\mu|$ 的 Radon 内正则性，
存在紧集 $K_0$ 使 $|\mu|(X\setminus K_0)<\delta$。由于 $G$ 有限且每个 $f\in G$ 在无穷远消失，
存在紧集 $K_1$ 使
\[
 |f(x)|<\delta\qquad(x\notin K_1,\ f\in G).
\]
令 $K=K_0\cup K_1$。对每个 $x\in K$，有限多个函数的连续性给出一个开邻域 $U_x$，使
\[
 |f(y)-f(z)|<\delta\qquad(y,z\in U_x,\ f\in G).
\]
从这些邻域取有限子覆盖 $U_1,\ldots,U_N$，并把 $K$ 分成 Borel 片
\[
 A_1=K\cap U_1,\qquad
 A_j=K\cap\left(U_j\setminus\bigcup_{i<j}U_i\right)\ (2\le j\le N),
 \qquad A_0=X\setminus K.
\]
删去空片，在每个余下的 $A_j$ 中选 $x_j$，并置
\[
 \nu_\alpha=\sum_j\mu(A_j)\delta_{x_j}.
\]
若 $A_0=\varnothing$，以下所有下标为 $0$ 的项均约定删去。
对 $f\in G$，有限分割给
\begin{align*}
 \left|\int f\,\dd(\nu_\alpha-\mu)\right|
 &\leq\sum_{j\geq1}\int_{A_j}|f(x_j)-f(x)|\,|\mu|(\dd x)
      +\int_{A_0}|f(x_0)-f(x)|\,|\mu|(\dd x)\\
 &\leq\delta|\mu|(K)+2\delta|\mu|(A_0)
 \leq2\delta M<\varepsilon,
\end{align*}
并且
\[
 |\nu_\alpha|(X)\leq\sum_j|\mu(A_j)|\leq|\mu|(X),
 \qquad
 \nu_\alpha(X)=\sum_j\mu(A_j)=\mu(X).
\]
若 $\mu\geq0$，所有系数都非负，所以 $\nu_\alpha\geq0$。

固定 $f\in C_0(X)$ 和 $\delta>0$ 后，所有
$\alpha\succeq(\{f\},\delta)$ 都满足
$|\int f\dd(\nu_\alpha-\mu)|<\delta$，这正是弱星收敛。

若 $C_0(X)$ 可分，取稠密列 $(f_j)$，令 $\nu_n$ 同时以误差 $1/n$ 匹配
$f_1,\ldots,f_n$。记 $M=|\mu|(X)$；当 $M=0$ 结论直接成立。给定 $f$，先选 $f_j$ 使
$\norm{f-f_j}_\infty<\delta/(4M)$。当 $n\ge j$ 时，
\begin{align*}
 \left|\int f\dd(\nu_n-\mu)\right|
 &\le\left|\int(f-f_j)\dd\nu_n\right|
   +\left|\int f_j\dd(\nu_n-\mu)\right|
   +\left|\int(f_j-f)\dd\mu\right|\\
 &\le2M\norm{f-f_j}_\infty+\frac1n<\delta
\end{align*}
对充分大 $n$ 成立。非可分情形没有一个可数稠密测试族可供这项对角化使用，故一般结论必须表述成网；
这并不排除某个特殊测度恰好拥有某条原子逼近序列。
\end{proof}

\begin{theorem}[有限测试的确定性构造]\label{theorem:7-4-2-2}
设 $\mu$ 为局部紧 Hausdorff 空间 $X$ 上的有限符号或复 Radon 测度。给定有限
$G\subset C_0(X)$ 与 $\varepsilon>0$，存在有限 Borel 分割
$X=A_0\dot\cup\cdots\dot\cup A_N$ 及点 $x_j\in A_j$（空片删去），使
\[
 \nu=\sum_{j=0}^N\mu(A_j)\delta_{x_j}
\]
满足
\begin{equation}\label{eq:7-4-2-finite-test-error}
 \left|\int f\dd(\nu-\mu)\right|<\varepsilon\quad(f\in G),
 \qquad |\nu|(X)\le|\mu|(X),
 \qquad \nu(X)=\mu(X).
\end{equation}
正测度情形中 $\nu$ 也为正测度。
\end{theorem}

\begin{proof}
令 $M=|\mu|(X)$。若 $M=0$，取 $\nu=0$。以下设 $M>0$，并取
$0<\eta<\varepsilon/(2M)$。Radon 内正则性给紧集 $K_0$，使
$|\mu|(X\setminus K_0)<\eta$。又因 $G$ 有限且每个函数在无穷远消失，存在紧集 $K_1$，使
\begin{equation}\label{eq:7-4-2-tail-small-functions}
 |f(x)|<\eta\qquad(x\notin K_1,\ f\in G).
\end{equation}
令 $K=K_0\cup K_1$。有限多个连续函数可在紧集 $K$ 上用同一个有限开覆盖控制振幅：对每个
$x\in K$ 可取开邻域
$U_x$，使
\begin{equation}\label{eq:7-4-2-common-oscillation}
 |f(y)-f(z)|<\eta
 \qquad(y,z\in U_x,\ f\in G).
\end{equation}
从 $(U_x)_{x\in K}$ 取有限子覆盖 $U_1,\ldots,U_N$，再令
\[
 A_1=K\cap U_1,\qquad
 A_j=K\cap\left(U_j\setminus\bigcup_{i<j}U_i\right),\qquad
 A_0=X\setminus K.
\]
删去空片，并在每个非空 $A_j$ 中选 $x_j$。这是有限 Borel 分割。若 $A_0\ne\varnothing$，
由 \eqref{eq:7-4-2-tail-small-functions}，$x_0$ 也使所有 $|f(x_0)|<\eta$。

对 $f\in G$，直接按分割计算
\begin{align*}
 \left|\int f\dd\nu-\int f\dd\mu\right|
 &=\left|\sum_{j=0}^N\int_{A_j}(f(x_j)-f(x))\dd\mu(x)\right|\\
 &\le\sum_{j=1}^N\eta|\mu|(A_j)+2\eta|\mu|(A_0)\\
 &\le2\eta M<\varepsilon.
\end{align*}
这里第一行只是有限和与积分的线性，没有交换无穷操作；第二行分别用了
\eqref{eq:7-4-2-common-oscillation} 和尾部两个函数值都小于 $\eta$。
此外，
\[
 |\nu|(X)\le\sum_{j=0}^N|\mu(A_j)|
 \le\sum_{j=0}^N|\mu|(A_j)=|\mu|(X),
\]
而 $\nu(X)=\sum_j\mu(A_j)=\mu(X)$。若 $\mu\ge0$，所有系数 $\mu(A_j)$ 非负，正性也被保留。
\end{proof}

在紧度量空间上，若 $G$ 中每个函数都是 $L$-Lipschitz，且分割片直径不超过 $h$，上述逐片计算立即给
\begin{equation}\label{eq:7-4-2-lipschitz-atomic-error}
 \left|\int f\dd(\nu-\mu)\right|
 \le Lh\,|\mu|(X).
\end{equation}
中点求积以网格直径控制误差，经验测度以大数律控制误差，而下一命题用有限维凸性把指定测试上的误差
直接降为零。

\begin{theorem}[经验测度强一致性]\label{theorem:7-4-2-3}
设 $X$ 为 Polish 空间，$X_1,X_2,\ldots$ 独立同分布，共同 Borel 概率分布为 $\mu$。则
\[
 L_n\Rightarrow\mu\qquad\text{a.s.}
\]
其中 $\Rightarrow$ 表示对所有有界连续函数的弱收敛。
\end{theorem}

\begin{proof}
先把所需的可数决定类真正构造出来。由 Polish 假设，取 $X$ 上一个兼容完备度量 $d$ 和一个可数
稠密集 $D=\{z_i:i\geq1\}$。对 $i\geq1$ 与有理数 $0<r<s$，定义帽函数
\[
 \psi_{i,r,s}(x)=
 \begin{cases}
  1,&d(x,z_i)\leq r,\\
  \dfrac{s-d(x,z_i)}{s-r},&r<d(x,z_i)<s,\\
  0,&d(x,z_i)\geq s.
 \end{cases}
\]
它取值于 $[0,1]$，且 Lipschitz 常数不超过 $(s-r)^{-1}$。令 $\mathcal F_0$ 由常数函数 $1$ 和任意
有限多个上述帽函数的逐点最大值组成；$\mathcal F_0$ 可数，并且其中每个函数都是有界 Lipschitz
函数。

这个固定族决定概率测度的弱收敛。事实上，设概率 $\rho_n$ 对 $\mathcal F_0$ 中每个函数的积分都收敛
到 $\rho$ 的相应积分。空开集的 Portmanteau 不等式没有内容。给定非空开集 $U$ 与 $x\in U$，若
$U\neq X$，记 $a=d(x,U^c)>0$，先取 $z_i$ 使 $d(x,z_i)<a/3$，再选有理数
$r<s$，使
\[
 d(x,z_i)<r<s<d(z_i,U^c).
\]
这样的有理数存在，因为
$d(z_i,U^c)\geq a-d(x,z_i)>2a/3>d(x,z_i)$。
于是 $\psi_{i,r,s}(x)=1$ 且 $\supp\psi_{i,r,s}\subset U$。因此，所有支撑包含于 $U$ 的
$\mathcal F_0$-函数可以枚举为 $(h_j)$；令 $H_m=h_1\vee\cdots\vee h_m$，则
$0\leq H_m\uparrow\1_U$。当 $U=X$ 时直接取 $H_m=1$。对每个固定 $m$，
\[
 \liminf_{n\to\infty}\rho_n(U)
 \geq\lim_{n\to\infty}\int H_m\dd\rho_n
 =\int H_m\dd\rho.
\]
令 $m\to\infty$ 并用单调收敛，得到
$\liminf_n\rho_n(U)\geq\rho(U)$。Portmanteau 定理~\ref{theorem:5-3-3-1} 因而给
$\rho_n\Rightarrow\rho$。这就由上述帽函数构造证明了所需的可数有界 Lipschitz 决定性质。

枚举 $\mathcal F_0=\{f_m:m\geq1\}$。固定 $m$。$f_m$ 有界，故 $f_m(X_k)$ i.i.d. 且可积。
强大数律给概率一事件 $\Omega_m$，使
\[
 \int f_m\dd L_n
 =\frac1n\sum_{k=1}^nf_m(X_k)
 \longrightarrow\E f_m(X_1)=\int f_m\dd\mu.
\]
可数交 $\Omega_0=\bigcap_{m\ge1}\Omega_m$ 仍有概率一。在每个
$\omega\in\Omega_0$ 上，上式对全部 $m$ 同时成立；收敛决定类的定义遂给
$L_n(\omega)\Rightarrow\mu$。

在 $X=\R$ 时，取半轴指标的连续上下夹逼，可把同一结论写成经验分布函数
$F_n(t)=n^{-1}\sum_k\1_{\{X_k\le t\}}$ 在 $F(t)=\mu(( -\infty,t])$ 的每个连续点收敛。
这里仍只需处理可数个决定点，再由单调性夹逼其余连续点；不能对不可数个 $t$ 直接取概率一事件的交。
\end{proof}

\begin{proposition}[Tchakaloff 有限矩匹配]\label{proposition:7-4-2-4}
设 $K$ 为紧 Hausdorff 空间，$\mu$ 为其上的有限正 Radon 测度，$V\subset C(K;\R)$ 为 $m$ 维线性空间。
则存在正原子测度
\[
 \nu=\sum_{j=1}^rc_j\delta_{x_j},\qquad
 0\leq r\leq m+1,
\]
使
\[
 \nu(K)=\mu(K),\qquad
 \int_Kf\dd\nu=\int_Kf\dd\mu\quad(f\in V).
\]
当 $r>0$ 时各 $c_j>0$；$r=0$ 表示空和，即零测度。
若 $1\in V$，可改进为 $r\le m$。这两个一般界都是最佳的。
\end{proposition}

\begin{proof}
若 $\mu(K)=0$，取 $r=0$ 的空和即得零测度。以下令 $M=\mu(K)>0$，并取 $V$ 的基
$f_1,\ldots,f_m$。把总质量坐标一并加入，定义连续映射
\[
 \Phi:K\to\R^{m+1},\qquad
 \Phi(x)=(f_1(x),\ldots,f_m(x),1),
\]
以及概率 $P=M^{-1}\mu$。向量
\[
 b=\int_K\Phi(x)\dd P(x)
\]
属于 $S:=\Phi(K)$ 的闭凸包。下面同时证明所需的凸包闭性，并给出有限维 \Caratheodory{} 约化步骤。
令 $\tau=\Phi_*P$。对每个 $n$，用有限个半径小于 $1/(2n)$ 的 Euclidean 球覆盖
紧集 $S$，再逐次作差得到有限 Borel 分割
$S=E_{n,1}\dot\cup\cdots\dot\cup E_{n,q_n}$，其中每个非空分片的直径小于 $1/n$。在每片选
$y_{n,j}\in E_{n,j}$，置
\[
 b_n=\sum_j\tau(E_{n,j})y_{n,j}\in\operatorname{conv}S.
\]
由于
\[
 \norm{b-b_n}
 =\left\|\sum_j\int_{E_{n,j}}(y-y_{n,j})\,\tau(\dd y)\right\|
 \leq\frac1n,
\]
故 $b\in\overline{\operatorname{conv}S}$。

再证明这个凸包已经闭。记 $d$ 为 $S$ 的仿射维数；由于 $S$ 落在末坐标等于 $1$ 的仿射超平面内，
$d\leq m$。若一个凸组合 $y=\sum_{j=1}^r\theta_jy_j$ 使用 $r>d+1$ 个正权重，则这些点仿射相关，
故存在不全为零的实数 $a_j$，使
\[
 \sum_{j=1}^ra_j=0,
 \qquad
 \sum_{j=1}^ra_jy_j=0.
\]
把 $(a_j)$ 乘以 $-1$（若有必要），可设某个 $a_j>0$。取
\[
 t=\min_{a_j>0}\frac{\theta_j}{a_j}>0,
 \qquad \theta_j'=\theta_j-ta_j.
\]
则每个 $\theta_j'\geq0$，至少一个为零，并且
$\sum_j\theta_j'=1$、$\sum_j\theta_j'y_j=y$。反复消去零权，便知每个凸包中的点都能用至多
$d+1$ 个点表示。于是 $\operatorname{conv}S$ 是紧集
\[
 S^{d+1}\times
 \left\{(\theta_1,\ldots,\theta_{d+1})\in[0,1]^{d+1}:\sum_j\theta_j=1\right\}
\]
在连续映射 $(y_1,\ldots,y_{d+1},\theta)\mapsto\sum_j\theta_jy_j$ 下的像，因而为闭集。
所以 $b\in\operatorname{conv}S$，并存在 $x_1,\ldots,x_r\in K$ 和 $\theta_j\geq0$，使
\[
 r\leq d+1\leq m+1,\qquad \sum_j\theta_j=1,\qquad
 b=\sum_{j=1}^r\theta_j\Phi(x_j).
\]
删去零权并令 $c_j=M\theta_j$，便同时得到总质量和基函数积分的匹配；线性性把它推广到全部 $V$。

若 $1\in V$，可选 $f_1=1$。这在“末坐标等于 $1$”之外再给出关系 $y_1=y_{m+1}$，故
$\Phi(K)$ 的仿射维数至多 $m-1$。把上面的仿射相关消元用于这个较小的仿射空间，只需至多 $m$ 个点。

最后核对界的准确性。不假定 $1\in V$ 时，令 $K=\{0,e_1,\ldots,e_m\}\subset\R^m$，令 $V$ 由
$m$ 个坐标函数在 $K$ 上的限制张成，并让 $\mu$ 在这 $m+1$ 个点上都给正质量。归一化积分向量位于
单纯形内部，不能由少于 $m+1$ 个顶点的凸组合表示。若 $1\in V$，取含 $m$ 个点的有限空间
$K$、$V=C(K)$ 及在每点都给正质量的 $\mu$；精确匹配所有 $V$ 中函数会逐点匹配质量，故必须使用
全部 $m$ 个原子。
\end{proof}

\begin{figure}[htbp]
\centering
\begin{tikzpicture}[node distance=11mm and 15mm]
  \node[book strong node] (space) {空间 $K$\\点 $x$};
  \node[book strong node,right=of space] (image) {有限观测像\\$\Phi(x)\in\R^m$};
  \node[book warm node,right=of image] (bar) {测度重心\\$b=\int\Phi\dd P$};
  \node[book strong node,below=of bar] (atom) {有限凸组合\\$\sum_{j=1}^r\theta_j\Phi(x_j)$};
  \node[book warm node,left=of atom] (emp) {经验重心\\$n^{-1}\sum_{k=1}^n\Phi(X_k)$};
  \draw[book accent arrow] (space)--node[above,book label]{$m$ 个测试}(image);
  \draw[book accent arrow] (image)--node[above,book label]{积分}(bar);
  \draw[book accent arrow] (atom)--node[right,book label]{精确匹配}(bar);
  \draw[book dashed arrow] (emp)--node[below,book label]{大数律}(atom);
  \draw[book dashed arrow] (space)--node[left,book label]{抽样}(emp);
\end{tikzpicture}
\caption{有限测试把测度问题降到有限维重心问题。确定性分割近似重心，经验测度随机逼近重心，Tchakaloff 用有限凸组合精确表示重心。}
\label{fig:7-4-2-finite-observation-barycenter}
\end{figure}

图~\ref{fig:7-4-2-finite-observation-barycenter} 只在有限测试像中宣称精确。若不断扩大测试族，必须另用
网、可数决定类或稠密性完成极限；一次 Tchakaloff 表示绝不意味着原测度本身已经是有限原子的。

\subsection{结果汇总：由实分析通往对偶空间、矩问题与谱理论}\label{subsec:7-4-3}
\index{结果汇总}\index{Feller 半群}\index{生成元}\index{预解算子}

前七章反复出现同一个四阶段结构：有限测试先给一致的数据，完备性、扩张或紧性制造无限对象，稠密类
或变换负责识别，核与半群再描述对象如何演化。下面先用三个例子说明从测试信息读取
表示对象时不可省略的假设，再由三道综合题分别展开 Fourier 完备化、矩泛函表示和随机半群这
三条路线。一般弱星紧性、闭算子与谱测度会在后续泛函分析中推广这些构造。

\begin{example}[变分到测度]\label{ex:7-4-3-1}
设 $a<x_0<b$，$f,h\in C^1[a,b]$，$\Phi\in C^1(\R)$。设
$L:[a,b]\times\R^2\to\R$ 连续，并且对后两个变量的偏导数 $L_u,L_p$ 存在且连续。对
\[
 S(f)=\Phi(f(x_0))+\int_a^bL(x,f(x),f'(x))\dd x
\]
考虑 $f+\varepsilon h$。映射
\[
 (x,\varepsilon)\longmapsto
 \bigl(x,f(x)+\varepsilon h(x),f'(x)+\varepsilon h'(x)\bigr)
\]
把紧集 $[a,b]\times[-\varepsilon_0,\varepsilon_0]$ 送到紧集，因而 $L_u,L_p$ 在该像上有共同有限上界；
差商由 $C(|h|+|h'|)$ 支配。支配收敛定理给
\begin{equation}\label{eq:7-4-3-first-variation}
 \delta S(f)[h]
 =\Phi'(f(x_0))h(x_0)
 +\int_a^b\bigl(L_u(x,f,f')h+L_p(x,f,f')h'\bigr)\dd x.
\end{equation}
若 $x\mapsto L_p(x,f(x),f'(x))$ 为 $C^1$，分部积分得到
\begin{align}\label{eq:7-4-3-variation-measure}
 \delta S(f)[h]
 &=\int_a^b\left(L_u-\frac{\dd}{\dd x}L_p\right)h\dd x
 +L_p(b,f(b),f'(b))h(b)-L_p(a,f(a),f'(a))h(a)\notag\\
 &\quad+\Phi'(f(x_0))h(x_0).
\end{align}
因此该一阶变分在上确界范数下由下列符号测度表示：
\[
 \nu=\left(L_u-\frac{\dd}{\dd x}L_p\right)\dd x
      -L_p(a,f(a),f'(a))\delta_a
      +L_p(b,f(b),f'(b))\delta_b
      +\Phi'(f(x_0))\delta_{x_0},
\]
其中密度中的 $L_u,L_p$ 均沿曲线
$x\mapsto(x,f(x),f'(x))$ 取值；式~\eqref{eq:7-4-3-variation-measure} 正是
$\delta S(f)[h]=\int h\dd\nu$。若只知道 \eqref{eq:7-4-3-first-variation} 而不能对 $L_p$ 分部积分，
含 $h'$ 的项只在 $C^1$ 范数下连续，未必由测度表示。例~\ref{ex:7-4-1-3} 的导数泛函正是这个边界。
\end{example}

\begin{example}[矩到弱星极限]\label{ex:7-4-3-2}
设 $K=[-1,1]$，概率测度 $\mu_N,\mu$ 都支持在 $K$，并且二者的 $0$ 到 $N$ 阶矩相同。给定
$f\in C(K)$ 和
$\varepsilon>0$，Stone--Weierstrass 定理给多项式 $p$，使
$\norm{f-p}_\infty<\varepsilon$；若 $N\ge\deg p$，则
\[
 \left|\int f\dd\mu_N-\int f\dd\mu\right|
 \le\int|f-p|\dd\mu_N+\int|f-p|\dd\mu<2\varepsilon.
\]
所以共同紧支撑和总质量控制把逐阶矩匹配升级为弱收敛。若删去共同紧支撑，仅有所有 Hankel 二次型非负
不能给紧性；例如 $\delta_N$ 的质量直接逃向无穷。
\end{example}

\begin{example}[半群到生成元]\label{ex:7-4-3-3}
在 $\R^d$ 上令 $P_tf(x)=\E f(x+B_t)$，其中 $B_t\sim N(0,tI_d)$。对 Fourier 模式
$e_\xi(x)=e^{i\xi\cdot x}$，Gaussian 特征函数给
\[
 P_te_\xi(x)=e^{-t|\xi|^2/2}e_\xi(x),
 \qquad
 \frac{P_te_\xi-e_\xi}{t}\longrightarrow-\frac{|\xi|^2}{2}e_\xi
 =\frac12\Delta e_\xi.
\]
这项计算确定了 Fourier 符号 $-|\xi|^2/2$，但 $e_\xi\notin C_0(\R^d)$，所以它不能替代生成元定义域的证明。
要把符号识别为生成元，还必须在生成元的定义域上证明上确界范数极限。综合题
\ref{problem:7-4-3-4} 对 $C_c^\infty(\R^d)$ 中的函数完成这一步。
\end{example}

\begin{remark}[假设的必要性]
表~\ref{tab:7-4-3-assumptions-counterexamples} 汇总了七个基本反例或失效情形。它们分别说明单调连续性、
局部紧性、连续性、可分性、可数决定类、支集正性和 Feller 条件为何出现在相应结论中。
\end{remark}

\begin{table}[htbp]
\centering
\small
\begin{tabularx}{\textwidth}{>{\raggedright\arraybackslash}p{0.19\textwidth}
 >{\raggedright\arraybackslash}p{0.22\textwidth}
 >{\raggedright\arraybackslash}p{0.31\textwidth}
 >{\raggedright\arraybackslash}X}
\toprule
结论 & 关键假设 & 反例或失效情形 & 相关主题 \\
\midrule
Daniell--Stone & $f_n\downarrow0\Rightarrow I(f_n)\downarrow0$
& 把收敛数列空间上的通常极限用 Hahn--Banach 保范延拓为
$I:\ell^\infty(\N)\to\R$；此延拓满足 $I(\mathbf 1)=1$。若 $0\le x\le M\mathbf 1$，则
$\|M\mathbf 1-x\|_\infty\le M$，从而
$I(M\mathbf 1-x)\le |I(M\mathbf 1-x)|\le M$，并有
$I(x)=M-I(M\mathbf 1-x)\ge0$，故 $I$ 正；但
$I(\1_{\{k\ge n\}})=1$，而这些函数逐点下降到零，故相应有限加法不能升级为可数加法
& 测度扩张、概率分布 \\
Riesz--Markov & 局部紧 Hausdorff 与局部 Urysohn 帽函数
& 没有紧开夹逼时，开集值的局部有限性与内正则性均无由推出
& $C_0(X)^*$、Radon 对偶 \\
测度表示导数 & $C_0$ 上确界范数连续
& $f_n(x)=n^{-1}\phi(nx)$ 使 $\norm{f_n}_\infty\to0$ 而 $f_n'(0)$ 不变
& 分布与 Sobolev 对偶 \\
原子逼近序列 & $C_0(X)$ 可分
& 非可分时弱星邻域没有可数稠密测试族；一般密度结论只能由有向网表达
& 弱星拓扑 \\
经验测度弱收敛 & Polish 空间的可数决定类
& 对不可数测试族分别取概率一事件，所得不可数交未必仍有概率一
& Varadarajan、经验过程 \\
$K$-支撑矩表示 & 对所有在 $K$ 上非负的多项式为正
& $K=[-1,1]$ 时 $L(p)=p(2)$ 满足 $L(q^2)\ge0$，却有
$L(1-x^2)=-3<0$；Hankel 正性不能保证表示测度支撑于 $K$
& Hausdorff 矩问题 \\
Feller 半群 & $K(C_0)\subset C_0$
& $Q(x)=\delta_0$（$x\in\Q$）、$Q(x)=\delta_1$（$x\notin\Q$）使
$Kf$ 对 $f(0)\ne f(1)$ 不连续
& 生成元理论 \\
\bottomrule
\end{tabularx}
\caption{结论、关键假设与反例。第三列说明每项假设在相应构造中所承担的作用。}
\label{tab:7-4-3-assumptions-counterexamples}
\end{table}

\begin{figure}[htbp]
\centering
\begin{tikzpicture}[node distance=8mm and 17mm]
  \node[book strong node] (obs) {有限观测\\邻域、测试、矩、有限维分布};
  \node[book strong node,below=of obs] (exist) {存在\\完备化、扩张、紧性};
  \node[book strong node,below=of exist] (id) {识别\\稠密、微分、变换};
  \node[book strong node,below=of id] (evo) {演化\\概率核、半群、预解算子};
  \node[book warm node,right=of exist] (dual) {对偶空间\\$C_0^*$、矩泛函};
  \node[book warm node,right=of evo] (spectral) {算子方向\\生成元、预解式};
  \draw[book accent arrow] (obs)--node[right,book label]{一致有限数据}(exist);
  \draw[book accent arrow] (exist)--node[right,book label]{候选对象}(id);
  \draw[book accent arrow] (id)--node[right,book label]{唯一结构}(evo);
  \draw[book dashed arrow] (obs.east)--(dual.north);
  \draw[book dashed arrow] (id.east)--(dual.south);
  \draw[book dashed arrow] (evo)--(spectral);
\end{tikzpicture}
\caption{有限观测到表示与演化的四层结构。三项综合题分别展开 Fourier 完备化、矩泛函表示和随机半群三条路线。}
\label{fig:7-4-3-interface-master}
\end{figure}

图~\ref{fig:7-4-3-interface-master} 中“存在”和“识别”不能互换：先有一族相容 Fourier 部分和，仍须
$L^2$ 完备性产生极限；先有矩截断解，仍须紧支撑和稠密多项式识别极限；先有转移核，仍须
$C_0$ 不变性与强连续性，才谈得上该 Banach 空间中的生成元。

\begin{problem}[综合题 A：三角矩与 $L^2$ 完备化]\label{problem:7-4-3-2}
给定 $(a_k)_{k\in\Z}\in\ell^2(\Z)$。仅用有限 Fourier 和、正交性、\Fejer{} 逼近与 $L^2$ 完备性，
证明存在唯一 $f\in L^2(\T)$ 使 $\widehat f(k)=a_k$ 对所有 $k\in\Z$ 成立，并证明
$\norm f_2^2=\sum_k|a_k|^2$。
\end{problem}

\begin{proof}[解答]
在 $\T=\R/(2\pi\Z)$ 上使用归一化 Haar 测度 $\dd m=\dd x/(2\pi)$，并置
$e_k(x)=e^{ikx}$。直接积分给
\[
 \langle e_j,e_k\rangle
 =\frac1{2\pi}\int_{-\pi}^{\pi}e^{i(j-k)x}\dd x
 =\begin{cases}1,&j=k,\\0,&j\ne k.\end{cases}
\]
定义有限和
\[
 S_N=\sum_{|k|\le N}a_ke_k.
\]
若 $M>N$，有限正交展开给出精确恒等式
\begin{equation}\label{eq:7-4-3-fourier-cauchy}
 \norm{S_M-S_N}_2^2=\sum_{N<|k|\le M}|a_k|^2.
\end{equation}
因为 $(a_k)\in\ell^2$，右端随 $M,N\to\infty$ 趋于零。因此 $(S_N)$ 是 $L^2(\T)$ 中的 Cauchy
列。$L^2$ 完备性给出 $f\in L^2(\T)$，使 $S_N\to f$ 于 $L^2$。这正是由有限 Fourier 和产生
极限函数的存在性步骤。

固定 $j\in\Z$。Cauchy--Schwarz 不等式和 $\norm{e_j}_2=1$ 给
\[
 |\widehat f(j)-\widehat{S_N}(j)|
 =|\langle f-S_N,e_j\rangle|
 \le\norm{f-S_N}_2\longrightarrow0.
\]
当 $N\ge|j|$ 时 $\widehat{S_N}(j)=a_j$，所以 $\widehat f(j)=a_j$。再由范数连续性与有限
Parseval 恒等式，
\[
 \norm f_2^2=\lim_{N\to\infty}\norm{S_N}_2^2
 =\lim_{N\to\infty}\sum_{|k|\le N}|a_k|^2
 =\sum_{k\in\Z}|a_k|^2.
\]

只剩唯一性。若 $g\in L^2(\T)$ 的全部 Fourier 系数为零，则它的 \Fejer{} 均值为
\[
 \sigma_Ng
 =\sum_{|k|\le N}\left(1-\frac{|k|}{N+1}\right)\widehat g(k)e_k=0.
\]
定理~\ref{theorem:7-2-1-1} 的 \Fejer{} 近似先对连续函数给一致收敛；连续函数在 $L^2$ 中稠密，
而卷积正核满足 $\norm{\sigma_Nh}_2\le\norm h_2$。这里的收缩并非另行调用 Plancherel：若
$K_N\geq0$ 是归一化 \Fejer{} 核，则 Jensen 不等式逐点给
\[
 |\sigma_Nh(x)|^2
 =\left|\int_\T K_N(y)h(x-y)\dd m(y)\right|^2
 \leq\int_\T K_N(y)|h(x-y)|^2\dd m(y).
\]
对 $x$ 积分并用 Tonelli 以及 Haar 测度的平移不变性，得到
$\norm{\sigma_Nh}_2^2\leq\norm h_2^2$。所以一个三项逼近给
$\sigma_Ng\to g$ 于 $L^2$。具体地，给定连续 $q$，
\[
 \norm{\sigma_Ng-g}_2
 \le\norm{\sigma_N(g-q)}_2+\norm{\sigma_Nq-q}_2+\norm{q-g}_2
 \le2\norm{g-q}_2+\norm{\sigma_Nq-q}_\infty.
\]
先选 $q$，再令 $N\to\infty$，右端趋于零。因此 $g=0$。若两个函数具有同一系数，对其差应用此论证
便得二者在 $L^2$ 中相同。

综上，正交性给出 Cauchy 估计~\eqref{eq:7-4-3-fourier-cauchy}，$L^2$ 完备性产生极限，
系数泛函的连续性恢复给定数据，\Fejer{} 逼近则给出唯一性。
\end{proof}

\begin{problem}[综合题 B：多项式矩与 Stieltjes 变换]\label{problem:7-4-3-3}
固定紧区间 $K=[a,b]$。设 $L:\R[x]\to\R$ 线性，并且每个在 $K$ 上非负的实多项式 $p$ 都满足
$L(p)\ge0$。证明 $L$ 由 $K$ 上唯一有限正测度 $\mu$ 表示；构造在 $K$ 上匹配 $0$ 到 $N$ 阶矩的
有限原子测度 $\mu_N$；证明 $\mu_N\Rightarrow\mu$，并证明其 Stieltjes 变换在
$\C\setminus K$ 上局部一致收敛。
\end{problem}

\begin{proof}[解答]
\textbf{第一步：把多项式泛函延到 $C(K)$。}
正性给 $L(1)\ge0$。对任意实多项式 $p$，令 $M_p=\norm p_{C(K)}$。多项式
$M_p\pm p$ 在 $K$ 上非负，故
\[
 0\le L(M_p\pm p)=M_pL(1)\pm L(p).
\]
于是
\begin{equation}\label{eq:7-4-3-polynomial-bound}
 |L(p)|\le L(1)\norm p_{C(K)}.
\end{equation}
若 $L(1)=0$，该式给 $L=0$，以下结论取零测度即可。设 $L(1)>0$。Stone--Weierstrass 定理说明
实多项式在 $C(K;\R)$ 中一致稠密；估计 \eqref{eq:7-4-3-polynomial-bound} 因而给唯一连续延拓
$\widetilde L:C(K)\to\R$，且 $\norm{\widetilde L}\le L(1)$。

还须验证延拓保持正性。给 $f\in C(K)$、$f\ge0$，取多项式 $p_n\to f$ 一致，并令
$\varepsilon_n=\norm{p_n-f}_\infty$。则 $p_n+\varepsilon_n\ge0$ 于 $K$，所以
\[
 L(p_n)+\varepsilon_nL(1)\ge0.
\]
令 $n\to\infty$ 得 $\widetilde L(f)\ge0$。定理~\ref{theorem:7-4-1-2} 应用于紧空间 $K$，得到唯一
有限正 Radon 测度 $\mu$，满足
\[
 L(p)=\int_Kp(x)\dd\mu(x),\qquad \mu(K)=\widetilde L(1)=L(1).
\]

\textbf{第二步：Tchakaloff 原子截断。}
对 $N\ge0$ 令
$V_N=\Span\{1,x,\ldots,x^N\}$。它的维数为 $N+1$ 且含常数。命题
\ref{proposition:7-4-2-4} 给出支撑仍在 $K$ 的正原子测度
\[
 \mu_N=\sum_{j=1}^{r_N}c_{N,j}\delta_{x_{N,j}},
 \qquad r_N\le N+1,\quad c_{N,j}>0,\quad x_{N,j}\in K,
\]
使
\begin{equation}\label{eq:7-4-3-moment-match}
 \mu_N(K)=\mu(K),\qquad
 \int_Kx^k\dd\mu_N=\int_Kx^k\dd\mu=L(x^k)\quad(0\le k\le N).
\end{equation}
这里的 $N+1$ 而不是 $N+2$，来自 $1\in V_N$ 所降低的一维仿射维数。

\textbf{第三步：弱收敛。}
记 $M=L(1)$。给定 $f\in C(K)$ 和 $\varepsilon>0$，选多项式 $p$ 使
$\norm{f-p}_\infty<\varepsilon/(2M+1)$。当 $N\ge\deg p$ 时，
\eqref{eq:7-4-3-moment-match} 和线性性给 $\int p\dd\mu_N=\int p\dd\mu$，从而
\begin{align*}
 \left|\int f\dd\mu_N-\int f\dd\mu\right|
 &\le\int|f-p|\dd\mu_N+\int|f-p|\dd\mu\\
 &\le2M\norm{f-p}_\infty<\varepsilon.
\end{align*}
故 $\mu_N\Rightarrow\mu$。共同紧支集使紧性自动成立；证明甚至不需要抽取子列。

\textbf{第四步：Stieltjes 变换局部一致收敛。}
采用
\[
 S_\rho(z)=\int_K\frac1{z-x}\dd\rho(x),\qquad z\in\C\setminus K
\]
的符号。固定紧集 $C\Subset\C\setminus K$，令
$d=\dist(C,K)>0$。核满足共同界
\begin{equation}\label{eq:7-4-3-stieltjes-kernel-bounds}
 \left|\frac1{z-x}\right|\le d^{-1},\qquad
 \left|\frac1{z-x}-\frac1{w-x}\right|
 \le d^{-2}|z-w|\quad(z,w\in C,x\in K).
\end{equation}
第一项是所有积分存在以及参数积分的支配；第二项给
\[
 |S_{\mu_N}(z)-S_{\mu_N}(w)|\le Md^{-2}|z-w|,
 \qquad
 |S_\mu(z)-S_\mu(w)|\le Md^{-2}|z-w|.
\]
对每个固定 $z\in C$，函数 $x\mapsto(z-x)^{-1}$ 连续，弱收敛给
$S_{\mu_N}(z)\to S_\mu(z)$。给定 $\varepsilon>0$，取 $C$ 的有限
$\eta$-网 $z_1,\ldots,z_q$，其中 $2Md^{-2}\eta<\varepsilon/2$。在这有限多个点上同时取足够大的
$N$，使差小于 $\varepsilon/2$。对任意 $z\in C$ 选 $z_j$ 满足 $|z-z_j|<\eta$，于是
\begin{align*}
 |S_{\mu_N}(z)-S_\mu(z)|
 &\le|S_{\mu_N}(z)-S_{\mu_N}(z_j)|
 +|S_{\mu_N}(z_j)-S_\mu(z_j)|\\
 &\quad+|S_\mu(z_j)-S_\mu(z)|<\varepsilon.
\end{align*}
所以收敛在每个 $C\Subset\C\setminus K$ 上一致。

最后说明正性假设为何必须精确写成“在 $K$ 上非负”。仅有 Hankel 正性
$L(q^2)\ge0$ 不足以保证表示测度支撑于 $K$：在 $K=[-1,1]$ 上取 $L(p)=p(2)$，所有平方都取非负值，但
$1-x^2\ge0$ 于 $K$ 而 $L(1-x^2)=-3$。相应表示测度是 $\delta_2$，支撑恰在 $K$ 外。
\end{proof}

\begin{problem}[综合题 C：随机演化与算子方法]\label{problem:7-4-3-4}
设 $X$ 为局部紧 Hausdorff 空间，$(Q_t)_{t\ge0}$ 是满足 Chapman--Kolmogorov 方程的 Borel 转移核族，
$Q_0(x,\cdot)=\delta_x$，并令
$P_tf(x)=\int f(y)Q_t(x,\dd y)$。假定 $P_tC_0(X)\subset C_0(X)$ 且
$\norm{P_tf-f}_\infty\to0$（$t\downarrow0$）对每个 $f\in C_0(X)$ 成立。证明 $(P_t)$ 是
$C_0(X)$ 上的正收缩强连续半群；计算 Brownian 半群在 $C_c^\infty(\R^d)$ 上的生成元，并计算
满足 Feller 条件的 Poisson 跳半群在整个 $C_0(X)$ 上的生成元；最后对任意 $\alpha,\beta>0$ 证明
Laplace 预解算子恒等式
\[
 R_\alpha-R_\beta=(\beta-\alpha)R_\alpha R_\beta,
 \qquad R_\alpha f=\int_0^\infty e^{-\alpha t}P_tf\dd t.
\]
\end{problem}

\begin{proof}[解答]
\textbf{第一步：从转移族得到 Feller 半群。}
对 $f\in C_0(X)$，概率核给
\[
 |P_tf(x)|\le\int|f(y)|Q_t(x,\dd y)\le\norm f_\infty,
\]
故 $P_t$ 为正收缩。Chapman--Kolmogorov 及有界函数的核积分结合律给
\begin{align*}
 P_sP_tf(x)
 &=\int_X\int_Xf(z)Q_t(y,\dd z)Q_s(x,\dd y)\\
 &=\int_Xf(z)Q_{s+t}(x,\dd z)=P_{s+t}f(x).
\end{align*}
两重积分的绝对值由 $\norm f_\infty$ 支配，概率核上的 Tonelli/Fubini 定理保证交换合法。
$P_0=I$。零时刻的强连续性已假定；收缩性把它传播到每个时刻。对 $h\downarrow0$，
\[
 \norm{P_{t+h}f-P_tf}_\infty
 \le\norm{P_hf-f}_\infty\to0,
\]
而 $0<h<t$ 时
\[
 \norm{P_tf-P_{t-h}f}_\infty
 \le\norm{P_hf-f}_\infty\to0.
\]
因此 $(P_t)$ 是 $C_0(X)$ 上的正收缩强连续半群。其生成元定义为
\[
 D(A)=\left\{f\in C_0(X):
 \lim_{t\downarrow0}\frac{P_tf-f}{t}\text{ 在 }\norm{\cdot}_\infty\text{ 中存在}\right\},
 \qquad Af=\lim_{t\downarrow0}\frac{P_tf-f}{t}.
\]

\textbf{第二步：Brownian 生成元。}
在 $X=\R^d$ 上令 $P_tf(x)=\E f(x+\sqrt t Z)$，其中 $Z\sim N(0,I_d)$。取
$f\in C_c^\infty(\R^d)$。二阶 Taylor 公式的积分余项给
\begin{align*}
 f(x+y)-f(x)
 &=\nabla f(x)\cdot y+\frac12y^TD^2f(x)y+r(x,y),\\
 |r(x,y)|
 &\le\frac12|y|^2\omega_{D^2f}(|y|),
\end{align*}
其中
$\omega_{D^2f}(r)=\sup_{|x-x'|\le r}\norm{D^2f(x)-D^2f(x')}$。因为
$D^2f\in C_c(\R^d)$，它在整个 $\R^d$ 上一致连续，故该模量随 $r\downarrow0$ 趋于零。
Gaussian 一、二阶矩为
$\E Z_i=0$、$\E Z_iZ_j=\delta_{ij}$，所以
\[
 \frac{P_tf(x)-f(x)}t-\frac12\Delta f(x)
 =\frac1t\E r(x,\sqrt t Z).
\]
并且一致于 $x$ 有
\begin{equation}\label{eq:7-4-3-brownian-generator-domination}
 \frac1t|r(x,\sqrt t Z)|
 \le\frac12|Z|^2\omega_{D^2f}(\sqrt t|Z|)
 \le\norm{D^2f}_\infty|Z|^2.
\end{equation}
右端可积，左端对每个固定 $Z$ 趋于零。支配收敛定理应用于
\eqref{eq:7-4-3-brownian-generator-domination}，并在取期望前先取 $x$ 的上确界，得到
\[
 \left\|\frac{P_tf-f}{t}-\frac12\Delta f\right\|_\infty\longrightarrow0.
\]
故 $C_c^\infty(\R^d)\subset D(A)$ 且 $Af=\frac12\Delta f$。这一步给的是定义域内的范数极限，
不是只在 Fourier 模式上的形式符号。

\textbf{第三步：Poisson 跳生成元。}
设 $Q(x,\dd y)$ 为单次跳跃核，
$Kf(x)=\int f(y)Q(x,\dd y)$，并明确假设 Feller 条件
$K(C_0(X))\subset C_0(X)$。$K$ 是正收缩。若跳跃时钟速率为 $\lambda>0$，则
\begin{equation}\label{eq:7-4-3-poisson-semigroup-series}
 P_tf=e^{-\lambda t}\sum_{n=0}^\infty\frac{(\lambda t)^n}{n!}K^nf.
\end{equation}
该级数在 $C_0$ 范数中绝对收敛，因为
$\norm{K^nf}_\infty\le\norm f_\infty$。它确实定义强连续半群。对 $s,t\geq0$，双重级数满足
\[
 \sum_{m,n\geq0}e^{-\lambda(s+t)}
 \frac{(\lambda s)^m(\lambda t)^n}{m!n!}\norm{K^{m+n}f}_\infty
 \leq\norm f_\infty,
\]
故可在 $C_0$ 中按绝对收敛重排。令 $k=m+n$ 并用二项式公式，得到
\begin{align*}
 P_sP_tf
 &=e^{-\lambda(s+t)}
   \sum_{k=0}^\infty
   \left(\sum_{m=0}^k\frac{(\lambda s)^m(\lambda t)^{k-m}}{m!(k-m)!}\right)K^kf\\
 &=e^{-\lambda(s+t)}\sum_{k=0}^\infty
   \frac{(\lambda(s+t))^k}{k!}K^kf
 =P_{s+t}f.
\end{align*}
各项系数非负且和为一，所以 $P_t$ 是正收缩；并且
\[
 \norm{P_tf-f}_\infty
 \leq2(1-e^{-\lambda t})\norm f_\infty\longrightarrow0,
\]
故它在零时刻强连续。现在分离 $n=0,1$ 项得
\begin{align*}
 \frac{P_tf-f}{t}
 &=\frac{e^{-\lambda t}-1}{t}f
 +e^{-\lambda t}\lambda Kf
 +\frac{e^{-\lambda t}}t\sum_{n=2}^\infty\frac{(\lambda t)^n}{n!}K^nf.
\end{align*}
余项满足
\[
 \left\|\frac{e^{-\lambda t}}t
 \sum_{n=2}^\infty\frac{(\lambda t)^n}{n!}K^nf\right\|_\infty
 \le\frac{1-e^{-\lambda t}(1+\lambda t)}t\norm f_\infty\longrightarrow0.
\]
前两项分别趋于 $-\lambda f$ 与 $\lambda Kf$，故在整个 $C_0(X)$ 上
\[
 Af=\lambda(K-I)f.
\]
若只有 $Q$ 的可测性而没有 $K(C_0)\subset C_0$，式
\eqref{eq:7-4-3-poisson-semigroup-series} 仍可在有界可测函数上逐点有意义，却不能据此宣称得到
$C_0$ 上的半群；表~\ref{tab:7-4-3-assumptions-counterexamples} 已给出相应反例。

\textbf{第四步：Laplace 预解算子及预解式。}
设现在 $(P_t)$ 为第一步的一般正收缩强连续半群。对 $\alpha>0$，映射
$t\mapsto e^{-\alpha t}P_tf$ 在 $C_0(X)$ 中连续，且
\[
 \norm{e^{-\alpha t}P_tf}_\infty\le e^{-\alpha t}\norm f_\infty,
 \qquad \int_0^\infty e^{-\alpha t}\norm f_\infty\dd t
 =\frac{\norm f_\infty}{\alpha}.
\]
所以 Bochner 积分 $R_\alpha f$ 存在，并有
$\norm{R_\alpha}\le1/\alpha$。

对 $\alpha,\beta>0$，收缩估计给
\[
 \int_0^\infty\int_0^\infty
 e^{-\alpha s-\beta t}\norm{P_{s+t}f}_\infty\dd t\dd s
 \le\frac{\norm f_\infty}{\alpha\beta}<\infty.
\]
因此 Bochner--Fubini 定理允许交换积分；半群律给
\begin{align*}
 R_\alpha R_\beta f
 &=\int_0^\infty\int_0^\infty
 e^{-\alpha s-\beta t}P_{s+t}f\dd t\dd s\\
 &=\int_0^\infty P_uf
 \left(\int_0^ue^{-\alpha s-\beta(u-s)}\dd s\right)\dd u.
\end{align*}
若 $\alpha\ne\beta$，括号内等于
$(e^{-\alpha u}-e^{-\beta u})/(\beta-\alpha)$，故
\[
 R_\alpha R_\beta f=\frac{R_\alpha f-R_\beta f}{\beta-\alpha},
\]
即所需预解式；$\alpha=\beta$ 时原式两端都为零。

还可直接核对它确实解生成元方程。由有界算子与 Bochner 积分可交换以及半群律，
\begin{align*}
 P_hR_\alpha f
 &=\int_0^\infty e^{-\alpha t}P_{t+h}f\dd t\\
 &=e^{\alpha h}\left(R_\alpha f-
   \int_0^he^{-\alpha u}P_uf\dd u\right).
\end{align*}
于是
\[
 \frac{P_hR_\alpha f-R_\alpha f}{h}
 =\frac{e^{\alpha h}-1}{h}R_\alpha f
 -e^{\alpha h}\frac1h\int_0^he^{-\alpha u}P_uf\dd u.
\]
第一项在范数中趋于 $\alpha R_\alpha f$。第二项的平均在范数中趋于 $f$，因为
\[
 \left\|\frac1h\int_0^he^{-\alpha u}P_uf\dd u-f\right\|_\infty
 \le\frac1h\int_0^h\norm{e^{-\alpha u}P_uf-f}_\infty\dd u\longrightarrow0;
\]
被积函数由 $2\norm f_\infty$ 支配并由强连续性趋于零。因此
\[
 R_\alpha f\in D(A),\qquad
 AR_\alpha f=\alpha R_\alpha f-f,
\]
即 $(\alpha I-A)R_\alpha=I$。这一恒等式是 Hille--Yosida 理论中生成元与预解算子关系的初等形式。
\end{proof}

\begin{exercise}
在定理~\ref{theorem:7-4-1-2} 的证明中，逐项标出局部 Urysohn 引理、有限单位分解、
\Caratheodory{} 定理和层蛋糕公式分别负责哪一个箭头，并说明删去任一项后证明停在哪里。
\end{exercise}

\begin{exercise}
设 $X$ 为紧度量空间，$\mu$ 为概率测度，$f_1,\ldots,f_m$ 都是 $L$-Lipschitz。构造直径至多 $h$
的有限 Borel 分割，并从头证明 \eqref{eq:7-4-2-lipschitz-atomic-error}；比较所得原子数与
Tchakaloff 的 $m+1$ 上界为何控制的是不同量。
\end{exercise}

\begin{exercise}
对 Poisson 跳半群证明 $P_sP_t=P_{s+t}$：在
\eqref{eq:7-4-3-poisson-semigroup-series} 中逐项相乘，写出绝对收敛支配，再用二项式恒等式合并
$K^{m+n}$ 的系数。
\end{exercise}

这些结果通向两个自然方向。Fourier 的 $L^2$ 完备化通向积分方程与 Hilbert 空间；
Riesz--Markov、矩匹配与半群预解式则通向对偶空间、矩问题和算子理论。Banach--Alaoglu
定理、弱紧性、紧算子谱论与谱测度将在后续泛函分析中进一步统一这些方向。

